% !TEX root = ../main.tex
%
% Appendix F — Eight build algorithms in pseudocode for the Dirac--Igusa
% pro-object \mathfrak D_X^{\mathrm{DI}}.  Authored from main.tex line
% ranges 5887--7460 (compact realization opening; finite algebraic
% Dirac--Igusa target system; chiral Koszul source datum) and
% 12147--14346 (Dirac--Igusa realization datum; primitive recognition
% theorem; cofinal finite-window primitive-recognition datum; finite
% retained K3 \times E Hall kernel; geometric hybrid local/wrapped/mixed
% atlas; geometric source chiral Koszul coalgebra).  Cross-repo anchors:
% \texttt{chiral-bar-cobar-vol2} Construction Problem 1 (the source of
% \operatorname{Pf}_{\mathrm{prot}}(\mathfrak D_X) = \Delta_5); \texttt{calabi-yau-quantum-groups}
% Hall--Drinfeld correspondence-programme remark (the K3-side BKM
% Hall--Drinfeld double).
%
% Loaded from main.tex via % !TEX root = ../main.tex
%
% Appendix F — Eight build algorithms in pseudocode for the Dirac--Igusa
% pro-object \mathfrak D_X^{\mathrm{DI}}.  Authored from main.tex line
% ranges 5887--7460 (compact realization opening; finite algebraic
% Dirac--Igusa target system; chiral Koszul source datum) and
% 12147--14346 (Dirac--Igusa realization datum; primitive recognition
% theorem; cofinal finite-window primitive-recognition datum; finite
% retained K3 \times E Hall kernel; geometric hybrid local/wrapped/mixed
% atlas; geometric source chiral Koszul coalgebra).  Cross-repo anchors:
% \texttt{chiral-bar-cobar-vol2} Construction Problem 1 (the source of
% \operatorname{Pf}_{\mathrm{prot}}(\mathfrak D_X) = \Delta_5); \texttt{calabi-yau-quantum-groups}
% Hall--Drinfeld correspondence-programme remark (the K3-side BKM
% Hall--Drinfeld double).
%
% Loaded from main.tex via % !TEX root = ../main.tex
%
% Appendix F — Eight build algorithms in pseudocode for the Dirac--Igusa
% pro-object \mathfrak D_X^{\mathrm{DI}}.  Authored from main.tex line
% ranges 5887--7460 (compact realization opening; finite algebraic
% Dirac--Igusa target system; chiral Koszul source datum) and
% 12147--14346 (Dirac--Igusa realization datum; primitive recognition
% theorem; cofinal finite-window primitive-recognition datum; finite
% retained K3 \times E Hall kernel; geometric hybrid local/wrapped/mixed
% atlas; geometric source chiral Koszul coalgebra).  Cross-repo anchors:
% \texttt{chiral-bar-cobar-vol2} Construction Problem 1 (the source of
% \operatorname{Pf}_{\mathrm{prot}}(\mathfrak D_X) = \Delta_5); \texttt{calabi-yau-quantum-groups}
% Hall--Drinfeld correspondence-programme remark (the K3-side BKM
% Hall--Drinfeld double).
%
% Loaded from main.tex via % !TEX root = ../main.tex
%
% Appendix F — Eight build algorithms in pseudocode for the Dirac--Igusa
% pro-object \mathfrak D_X^{\mathrm{DI}}.  Authored from main.tex line
% ranges 5887--7460 (compact realization opening; finite algebraic
% Dirac--Igusa target system; chiral Koszul source datum) and
% 12147--14346 (Dirac--Igusa realization datum; primitive recognition
% theorem; cofinal finite-window primitive-recognition datum; finite
% retained K3 \times E Hall kernel; geometric hybrid local/wrapped/mixed
% atlas; geometric source chiral Koszul coalgebra).  Cross-repo anchors:
% \texttt{chiral-bar-cobar-vol2} Construction Problem 1 (the source of
% \operatorname{Pf}_{\mathrm{prot}}(\mathfrak D_X) = \Delta_5); \texttt{calabi-yau-quantum-groups}
% Hall--Drinfeld correspondence-programme remark (the K3-side BKM
% Hall--Drinfeld double).
%
% Loaded from main.tex via \input{appendices/F_build_algorithms} after
% \appendix has been opened at main.tex line 17081, so that
% \thesection = \Alph{section}.  No \chapter, no second \appendix.

\section{Eight build algorithms for the Dirac--Igusa pro-object}
\label{app:eight-build-algorithms}

The Dirac--Igusa pro-object
\[
\mathfrak D_X^{\mathrm{DI}}
=\varprojlim_R
\bigl(
\mathcal A^{E_3}_{X,R},
\mathcal F_{X,R}^{\mathrm{hyb}},
\Gamma_{X,R},
\Pi_X,
\Phi_R,
o_R,
\mathcal H_R,
P_R^\Pi,
C_{X,R},
\Theta_{\mathrm{Kos},R},
\mathscr L_{\mathrm{Pf},R},
\operatorname{pf}_{X,R},
\epsilon_{o,R},
\mathrm{Rec}_R
\bigr)
\]
of Definition~\ref{thm:finite-stage-dirac-igusa-pro-object} and
Definition~\ref{def:O1-strong-reduced-orientation} is the conditional source
operator-level package whose protected Pfaffian, after the four
obstructions \((D_0)\), \((O1)\), \((O1)^+\), \((O2)\) are discharged,
realizes the Igusa cusp form,
\[
\operatorname{Pf}_{\mathrm{prot}}(\mathfrak D_X^{\mathrm{DI}})=\Delta_5,\qquad
\epsilon_o=\nu_{\Delta_5}\big|_{W^{(2)}(\La^{2,1}_{II})},\qquad
P_X^{\Pi,\mathrm{red}}\cong\fg_{\Delta_5}.
\]
Construction at finite Harder--Narasimhan height \(R\) proceeds through
eight construction steps.  Each step is a procedure on input data
already supplied by earlier steps and feeds named output data into the
next.  The eight steps stratify the realization datum
\(\mathfrak K_X=(L_X,H_X,O_X,D_X,R_X)\) of
Definition~\ref{thm:finite-stage-dirac-igusa-pro-object}: steps 1--3
realize \((L_X)\) and \((H_X)\); step 4 realizes the local content of
\((O_X)\); step 5 produces the Hall correspondence object; step 6 fixes
the Pfaffian--Dirac line and feeds the orientation character; steps 7--8
realize \((D_X)\) and \((R_X)\).

Throughout, \(X=S\times E\) with \(S\) a smooth complex K3 surface and
\(E\) a smooth elliptic curve.  \(\sigma=\sigma_{s,t}\) is Liu's product
stability condition on the Abramovich--Polishchuk heart in
\(\mathsf{D}^b(\mathrm{Coh}(X))\)
\cite{AbramovichPolishchukTStructures,LiuProductStability}.  The HN
sector \((\sigma,S)\) is fixed and the cofinal subsystem
\(\mathcal R_{\mathrm{HN}}=\{R_0<R_1<\cdots\}\) is the one of
Definition~\ref{def:O1-strong-reduced-orientation}, governed by
\textup{[K3$\times$E-fin]}, \textup{[K3$\times$E-atlas]}, and
\textup{[ML]}.

Each construction step carries a licensing tag pulled from the
five-type universal stage chain
\(\mathsf P\to\mathsf C\to\mathsf S\to\mathsf Z\to\mathsf A\) of
\texttt{chiral-bar-cobar-vol2:CLAUDE.md~§9}: \(\alpha\) primitive,
\(\beta\) functorial passage, \(\gamma\) shadow specialization,
\(\delta\) endpoint convergence, \(\epsilon\) terminal scalar.  The
ordering charge\(\to\)moduli\(\to\)hybrid\(\to\)orientation\(\to\)Hall
\(\to\)Dirac+Pfaffian\(\to\)Koszul\(\to\)recognition is the unique
ordering compatible with input/output type discipline; reordering any
two adjacent steps fails an input check.

The four obstructions are attributed to specific steps:
\((D_0)\) is the cofinal-trivialization datum of the entire eight-step
construction at every finite \(R\) (steps 1--6 supply its components);
\((O1)\) lives at step 4; \((O1)^+\) lives at step 4 with a Coxeter
coherence check feeding step 6; \((O2)\) lives at step 6.  Step 8
inherits the obstruction tower discharged by steps 1--7 and adds the
exact-completion residual \textup{(R8)} of
Definition~\ref{def:cofinal-primitive-recognition-datum}.

\subsection{Construction~1 — Charge window}
\label{app:build-charge-window}

\paragraph{Tag.} \(\alpha\)-primitive (lattice datum).

\paragraph{Input.}
HN height \(R\in\mathcal R_{\mathrm{HN}}\); the formal dyonic Mukai
lattice \(\Gamma_X^{\mathrm{form}}\) of
Definition~\ref{def:gram-index-lattice}; the Gram map
\(\Pi_X:\Gamma_X^{\mathrm{form}}\to\Gamma_{\mathrm{gram}}\) and the
cocycle \(B\) of Lemma~\ref{lem:gram-additivity-defect}; the active
support \(\Gamma_{\mathrm{act}}=\{\gamma\in\Gamma_{\mathrm{eff}}\mid
f(nm,l)\ne0\}\) of the Igusa Jacobi form \(\phi_{0,1}\); the cofinal
subsystem \(\mathcal R_{\mathrm{HN}}\).

\paragraph{Output.}
A finite cusp-positive active window
\[
A_R\subset\Gamma_{\mathrm{act}}\subset\Gamma_{\mathrm{gram}},
\qquad
A_R\hbox{ downward saturated for the cofinal target-window order},
\]
together with the normal-ordered extension
\[
\widehat\Gamma_R=
\{(c,T)\mid c\in\Gamma_R^{HN},\ T\in\mathcal T_R(c)\}
\subset\widehat\Gamma_X,
\qquad
\Gamma_R^\Pi=\overline\Pi_X(\widehat\Gamma_R)\subset\Gamma_{\mathrm{gram}},
\]
the lift \(L:\Gamma_{\mathrm{gram}}\to\Gamma_X^{\mathrm{form}}\) of
Lemma~\ref{lem:primitive-lift-every-gram-triple}, and the finite test
window \(\Gamma_R^{\mathrm{test}}\) of
Definition~\ref{def:O1-strong-reduced-orientation}.

\paragraph{Procedure.}
\begin{enumerate}
\item Compute the active support cone
\[
\Gamma_{\mathrm{act}}=\{(n,l,m)\in\Gamma_{\mathrm{eff}}\mid f(nm,l)\ne 0\}
\]
from the Fourier expansion of the index-one weak Jacobi form
\(\phi_{0,1}(\tau,z)=\sum_{n,l}c(n,l)q^nr^l\) of
\cite{EZ:1}, with the index substitution
\(f(nm,l):=c(nm,l)\) tabulated in
Appendix~\ref{sec:n1-boundary-compatibility-conditions} for the
\(N=1\) Igusa boundary.

\item Set
\[
A_R^{\circ}
=\{\gamma=(n,l,m)\in\Gamma_{\mathrm{act}}\mid n+|l|+m\le R\}.
\]
Form its downward saturation
\[
A_R
=\{\gamma'\in\Gamma_{\mathrm{act}}\mid
\exists\gamma\in A_R^{\circ},\ 0<\gamma'\le\gamma,\
\gamma-\gamma'\in Q_+\},
\]
in the sense of Definition~\ref{def:cofinal-primitive-recognition-datum}.

\item Normal-ordered lift.  For each \(c\in\Gamma_R^{HN}\) compute the
\(B\)-translate orbit
\[
\mathcal T_R(c)
=\Bigl\{\sum_{i<j}B(c_i,c_j)\Bigm|c=\sum_i c_i,\
c_i\in F_{\sigma,S}(R),\ \sum_i h_S(c_i)\le R\Bigr\}.
\]
Set \(\widehat\Gamma_R=\{(c,T):c\in\Gamma_R^{HN},\ T\in\mathcal T_R(c)\}\).

\item Finite target test window:
\(\Gamma_R^{\mathrm{test}}=\overline\Pi_X(\widehat\Gamma_R)\), compatibly
with \(R\to R'\) by inclusion.

\item Output \((A_R,\widehat\Gamma_R,\Gamma_R^\Pi,\Gamma_R^{\mathrm{test}},L)\).
\end{enumerate}

\paragraph{Type checks.}
\(A_R\) is finite and downward saturated;
\(\overline\Pi_X(\widehat\Gamma_R)=\Gamma_R^\Pi\) lies in
\(\Gamma_{\mathrm{gram}}\); \(\Gamma_R^{HN}\subset
F_{\sigma,S}(R)\) is finite by \textup{[K3$\times$E-fin]}; the cocycle
identity
\(\Pi_X(c+c')=\Pi_X(c)+\Pi_X(c')+B(c,c')\) is an axiom of the cocycle.

\paragraph{Obstructions discharged or invoked.}
None; this step is purely lattice combinatorics from the imported
Gritsenko--Nikulin Borcherds product
\cite{GN,GNPartI,GNII,BorcherdsGKM88,Borcherds95}.
The downward saturation is the input requirement of \textup{(R0)} in
Definition~\ref{def:cofinal-primitive-recognition-datum}.

\subsection{Construction~2 — Finite moduli}
\label{app:build-finite-moduli}

\paragraph{Tag.} \(\alpha\)-primitive (geometric source).

\paragraph{Input.}
The active window \(A_R\) and lift \(L\) of
Construction~\ref{app:build-charge-window}; Liu's product stability
condition \(\sigma_{s,t}\); the standing hypotheses
\textup{[K3$\times$E-fin]} and \textup{[K3$\times$E-atlas]}.

\paragraph{Output.}
A finite class set \(C_R\) and lift
\(\ell_R:\Gamma_R^\Pi\to C_R\) with
\(\Pi_X(\ell_R(\gamma))=\gamma\); the family of finite-type quasi-smooth
derived Artin substacks
\[
\mathfrak M^{ss}_{R,c}
=\mathfrak M^{ss}_{\sigma_{s,t}}(X,c)_{\le R}
\subset\operatorname{Perf}(X),
\qquad c\in C_R,
\]
of \(\sigma_{s,t}\)-semistable objects of full Liu numerical class
\(c\) and HN height \(\le R\); the scalar-rigidified truncation
\(\mathfrak M^{\mathrm{sc}}_{R,c}\); the universal perfect complex
\(\mathcal A_{R,c}\) on \(X\times\mathfrak M^{\mathrm{sc}}_{R,c}\); the
finite stratification by retained closed substacks of finite residual
inertia.

\paragraph{Procedure.}
\begin{enumerate}
\item Form \(C_R=\{c=L\gamma:\gamma\in A_R\}\) closed under retained
extension sums (taken in \(\Gamma_X^{\mathrm{form}}\)).

\item For each \(c\in C_R\) cut the substack of
\(\sigma_{s,t}\)-semistable objects of class \(c\) by HN height \(\le
R\):
\[
\mathfrak M^{ss}_{R,c}=
\bigl\{[A]\in\operatorname{Perf}(X)\mid
\sigma_{s,t}\hbox{-semistable},\
\hn(A)\le R,\
\nu(A)=c\bigr\}.
\]

\item Verify finite-typeness: \textup{[K3$\times$E-fin]} (Liu product
stability + product-boundedness theorem at fixed full Liu charge with
finite fibres over the active charge) gives the finite-type assertion.
PTVV \cite{PTVV2013} on the Calabi--Yau threefold \(X\) gives the
\((-1)\)-shifted symplectic structure on
\(\operatorname{Perf}(X)\), restricted to the retained semistable
sector.

\item Quasi-smoothness: the universal perfect complex
\(\mathcal A_{R,c}\) supplies a relative cotangent complex of bounded
Tor-amplitude on the retained charts.  Tangent at \(A\) is
\(T_A\mathfrak M^{ss}_{R,c}\simeq\operatorname{RHom}_X(A,A)[1]\).

\item Scalar rigidification: form \(\mathfrak M^{\mathrm{sc}}_{R,c}\) by
killing the scalar \(\mathbb G_m\)-automorphism of
\(c\)-stable factors.

\item Two-step flag stack
\(\mathfrak F^{(2)}_{R,c,c',c''}\) and extension stack
\(\mathfrak E_{R,c,c'}\) constructed as iterated derived vector stacks
of \(R\pi_*\mathcal RHom_X(\mathcal B,\mathcal A[1])\); both finite type
and quasi-smooth in the retained sector by
Proposition~\ref{prop:finite-dcritical-hall-product}.

\item Output \((C_R,\ell_R,\{\mathfrak M^{ss}_{R,c}\},
\{\mathfrak E_{R,c,c'}\},\{\mathfrak F^{(2)}_{R,c,c',c''}\})\) with
universal perfect complexes and the finite stratification.
\end{enumerate}

\paragraph{Type checks.}
\(\mathfrak M^{ss}_{R,c}\) is a finite-type quasi-smooth derived Artin
substack of \(\operatorname{Perf}(X)\); the scalar rigidification leaves
only finite residual inertia
\(H_S\subset E_{\mathrm{tors}}\) on each connected stabilizer
component, by
Proposition~\ref{prop:finite-dcritical-hall-product};
PTVV \((-1)\)-shifted symplectic structure restricts.

\paragraph{Obstructions discharged or invoked.}
\textup{[K3$\times$E-fin]} is invoked here in its product-boundedness
form and supplies the finite-type assertion not in the cited literature
for the K3\(\times\)E threefold; it is part of the standing datum.
\textup{[K3$\times$E-atlas]} supplies the quasi-smooth d-critical
charts.  These enter \((D_0)\) at the level of the finite Hall atlas
of Theorem~\ref{thm:finite-stage-dirac-igusa-pro-object}.

\subsection{Construction~3 — Hybrid wrapped factorization carrier}
\label{app:build-hybrid-wrapped}

\paragraph{Tag.} \(\beta\)-functorial passage (Stage-1 carrier).

\paragraph{Input.}
The finite moduli stacks
\(\{\mathfrak M^{\mathrm{sc}}_{R,c}\}\) of
Construction~\ref{app:build-finite-moduli}; the
projection-finite/wrapped split by elliptic degree
\(b_R^{\mathrm{geom}}:\Gamma_{\sigma,S}^{HN}(R)\to\Z_{\ge0}\) of
Definition~\ref{def:hybrid-wrapped-factorization-datum}; a fixed
origin \(0_E\in E\); the Abramovich--Polishchuk/Liu wrapped anchor
candidate.

\paragraph{Output.}
The finite hybrid local/wrapped factorization datum
\[
\mathcal F_{X,\sigma,S,\le R}^{\mathrm{hyb}}
=
(\mathcal C_{\sigma,S,\le R},
b_R^{\mathrm{geom}},
\mathcal F_R^{\mathrm{loc}},
\mathcal G_R^{\mathrm{wr}},
\mathfrak E_R^{\mathrm{loc}},
\mathfrak E_R^{\mathrm{hyb}},
\mathfrak E_R^{\mathrm{wr}},
\mathfrak{Flag}_R^{\mathrm{hyb}},
\mathsf{Corr}_R^{E,\mathrm{hyb}},
\lambda_R,
Q_{E,R},
\theta^Q_R,
o_R,
I_R^{\mathrm{prot}})
\]
of Definition~\ref{def:hybrid-wrapped-factorization-datum},
together with the hybrid Ran base
\(\operatorname{Ran}^{\mathrm{hyb}}(E)\), the eight-word two-step
binary flag atlas \(\{w\in\{\mathrm L,\mathrm W\}^3\}\), and the
quotient-after-correspondence functor \(Q_{E,R}\).

\paragraph{Procedure.}
\begin{enumerate}
\item \emph{Elliptic-degree split.}  Decompose
\(\Gamma_{\sigma,S}^{HN}(R)
=\Gamma_R^{\mathrm{loc}}\sqcup\Gamma_R^{\mathrm{wr}}\) where
\(\Gamma_R^{\mathrm{loc}}=\{b_R^{\mathrm{geom}}=0\}\) and
\(\Gamma_R^{\mathrm{wr}}=\{b_R^{\mathrm{geom}}>0\}\).  Verify additivity
\(b_R^{\mathrm{geom}}(\gamma+\gamma')
=b_R^{\mathrm{geom}}(\gamma)+b_R^{\mathrm{geom}}(\gamma')\) on retained
extensions remaining in the quotient.

\item \emph{Projection-finite local sector.}  For
\(\alpha\in\Gamma_R^{\mathrm{loc}}\) and finite index set \(I\), build
the closed-support incidence stack over \(E^I\):
\[
\mathfrak M^{\mathrm{loc,cl}}_{\alpha,R,I}
=
\Bigl\{(A,\mathbf e)\in\mathfrak M^{\mathrm{sc}}_{R,L\alpha}\times E^I
\Bigm|p_E(\mathrm{Supp}\,A)\subset\bigcup_{i\in I}S\times\{e_i\}\Bigr\}.
\]
Form
\(\mathscr H^{\mathrm{loc}}_{\alpha,R,I}=
(\pi_{\alpha,R,I})_!(\Phi_{\alpha,R,I}^{\mathrm{loc}}\otimes
o_{\alpha,R,I}^{\mathrm{loc}})\) and the open-support sheaf
\(\mathcal F_{\alpha,R}^{\mathrm{loc}}(U)\) over \(U\subset E\) by
restricted Borel--Moore configuration sums.

\item \emph{Wrapped prequotient sector.}  For
\(\eta\in\Gamma_R^{\mathrm{wr}}\), form the rigidified prequotient
\(\mathcal M_{\eta,R}^{\mathrm{wr,rig}}\to\mathcal M_{\eta,R}^{\mathrm{wr}}\)
by the Abramovich--Polishchuk / determinant anchor
\[
\lambda_\eta:\mathcal M_{\eta,R}^{\mathrm{wr,rig}}\longrightarrow E,
\qquad
\lambda_\eta(A)=\det Rp_{E,*}A
\otimes\mathcal O_E(-\chi(A)\cdot 0_E),
\]
divided by a finite cover when \(\chi(A)\ne 1\) so that
\(\lambda_\eta(t\cdot A)=\lambda_\eta(A)+t\) holds.  Verify the
losslessness residual \(\mathfrak o^{\lambda,\mathrm{loss}}_R=0\) by
the Ext-distinguishing test of
Definition~\ref{def:hybrid-wrapped-factorization-datum}, separating
\(W_0=i_{s,*}(\mathcal O_E\oplus\mathcal O_E)\) from
\(W_L=i_{s,*}(L\oplus L^{-1})\).

\item \emph{Mixed local/wrapped correspondences.}  For an ordered
hybrid input
\(\alpha_-\in\Gamma_R^{\mathrm{loc}}\),
\(\eta\in\Gamma_R^{\mathrm{wr}}\),
\(\beta_+\in\Gamma_R^{\mathrm{loc}}\), form
\(\mathfrak E^{\mathrm{mix}}_{\alpha_-;\eta;\beta_+,R}\) before the
reduced \(E\)-quotient.

\item \emph{Eight-word two-step flag atlas.}  Build flag stacks
\(\mathfrak F^{(2),w}_R\) for each
\(w\in\{\mathrm L,\mathrm W\}^3=
\{\mathrm{LLL},\mathrm{LLW},\mathrm{LWL},\mathrm{WLL},\mathrm{LWW},
\mathrm{WLW},\mathrm{WWL},\mathrm{WWW}\}\), classifying filtrations
with three associated graded factors of those types.  Verify the
four-input pentagon coherence and the quotient-after-correspondence
descent.

\item \emph{Hybrid correspondence-module functor.}  Assemble the
correspondence category
\(\mathsf{Corr}^{E,\mathrm{hyb}}_R\) and the functor
\(\mathcal B_R^{E,\mathrm{hyb}}:\mathsf{Corr}^{E,\mathrm{hyb}}_R
\to\mathsf{Ch}_{\C}\) by
\(\mathcal B_R^{E,\mathrm{hyb}}(e)=q_!\,\mathrm{TS}^{\mathrm{red}}_e\,p^*\)
for every generating arrow
\(e=(Y\xleftarrow{p}\mathfrak E_e\xrightarrow{q}Z)\).

\item \emph{Quotient functor \(Q_{E,R}\).}  Apply
\(E\)-equivariant compactly supported vanishing-cycle chains and the
orientation descent of step~4 to produce the reduced object
\[
\mathcal F^{\mathrm{geom}}_{X,\sigma,S,\le R}
=Q_{E,R}\circ\mathcal B_R^{E,\mathrm{hyb}}.
\]

\item Output the hybrid datum with all coherences and the protected
integration map \(I_R^{\mathrm{prot}}\), with homogeneous
\(s\)-degree the Borcherds trace exponent
\(m_R=\operatorname{pr}_3\overline\Pi_X\).
\end{enumerate}

\paragraph{Type checks.}
The carrier
\(\mathcal F_{X,\sigma,S,\le R}^{\mathrm{hyb}}\) is a hybrid
correspondence-module functor on
\(\mathsf{Corr}^{E,\mathrm{hyb}}_R\) with values in
\(\mathsf{Ch}_{\C}\); the four-input pentagon is supplied by
the eight-word two-step flag atlas; the anchor residual tuple
\(\mathfrak o^\lambda_R=
(\mathfrak o^{\lambda,\mathrm{ex}}_R,
\mathfrak o^{\lambda,\mathrm{unit}}_R,
\mathfrak o^{\lambda,\mathrm{loss}}_R,
\mathfrak o^{\lambda,\mathrm{multi}}_R,
\mathfrak o^{\lambda,\mathrm{tr}}_{R'R})\) vanishes; the hybrid
residual \(\mathfrak O_{\mathrm{hyb},R}=0\) at every height \(R\),
compatibly with transition maps.

\paragraph{Obstructions discharged or invoked.}
The middle type-II wall lies in
\(\delta_2\leftrightarrow(0,1,1)\), elliptic degree
\(d=m+1=2\) by Lemma~\ref{app:pfaffian-wall-three-walls}; it
cannot be represented by a projection-finite local stack and demands
the wrapped prequotient.  This step's correctness — the existence of
the wrapped representative
\(x_{\delta_2,R}\in\mathcal M_{\eta,R}^{\mathrm{wr,rig}}\) with
\(\overline\Pi_X(x_{\delta_2,R})=(0,1,1)\) and its compatible anchor
\(\lambda_{\eta,R}\) — is precisely the hybrid input to \((O2)\) at
the middle wall.

\subsection{Construction~4 — Reduced d-critical orientation
and Joyce--Upmeier transport}
\label{app:build-orientation}

\paragraph{Tag.} \(\beta\)-functorial passage (oriented carrier).

\paragraph{Input.}
The hybrid carrier
\(\mathcal F_{X,\sigma,S,\le R}^{\mathrm{hyb}}\) of
Construction~\ref{app:build-hybrid-wrapped}; the K3 holomorphic
\(2\)-form \(\sigma_S\) on \(S\); BBDJS d-critical chart machinery
\cite{BBDJS2015}; the Joyce--Upmeier multiplicative orientation
formalism \cite{JoyceUpmeier,JoyceTanakaUpmeier}; the cosection
localization machinery of Kiem--Li and the cosection-reduced
obstruction theory; the type-II Weyl group
\(W^{(2)}(\La^{2,1}_{II})\) generated by
\(s_{\delta_1},s_{\delta_2},s_{\delta_3}\).

\paragraph{Output.}
The cosection
\(\sigma=p_S^*\sigma_S\); the reduced tangent complex
\[
\operatorname{RHom}_{\mathrm{red}}(A,A)
=
\operatorname{Cone}\!\Bigl(
\operatorname{RHom}_X(A,A)
\xrightarrow{(\mathrm{tr},\mathrm{cs})}
R\Gamma(X,\mathcal O_X)\oplus H^2(S,\mathcal O_S)[-1]
\Bigr)[-1];
\]
the BBDJS vanishing-cycle complex \(\Phi_{R,c}\); the reduced
orientation line
\[
o_{R,c}
=\det\operatorname{Ext}^1_{\mathrm{red}}(A,A)^{-1}
\quad\text{on Darboux charts},\qquad
o_{R,c}^{\otimes 2}\simeq\det\operatorname{RHom}_{\mathrm{red}}(A,A);
\]
the Joyce--Upmeier multiplicative transport
\(o_{c\oplus c'}\simeq o_c\otimes o_{c'}\) compatible with the reduced
extension correspondences; the quotient orientation
\(\mathfrak Q^{or}_R\) of
Theorem~\ref{prop:generic-local-sign-orientation-character}; the
Weyl-equivariant lifts
\(\tau_{i,R,S}:o_{R,S}\to s_{\delta_i}^*o_{R,S}\) for \(i=1,2,3\) of
\((O1)^+\); the local generic wall-sign data \(\epsilon_{o,R}^{\mathrm{loc}}\).

\paragraph{Procedure.}
\begin{enumerate}
\item \emph{Cosection.}  Pull back the K3 semiregularity cosection
\(\mathrm{cs}\) along \(p_S:X\to S\) to obtain the 2-step cosection
\(\sigma=p_S^*\sigma_S\) on the K3 factor; verify additivity in exact
triangles by \textup{(RedOr)} of
Theorem~\ref{thm:finite-stage-dirac-igusa-pro-object}.

\item \emph{Reduced obstruction.}  Form the reduced tangent complex by
the displayed cone above (Kiem--Li cosection localization).  PTVV
\cite{PTVV2013} \((-1)\)-shifted symplectic on
\(\operatorname{RHom}_X(A,A)\) descends, after the cosection
contraction, to the reduced complex.

\item \emph{BBDJS d-critical chart.}  On each Darboux chart of
\(\mathfrak M^{\mathrm{sc}}_{R,c}\), build the BBDJS vanishing-cycle
complex \(\Phi_{R,c}\) \cite{BBDJS2015}; verify reduced
Thom--Sebastiani transport coherent on two-step flag stacks of
Construction~\ref{app:build-hybrid-wrapped}.

\item \emph{Reduced orientation line.}  Define
\(o_{R,c}=\det\operatorname{Ext}^1_{\mathrm{red}}(A,A)^{-1}\); verify
the squared isomorphism
\(o_{R,c}^{\otimes 2}\simeq\det\operatorname{RHom}_{\mathrm{red}}(A,A)\)
by \(3\)-Calabi--Yau Serre duality
\(\operatorname{Ext}^2_{\mathrm{red}}(A,A)
\simeq\operatorname{Ext}^1_{\mathrm{red}}(A,A)^\vee\).

\item \emph{Quotient-orientation descent on the K3\(\times\)E-quotient
self-Ext frame bundle.}  Compute the four Cech--Borel obstruction
classes for every retained stratum \(S\):
\[
\alpha^{\mathrm{red}}_{R,S}=w_2(\mathscr D_{R,S}),
\quad
\alpha^{E,\mathrm{free}}_{R,S}=
\mathrm{edge}_{2,0}(w_2^E(\mathscr D_{R,S})),
\quad
\widetilde\beta^H_{R,S}=w_2^G(\mathscr D_{R,S}),
\quad
\lambda^H_{R,S}\in H^1_G(S;\mathbb F_2),
\]
and verify simultaneous vanishing with compatible
null-trivializations on overlaps, on Hall correspondence pullbacks,
on two-step flag stacks, and under successor maps.  The full
calculation proceeds primary stabilizer by stabilizer; for cyclic
\(\Z/2^a\) the absent mixed coefficient requires the rank-two
two-primary check.  Output the descended quotient orientation system
\(\mathfrak Q^{or}_R\).

\item \emph{Multiplicative Joyce--Upmeier transport.}  Transport
\(o_{c\oplus c'}\simeq o_c\otimes o_{c'}\) coherently across two-step
flag stacks
\cite{JoyceUpmeier,JoyceTanakaUpmeier}; verify the reduced
Thom--Sebastiani identification on three-input flags.

\item \emph{Weyl-equivariant lifts \((O1)^+\).}  For each
\(i\in\{1,2,3\}\) construct the lift
\(\tau_{i,R,S}:o_{R,S}\to s_{\delta_i}^*o_{R,S}\); verify the
projective Weyl cocycle
\(c_{o,R}\in Z^2(\mathcal G_R;\underline{\mathbb F}_2)\) on the
retained action groupoid; null-trivialize \([c_{o,R}]\) and choose a
\(1\)-cochain killing it; verify \(\tau_{i,R,S}^2=1\); verify the
Coxeter relations
\((s_{\delta_i}s_{\delta_j})^3=1\) for \(i\ne j\) (the
\((3,3,3)\)-hyperbolic triangle relations of
Appendix~\ref{app:pfaffian-wall-three-walls}); verify that the torsor
defects \(\omega_{i,\mathcal C}\) of
Definition~\ref{thm:finite-stage-dirac-igusa-pro-object}~$(O_X)$ vanish.

\item \emph{Local wall-sign datum.}  At each generic type-II wall point
\(\Hpl_{\delta_i}\), with the rank-one local model
\(K_{\delta_i}^{\mathrm{cross}}=[\R\xrightarrow{u_i}\R]\), record the
local generic sign \(\epsilon_{o,R}^{\mathrm{loc}}(s_{\delta_i})=-1\)
of Proposition~\ref{prop:rank-one-crossing-o2-sign}.

\item Output
\((o_R,\mathfrak Q^{or}_R,\{\tau_{i,R,S}\},\epsilon_{o,R}^{\mathrm{loc}})\).
\end{enumerate}

\paragraph{Type checks.}
\(o_R\) is a Picard-line section squaring to the reduced determinant;
the four Cech--Borel classes vanish with compatible
null-trivializations; the lifts \(\tau_{i,R}\) satisfy
\(\tau_{i,R}^2=1\) and the \((3,3,3)\)-Coxeter relations; multiplicative
transport agrees on flag stacks.

\paragraph{Obstructions discharged or invoked.}
This step is the entire site of \((O1)\) (strong reduced orientation
on the K3\(\times\)E-quotient self-Ext frame bundle) and of
\((O1)^+\) (Weyl-equivariant transport along
\(W^{(2)}(\La^{2,1}_{II})\) reflections).  Discharging \((O1)\)
requires the simultaneous vanishing and compatible
null-trivialization of the four Cech--Borel classes on every retained
stratum; discharging \((O1)^+\) requires \([c_{o,R}]=0\), the
Coxeter coherence, and the vanishing of all torsor defects
\(\omega_{i,\mathcal C}\).  The local wall-sign datum is the input to
\((O2)\) at step~6.

\subsection{Construction~5 — Hall correspondences and primitive part}
\label{app:build-hall}

\paragraph{Tag.} \(\gamma\)-shadow specialization
(Stage-2 chiral Hall shadow).

\paragraph{Input.}
The oriented hybrid carrier from
Construction~\ref{app:build-orientation}; the extension stack
\(\mathfrak E^{\mathrm{geom}}_{R,\widehat c,\widehat c'}\); the
two-step flag stack
\(\mathfrak F^{(2),\mathrm{geom}}_{R,\widehat c,\widehat c',\widehat c''}\);
the BBDJS reduced Thom--Sebastiani transport; six-functor admissibility
of \(p^*,q_!,p_!\) (Construction~\ref{app:build-finite-moduli}~step
\textup{(Six)}).

\paragraph{Output.}
The finite Hall object
\[
\mathcal H_R^{\mathrm{geom}}
=\bigoplus_{\widehat c\in\widehat\Gamma_R}
H_c^\bullet\!\bigl(
\mathfrak M^{\mathrm{geom}}_{R,\widehat c},
\Phi_{R,c}\otimes o_{R,c}\bigr);
\]
the Hall product \(m_R=q_!\,\mathrm{TS}^{\mathrm{red}}\,p^*\) and Hall
coproduct \(\Delta_R=p_!\,(\mathrm{TS}^{\mathrm{red}})^{-1}\,q^*\); the
finite Hall bialgebra
\((\mathcal H_R^{\mathrm{geom}},m_R,\Delta_R,\eta_R,\epsilon_R)\); the
primitive part
\[
P_R^{\mathrm{geom}}
=\ker(\Delta_R-1\otimes\mathrm{id}-\mathrm{id}\otimes 1)\cap\ker\epsilon_R;
\]
the normal-ordered Gram descent
\(P_R^{\Pi,\mathrm{geom}}
=\overline\Pi_{X,*}^{\Theta_R}P_R^{\mathrm{geom}}\); the
\(\Gamma_R^\Pi\)-graded super-Lie bracket
\([\,\bar x,\bar y\,]_{\mathrm{red}}\) of
Proposition~\ref{thm:hall-primitives-formal};
the Hopf pairing matrix \(G_{\gamma,ab}\) and radical \(K_\gamma\).

\paragraph{Procedure.}
\begin{enumerate}
\item \emph{Pullback.}  For
\(\widehat c,\widehat c'\in\widehat\Gamma_R\), pull back along
\(p:\mathfrak E^{\mathrm{geom}}_{R,\widehat c,\widehat c'}
\to\mathfrak M^{\mathrm{geom}}_{R,\widehat c}
\times\mathfrak M^{\mathrm{geom}}_{R,\widehat c'}\); apply reduced
Thom--Sebastiani \(\mathrm{TS}^{\mathrm{red}}\) to combine vanishing
cycle data \cite{BBDJS2015}; pushforward exceptionally along
\(q:\mathfrak E^{\mathrm{geom}}_{R,\widehat c,\widehat c'}
\to\mathfrak M^{\mathrm{geom}}_{R,\widehat c\star\widehat c'}\) to
obtain
\(m_R(x\otimes y)=q_!\,\mathrm{TS}^{\mathrm{red}}\,p^*(x\boxtimes y)\).

\item \emph{Coproduct.}  Dually pull back along \(q\), apply
inverse reduced Thom--Sebastiani, pushforward along \(p\) to obtain
\(\Delta_R(z)=p_!\,(\mathrm{TS}^{\mathrm{red}})^{-1}\,q^*(z)\).
The compact geometric coproduct correspondence
\(\mathfrak C^{\mathrm{geom}}_{\rho;\mu,\nu}\) of
Proposition~\ref{thm:hall-primitives-formal} is the
support data for splitting.

\item \emph{Bialgebra check.}  Verify associativity of \(m_R\) and
coassociativity of \(\Delta_R\) on the two-step flag stack
\(\mathfrak F^{(2),\mathrm{geom}}_{R,\widehat c,\widehat c',\widehat c''}\)
by base change and Thom--Sebastiani coherence.  Verify the bialgebra
identity (Hopf compatibility) via the Massey--Saito Joyce--Upmeier
multiplicative orientation transport.

\item \emph{Primitive extraction.}  Compute
\(P_R^{\mathrm{geom}}=\ker(\Delta_R-1\otimes\mathrm{id}
-\mathrm{id}\otimes 1)\cap\ker\epsilon_R\); since
\(\Delta_R-1\otimes\mathrm{id}-\mathrm{id}\otimes 1\) is degree-graded
on \(\widehat\Gamma_R\), this reduces to a finite linear computation
in each charge.

\item \emph{Normal-ordered descent.}  Apply
\(\overline\Pi_{X,*}^{\Theta_R}\) of
Theorem~\ref{thm:hall-product} to obtain the
\(\Gamma_R^\Pi\)-graded
\(P_R^{\Pi,\mathrm{geom}}\); verify additivity
\(\overline\Pi_X(\widehat c\star\widehat c')
=\overline\Pi_X(\widehat c)+\overline\Pi_X(\widehat c')\).

\item \emph{Bracket and pairing matrices.}  Compute the matrices
\[
M^{\alpha+\beta,c}_{\alpha,a;\beta,b}
=\int_{\mathfrak E^{\mathrm{geom}}_{\alpha,\beta;\alpha+\beta}}
q^!(e_{\alpha+\beta,c}^\vee)\cap
\mathrm{TS}^{\mathrm{red}}(p^*(e_{\alpha,a}\boxtimes e_{\beta,b})),
\quad
B_{\alpha,\beta}(x\otimes y)
=M_{\alpha,\beta}(x\otimes y)
-(-1)^{|x||y|}M_{\beta,\alpha}\mathsf{sw}(x\otimes y),
\]
\[
G_{\gamma,ab}=
\mathrm{tr}_0(q_!\,\mathrm{TS}^{\mathrm{red}}\,p^*
(e_{\gamma,a}\boxtimes e_{-\gamma,b}))
\quad\hbox{on}\quad
\mathfrak P^{\mathrm{geom}}_{\gamma,-\gamma;0}.
\]
Form the radical
\(K_{\gamma,\bar p}
=\ker G_{\gamma,\bar p}\cap\ker G_{-\gamma,\bar p}^t\) and the
quotient \(L_{\gamma,\bar p}\); fix splittings.

\item Output
\((\mathcal H_R^{\mathrm{geom}},m_R,\Delta_R,P_R^{\mathrm{geom}},
P_R^{\Pi,\mathrm{geom}},
\{B_{\alpha,\beta}\},\{G_{\gamma,ab}\},\{K_\gamma\},\{L_\gamma\})\).
\end{enumerate}

\paragraph{Type checks.}
\(\mathcal H_R^{\mathrm{geom}}\) is a finite \(\widehat\Gamma_R\)-graded
Hall bialgebra; \(m_R\) is associative on the two-step flag stack;
\(\Delta_R\) is coassociative; the supercommutator of primitives is
primitive (the bialgebra identity); the descended bracket is
\(\Gamma_R^\Pi\)-graded after \(\overline\Pi_{X,*}^{\Theta_R}\);
finite-dimensional in each homogeneous parity piece by
finite-typeness.

\paragraph{Obstructions discharged or invoked.}
\((D_0)\) is partially discharged: the finite Hall package
\((\mathcal H_R^{\mathrm{geom}},m_R,\Delta_R,P_R)\) is constructed.
The seven obstruction classes
\(\mathfrak o_\Theta^{\mathrm{top}},
\mathfrak o_\Theta^{\mathrm{Hoch}},
\mathfrak o_\Theta^{\mathrm{tri}},
\mathfrak o_\Theta^{\mathrm{Jac}},
\mathfrak o_F,
\mathfrak o_\Theta^{\mathrm{cyc}},
\mathfrak o_\Theta^{\mathrm{rad}}\) of
Theorem~\ref{thm:hall-product} are required to vanish on
this finite Hall stage and compatibly in the inverse limit.

\paragraph{Cross-volume anchor.}  The output
\((\mathcal H_R^{\mathrm{geom}},m_R)\) is the finite-stage realization
of the K3-side BKM Hall--Drinfeld double
\(\mathcal D_\hbar(\mathrm{CoHA}_{K3\times E})\) of
\texttt{calabi-yau-quantum-groups:CLAUDE.md~§6.7}.  In the universal
stage chain \(\mathsf P\to\mathsf C\to\mathsf S\) the Hall positive half
\(Y^+(X)\) is precisely \(P_R^{\mathrm{geom}}\) at finite \(R\); the
Drinfeld double \(G(X)=D(Y^+(X))\) requires Hopf pairing, completion,
integral form, and stable-envelope transport which are supplied by
\(G_{\gamma,ab}\) and \((\mathrm{ML})\).  The κ-tuple coordinate
\(\kappa_{\mathrm{BKM}}=5\) is the weight of \(\Delta_5\) and the
output of the level-4 scalar trace; it is not a level-2 invariant of
the carrier itself.

\subsection{Construction~6 — Dirac operator and Pfaffian line}
\label{app:build-dirac-pfaffian}

\paragraph{Tag.} \(\gamma\)-shadow specialization
(orientation gerbe + Pfaffian).

\paragraph{Input.}
The oriented Hall package
\((\mathcal H_R^{\mathrm{geom}},o_R,\{\tau_{i,R,S}\})\) of
Constructions~\ref{app:build-orientation}--\ref{app:build-hall}; the
local wall-sign datum
\(\epsilon_{o,R}^{\mathrm{loc}}\) of
Construction~\ref{app:build-orientation}; the rank-one type-II
Pfaffian crossing datum of
Definition~\ref{def:rank-one-type-ii-crossing} and
Appendix~\ref{app:pfaffian-wall-rank-one}; the half-Hilbert orbit
under \(W^{(2)}(\La^{2,1}_{II})\) of size \(2^{12}=4096\).

\paragraph{Output.}
The Dirac operator on the orientation gerbe
\(D_R\); the Pfaffian line
\(\mathscr L_{\mathrm{Pf},R}=
\det\operatorname{Ext}^1_{\mathrm{red}}(A,A)^{-1}\) (squaring to the
reduced determinant); the protected Pfaffian section
\(\operatorname{pf}_{X,R}\); the global wall-sign character
\(\epsilon_{o,R}:W^{(2)}(\La^{2,1}_{II})\to\{\pm1\}\); the Pfaffian
section asymptotic
\[
\operatorname{pf}_{X,R}
=
64q^{1/2}r^{1/2}s^{1/2}
\prod_{\gamma=(n,l,m)\in A_R}
(1-q^nr^ls^m)^{f(nm,l)}+\hbox{higher-window correction},
\]
which exhausts to \(\Delta_5\) as \(R\to\infty\).

\paragraph{Procedure.}
\begin{enumerate}
\item \emph{Dirac operator.}  On each Darboux chart of
\(\mathfrak M^{\mathrm{sc}}_{R,c}\), with reduced normal self-Ext
splitting \(K^{\mathrm{nor}}_{R,c}\), build the odd Dirac operator
\(D_{R,c}\) acting on \(\Omega^{\mathrm{sp}}\otimes o_{R,c}\)
(spinors on the orientation gerbe); verify that, at a generic
type-II wall point, the local rank-one normal block is
\(D_{\delta_i}=\begin{psmallmatrix}0&u\\-u&0\end{psmallmatrix}\) with
\(\operatorname{Pf}(D_{\delta_i})=u\)
(Definition~\ref{def:rank-one-type-ii-crossing}).

\item \emph{Pfaffian line.}  Set
\(\mathscr L_{\mathrm{Pf},R}=
\det\operatorname{Ext}^1_{\mathrm{red}}(A,A)^{-1}\); verify
\(\mathscr L_{\mathrm{Pf},R}^{\otimes 2}\simeq
\det\operatorname{RHom}_{\mathrm{red}}(A,A)\) by
\(3\)-Calabi--Yau Serre duality.

\item \emph{Local wall sign.}  At each generic point of
\(\Hpl_{\delta_i}\) with normal model
\(K^{\mathrm{nor}}_{\delta_i,R}=[\R\xrightarrow{u_i}\R]\), compute
the monodromy of the Pfaffian section
\(\operatorname{pf}_{\delta_i,R}=\upsilon_iu_i\) along the
wall-reflection loop \(\ell_{\delta_i,R,S}\) closed by the
\((O1)^+\) lift \(\tau_{i,R,S}\):
\[
M_{\ell_{\delta_i,R,S}}(\upsilon_iu_i)
=\upsilon_i(-u_i)=-\upsilon_iu_i,
\qquad
\epsilon_{o,R}^{\mathrm{loc}}(s_{\delta_i})
=(-1)^{N^{\mathrm{Pf}}_{\delta_i,R}}=-1.
\]

\item \emph{Half-Hilbert orbit transport.}  The reflection
\(s_{\delta_i}\) acts on the half-Hilbert orbit (compact source labels
of size \(2^{12}=4096\) parametrized by 12 binary parities tied to the
\(2^{12}\) odd-spin structures on the lattice
\(\La^{2,1}_{II}/2\La^{2,1}_{II}\)).  Sum the local rank-one
contributions over the orbit:
\[
\epsilon_{o,R}(s_{\delta_i})
=\prod_{x\in\hbox{orbit}/2^{12}}
\epsilon_{o,R}^{\mathrm{loc}}(s_{\delta_i};x).
\]
The output is the global character
\(\epsilon_{o,R}:W^{(2)}(\La^{2,1}_{II})\to\{\pm 1\}\).

\item \emph{Compatibility package.}  Verify the component, boundary,
overlap, quotient, and protected integration compatibility package of
Proposition~\ref{prop:appE-local-pfaffian-sign} on every
retained type-II wall stratum.

\item \emph{Pfaffian section.}  Define
\(\operatorname{pf}_{X,R}\) as the protected integration of the
oriented Hall product over
\(\mathcal F_{X,\sigma,S,\le R}^{\mathrm{hyb}}\):
\[
\operatorname{pf}_{X,R}
=
64q^{1/2}r^{1/2}s^{1/2}
I_R^{\mathrm{prot}}\!\left(
\bigotimes_{\gamma\in A_R}\operatorname{pf}_{\gamma,R}\right),
\qquad
\operatorname{pf}_{\gamma,R}=(1-q^nr^ls^m)^{f(nm,l)}.
\]

\item Output
\((D_R,\mathscr L_{\mathrm{Pf},R},\operatorname{pf}_{X,R},
\epsilon_{o,R})\), with the protected integration compatibility data.
\end{enumerate}

\paragraph{Type checks.}
\(\mathscr L_{\mathrm{Pf},R}\) is a Picard-line whose square
isomorphism is supplied by reduced Serre duality;
\(\operatorname{pf}_{X,R}\) is a section of
\(\mathscr L_{\mathrm{Pf},R}\) with the displayed cusp expansion; the
character \(\epsilon_{o,R}\) takes values in \(\{\pm 1\}\) and
satisfies the Coxeter relations of
\(W^{(2)}(\La^{2,1}_{II})\); on each simple type-II reflection
\(\epsilon_{o,R}(s_{\delta_i})=-1\) by the rank-one local model.

\paragraph{Obstructions discharged or invoked.}
This step is the entire site of \((O2)\) (local Pfaffian wall normal
form on type-II walls; the \(2^{12}=4096\) sign sum on the
half-Hilbert orbit transport).  The local rank-one Pfaffian crossing
of Appendix~\ref{app:pfaffian-wall-rank-one} discharges \((O2)\)
generically; the global half-Hilbert orbit transport remains the open
\((O2)\) component of
Open~Problem~\ref{def:O1-plus-weyl-transport}.  The character
identity
\(\epsilon_o=\nu_{\Delta_5}\big|_{W^{(2)}(\La^{2,1}_{II})}\) becomes a
theorem only after \((D_0)\), \((O1)\), \((O1)^+\), and \((O2)\) are
discharged compatibly in \(R\).

\subsection{Construction~7 — Chiral Koszul source coalgebra}
\label{app:build-koszul}

\paragraph{Tag.} \(\gamma\)-shadow specialization
(\emph{Bar = twisting; not bulk}; the bulk is
\(Z^{\mathrm{der}}_{\mathrm{ch}}(\mathcal F_X^{\mathrm{hyb}})\),
constructed elsewhere).

\paragraph{Input.}
The oriented Hall primitives
\(P_R^{\Pi,\mathrm{geom}}\) of
Construction~\ref{app:build-hall}; the hybrid carrier
\(\mathcal F_{X,\sigma,S,\le R}^{\mathrm{hyb}}\) of
Construction~\ref{app:build-hybrid-wrapped}; the augmentation
\(\varepsilon_{X,R}^{\mathrm{hyb}}:
\mathcal F_{X,\sigma,S,\le R}^{\mathrm{hyb}}\to k\mathbf 1\); the
hybrid residual \(\mathfrak O_{\mathrm{hyb},R}=0\); the bounded
length \(N_R\) (positivity of HN height on every non-vacuum colour);
the Beilinson--Drinfeld/Francis--Gaitsgory bar-cobar formalism
\cite{BD,FrancisGaitsgoryCKD}; the conilpotent completed domain on
the elliptic curve \(E\).

\paragraph{Output.}
The conilpotent chiral coalgebra on \(E\):
\[
C_{X,R}^{\mathrm{geom}}
=\bar B_E^{\mathrm{ch}}
(\mathcal F_{X,\sigma,S,\le R}^{\mathrm{hyb}},
\varepsilon_{X,R}^{\mathrm{hyb}})
=k\mathbf 1\oplus\bigoplus_{1\le p\le N_R}
\bar B_E^{\mathrm{ch},p}(\overline{\mathcal F}_{X,R}^{\mathrm{geom}});
\]
the bar-length filtration
\(F^{\mathrm{co}}_{R,p}\); the chiral collision coproduct
\(\Delta_R^{\mathrm{ch}}\); the bar comparison
\(b_{X,R}=\mathrm{id}_{\bar B_E^{\mathrm{ch}}(\cdot)}\); the source
cobar counit
\(\varepsilon^{\mathrm{src}}_{\mathrm{bar},R}:
\Omega_E^{\mathrm{ch}}C_{X,R}^{\mathrm{geom}}
\xrightarrow{\simeq}
\mathcal F_{X,\sigma,S,\le R}^{\mathrm{hyb}}\); after a separate
source-to-target quasi-isomorphism
\(\vartheta_R:\mathcal F_{X,\sigma,S,\le R}^{\mathrm{hyb}}
\xrightarrow{\simeq}\mathsf{Den}_{\Delta_5,E,\le R}\), the cobar
comparison
\(\Theta_{\mathrm{Kos},R}=\vartheta_R\circ\varepsilon^{\mathrm{src}}_{\mathrm{bar},R}\).

\paragraph{Procedure.}
\begin{enumerate}
\item \emph{Augmentation kernel.}  Set
\(\overline{\mathcal F}_{X,R}^{\mathrm{geom}}
=\ker(\varepsilon_{X,R}^{\mathrm{hyb}})\); verify the bounded-length
hypothesis (positivity of HN height on every non-vacuum surviving
colour, hence at most \(N_R\) factors in any non-zero word).

\item \emph{Bar coalgebra.}  Form the reduced chiral bar
\[
C_{X,R}^{\mathrm{geom}}
=k\mathbf 1\oplus\bigoplus_{1\le p\le N_R}
\bar B_E^{\mathrm{ch},p}
(\overline{\mathcal F}_{X,R}^{\mathrm{geom}}),
\]
length \(p\) summand generated by ordered \(p\)-fold inputs from the
augmentation ideal of total HN height \(\le R\) (height \(>R\) outputs
zero); coaugmentation by the inclusion of the vacuum summand; counit by
projection.

\item \emph{Bar-length filtration.}
\[
F^{\mathrm{co}}_{R,p}C_{X,R}^{\mathrm{geom}}
=k\mathbf 1\oplus\bigoplus_{1\le j\le p}
\bar B_E^{\mathrm{ch},j}
(\overline{\mathcal F}_{X,R}^{\mathrm{geom}}),\qquad
0\le p\le N_R.
\]

\item \emph{Collision coproduct.}  For
\(x=[x_1|\cdots|x_p]\), define the cut-map coproduct
\[
\Delta_R^{\mathrm{ch}}(x)
=\mathbf1\otimes x+x\otimes\mathbf1
+\sum_{i=1}^{p-1}[x_1|\cdots|x_i]\otimes[x_{i+1}|\cdots|x_p],
\]
with each cut realized by the corresponding finite collision/flag
pull-push map in
\(\mathsf{Corr}^{E,\mathrm{hyb}}_{R,\mathrm{geom}}\) of
Construction~\ref{app:build-hybrid-wrapped}.  Verify
coassociativity by the four-input pentagon and the eight-word flag
coherences; verify counit identities via vacuum collision flags.

\item \emph{Conilpotence.}  Verify
\((\overline\Delta_R^{\mathrm{bar}})^{N_R+1}=0\): each reduced cut
strictly lowers bar length on each branch.

\item \emph{Bar comparison.}  Take \(b_{X,R}\) to be the identity of
\(\bar B_E^{\mathrm{ch}}
(\mathcal F_{X,\sigma,S,\le R}^{\mathrm{hyb}},
\varepsilon_{X,R}^{\mathrm{hyb}})\), eliminating the coalgebra-side
finite Koszul residuals
\(\mathfrak o^C_R, \mathfrak o^{\mathrm{fil}}_R,
\mathfrak o^{\mathrm{conil}}_R, \mathfrak o^{\Delta,\mathrm{ch}}_R\)
of Definition~\ref{def:chiral-koszul-source-datum}.

\item \emph{Source cobar counit.}  Apply Francis--Gaitsgory chiral
Koszul duality
\cite[Theorem~6.2.2]{FrancisGaitsgoryCKD} in the conilpotent completed
domain to obtain
\(\varepsilon^{\mathrm{src}}_{\mathrm{bar},R}:
\Omega_E^{\mathrm{ch}}C_{X,R}^{\mathrm{geom}}
\xrightarrow{\simeq}
\mathcal F_{X,\sigma,S,\le R}^{\mathrm{hyb}}\); verify it is a
source-internal quasi-isomorphism of chiral algebras on \(E\), not a
homotopy inverse of the target counit.

\item \emph{Source-to-target identification.}  Construct
\(\vartheta_R\) from the geometric primitive representatives, Hall
products, coproducts, pairings, and HN transition compatibility;
verify the residual
\(\mathfrak o^\vartheta_R\) of
Corollary~\ref{cor:source-bar-remaining-cobar-obstruction} vanishes;
set \(\Theta_{\mathrm{Kos},R}=\vartheta_R\circ
\varepsilon^{\mathrm{src}}_{\mathrm{bar},R}\).

\item \emph{Inverse limit.}  Verify the Koszul Mittag--Leffler
condition \(R^1\varprojlim=0\) for the inverse systems of source
coalgebras, completed chiral bar/cobar constructions, target
truncations, and the cones of \(b_{X,R}\) and
\(\Theta_{\mathrm{Kos},R}\) \cite[\S3.5]{WeibelHA}.  Form
\[
C_X^{\mathrm{geom}}
=\varprojlim_R C_{X,R}^{\mathrm{geom}},\qquad
\Theta_{\mathrm{Kos}}:
\Omega_E^{\mathrm{ch}}C_X^{\mathrm{geom}}
\xrightarrow{\simeq}\mathsf{Den}_{\Delta_5,E}.
\]

\item Output
\((C_{X,R}^{\mathrm{geom}},F^{\mathrm{co}}_{R,\bullet},
\Delta_R^{\mathrm{ch}},b_{X,R},\Theta_{\mathrm{Kos},R})\) at every
\(R\), with successor compatibility.
\end{enumerate}

\paragraph{Type checks.}
\(C_{X,R}^{\mathrm{geom}}\) is a coaugmented conilpotent chiral
coalgebra on \(E\) with bounded bar-length filtration of length at
most \(N_R\); \(\Delta_R^{\mathrm{ch}}\) is coassociative and counital
on the chosen flag stacks; the bar-cobar inversion is supplied by
Francis--Gaitsgory chiral Koszul; the cobar comparison
\(\Theta_{\mathrm{Kos},R}\) is source-internal, not the target
bar-cobar counit (the source/target separation of
Lemma~\ref{lem:source-target-koszul-separation}).

\paragraph{Obstructions discharged or invoked.}
The Koszul Mittag--Leffler exactness hypothesis is the
\textup{(D0)}-component for the chiral Koszul lane.  The total finite
Koszul residual
\(\mathfrak O_{\mathrm{Kos},R}=
(\mathfrak o^C_R,\mathfrak o^{\mathrm{fil}}_R,
\mathfrak o^{\mathrm{conil}}_R,\mathfrak o^{\Delta,\mathrm{ch}}_R,
\{\mathfrak o^{C,\mathrm{tr}}_{R'R}\},\mathfrak o^b_R,
\mathfrak o^\Omega_R,\mathfrak o^{\mathrm{hyb}}_R,
\mathfrak o^\Pi_R,\mathfrak o^{\mathrm{sep}}_R,
\mathfrak o^{\Prim}_R)\)
vanishes after this step.

\paragraph{Cross-volume firewall.}  The bar coalgebra
\(\bar B_E^{\mathrm{ch}}
(\mathcal F_{X,R}^{\mathrm{hyb}})\) is a \emph{twisting / coupling}
coalgebra (Maurer--Cartan classification), \emph{not} the bulk
\(Z^{\mathrm{der}}_{\mathrm{ch}}(\mathcal F_X^{\mathrm{hyb}})
=\mathrm{ChirHoch}^\bullet(\mathcal F_X^{\mathrm{hyb}},
\mathcal F_X^{\mathrm{hyb}})\); cf.\
\texttt{chiral-bar-cobar-vol2:CLAUDE.md~§6} (universal stage chain
\(\mathsf P\to\mathsf C\to\mathsf S\to\mathsf Z\)).
The chiral Koszul comparison
\(\Theta_{\mathrm{Kos}}:\Omega_E^{\mathrm{ch}}C_X^{\mathrm{geom}}
\xrightarrow{\simeq}\mathsf{Den}_{\Delta_5,E}\) is a comparison
between two algebras of the same level; the homotopy inverse to the
target-internal counit
\(\Omega_E^{\mathrm{ch}}\bar B_E^{\mathrm{ch}}
(\mathsf{Den}_{\Delta_5,E})\to\mathsf{Den}_{\Delta_5,E}\) is not used
in either direction.

\subsection{Construction~8 — Primitive recognition}
\label{app:build-recognition}

\paragraph{Tag.} \(\delta\)-endpoint convergence
(source-to-target identification).

\paragraph{Input.}
The descended source primitive Lie superalgebra
\(P_X^{\Pi,\mathrm{red}}\) of
Construction~\ref{app:build-hall} (after the radical quotient
\(\overline{\mathrm{Rad}\langle\,\cdot,\cdot\rangle_X}\)); the imported
target denominator algebra
\(\fg_{\Delta_5}=G(\mathscr B_{\Delta_5})\) on
\(\La^{2,1}_{II}\) of Definition~\ref{def:bkm-imaginary-simple-roots} and
Definition~\ref{def:bkm-relations}; the increasing cofinal
sequence of finite downward-saturated windows
\(W_1\subset W_2\subset\cdots\) with
\(\bigcup_\nu W_\nu=\mathcal R_+\); the Borcherds--Kac--Moody data
\((H,I,p,\alpha_i,(\,,\,))\) of
Theorem~\ref{thm:primitive-recognition}.

\paragraph{Output.}
The cofinal finite-window primitive-recognition datum
\(\{\textup{(R0),(R1),(R2),(R3),(R4),(R5),(R6),(R7),(R8)}\}\) of
Definition~\ref{def:cofinal-primitive-recognition-datum} for every
\(W_\nu\), compatible under \(W_\nu\subset W_{\nu+1}\); the source
representatives
\(e_i^X,f_i^X,h_i^X\); the source presentation map
\(\pi_W:\widehat{\mathfrak F}(\mathscr B_{\Delta_5})_W
\to P_X^{\Pi,\mathrm{red}}\big|_W\); the kernel equality
\(\ker\pi_W=
(\overline{J_{\mathrm{BK}}+\mathrm{Rad}_{\mathrm{GN}}})_W\); the
isomorphism
\[
\boxed{\quad
P_X^{\Pi,\mathrm{red}}\cong\fg_{\Delta_5},\quad
\Omega_E^{\mathrm{ch}}C_X\simeq\mathsf{Den}_{\Delta_5,E}.
\quad}
\]

\paragraph{Procedure.}
For every cofinal \(W_\nu\):
\begin{enumerate}
\item[\textup{(R0)}] \emph{Compact source labels.}  Verify every
finite source basis vector in \(W_\nu\cup(-W_\nu)\cup\{0\}\) satisfies
the compact source-admissibility condition of
Theorem~\ref{thm:finite-stage-dirac-igusa-pro-object}.

\item[\textup{(R1)}] \emph{Representatives and parities.}  Supply
parity-homogeneous primitive representatives
\(e_i^X\in P^{\Pi,\mathrm{red}}_{X,\alpha_i}\),
\(f_i^X\in P^{\Pi,\mathrm{red}}_{X,-\alpha_i}\),
\(h_i^X=\iota_H^{-1}(\alpha_i)\) for every simple-root index \(i\) with
\(\alpha_i\in W_\nu\), with the Gritsenko--Nikulin parity fibres
\(\mu_{\overline 0}(a),\mu_{\overline 1}(a)\) for imaginary simple
roots and even real representatives for \(\delta_1,\delta_2,\delta_3\).

\item[\textup{(R2)}] \emph{Relations.}  Verify in the protected Hall
super-bracket the Cartan, Chevalley, real Serre, Borcherds
orthogonality, and super sign relations for every relation whose
displayed source and target degrees lie in
\(W_\nu\cup(-W_\nu)\cup\{0\}\).  In the real rank-three subsystem the
Cartan matrix is
\(((\delta_i,\delta_j))=\mathrm{diag}(2)-2(J-\mathrm{diag}(1))\) with
real Serre exponent~3.

\item[\textup{(R3)}] \emph{Full finite parity dimensions.}  For every
\(\beta\in W_\nu\) verify
\(\dim P^{\Pi,\mathrm{red}}_{X,\beta,\overline p}
=\rootmult_{\overline p}^{\mathrm{GN}}(\beta)\) for both parities
\(p=0,1\).  In the first timelike channel
\(\beta=\delta_{123}\),
Proposition~\ref{prop:bkm-delta123-presentation-count} requires
\((d^{\mathrm{GN}}_{(1,1,1),\overline 0},
d^{\mathrm{GN}}_{(1,1,1),\overline 1})=(29,93)\).  Determinant
information \(29-93=-64=f(1,1)\) is consistent but not the
recognition.

\item[\textup{(R4)}] \emph{Hopf pairing and radical.}  Compute the
positive-negative pairing matrices in parity blocks; verify the
kernel agrees with the
Gritsenko--Nikulin/Kac invariant-pairing kernel; verify Lie-ideal and
coideal stability under all finite brackets; verify carrying under
HN transitions to the next window.

\item[\textup{(R5)}] \emph{No extra relations.}  Verify the kernel
equality
\(\ker\pi_\nu=(\overline{J_{\mathrm{BK}}+\mathrm{Rad}_{\mathrm{GN}}})_{\le W_\nu}\) on the finite triangular span generated by
\(W_\nu\cup(-W_\nu)\cup\{0\}\).

\item[\textup{(R6)}] \emph{Generation.}  Verify that every source
primitive root space in \(W_\nu\cup(-W_\nu)\), after the finite
radical quotient, is generated by the Cartan and the supplied
simple-root representatives through brackets whose intermediate
positive degrees stay in \(W_\nu\).

\item[\textup{(R7)}] \emph{Completed PBW compatibility.}  Verify the
PBW-graded isomorphism
\(\mathrm{gr}_{\mathrm{PBW}}U(P_X^{\Pi,\mathrm{red}})_{\le W_\nu}
\xrightarrow{\sim}
\mathrm{gr}_{\mathrm{PBW}}U(\fg_{\Delta_5})_{\le W_\nu}\); verify
strict preservation under \(W_{\nu+1}\to W_\nu\) transition maps.

\item[\textup{(R8)}] \emph{Exact triangular completion.}  Verify
strictness of the transition maps and closedness of images on every
finite window, so that passage to the inverse limit commutes with the
finite kernels, cokernels, radicals, PBW associated gradeds, and
parity pieces.  Equivalently the relevant \(\lim^1\) terms vanish for
these triangular finite-window systems.

\item Output the global isomorphism
\(P_X^{\Pi,\mathrm{red}}\cong\fg_{\Delta_5}\) by passage to the inverse
limit (compatible cofinal datum); together with the chiral Koszul
comparison of Construction~\ref{app:build-koszul} this gives
\(\Omega_E^{\mathrm{ch}}C_X\simeq\mathsf{Den}_{\Delta_5,E}\) and hence
the protected Pfaffian identification
\(\operatorname{Pf}_{\mathrm{prot}}(\mathfrak D_X^{\mathrm{DI}})=\Delta_5\).
\end{enumerate}

\paragraph{Type checks.}
The source presentation map \(\pi:\widehat{\mathfrak F}(\mathscr B_{\Delta_5})\to P_X^{\Pi,\mathrm{red}}\) is
continuous and degree-preserving; \textup{(R5)} gives injectivity
after the completed radical quotient; \textup{(R6)} gives
surjectivity; \textup{(R7)} gives PBW compatibility; \textup{(R8)}
gives exactness in the inverse-limit step.  No claim on signed
superdimensions \(\sdim=f(nm,l)\) substitutes for the two-integer
parity dimensions of \textup{(R3)}.

\paragraph{Obstructions discharged or invoked.}
This step inherits the obstruction tower discharged by
Constructions~1--7 and adds the cofinal exact-completion residual
\textup{(R8)}.  The Beilinson--Drinfeld bar-cobar firewall is
respected throughout: \textup{(R5)} is a source kernel-equality
theorem, not an inversion of the target bar-cobar counit;
\textup{(R7)} is a PBW graded comparison of source and target
enveloping algebras, not a graph-equality argument from the scalar
square.

\paragraph{Cross-volume anchor.}  The output is
\texttt{chiral-bar-cobar-vol2:CLAUDE.md~§9} Construction
Problem~1: the operator-level construction of
\(\mathfrak D_X^{\mathrm{DI}}\) for \(K3\times E\) with
\(\operatorname{Pf}_{\mathrm{prot}}(\mathfrak D_X^{\mathrm{DI}})=\Delta_5\), under the
four obstructions \((D_0),(O1),(O1)^+,(O2)\).  The Drinfeld doubling
\(Y^+(X)\to G(X)=D(Y^+(X))\) of
\texttt{calabi-yau-quantum-groups:CLAUDE.md} requires the Hopf
pairing of \textup{(R4)}, completion by \textup{(R8)}, integral form,
and stable-envelope transport which are supplied by the same
data.  The κ-tuple coordinate \(\kappa_{\mathrm{BKM}}=5\) is the
weight of \(\Delta_5\) and lives at level~4 of the universal stage
chain; it is a consequence of the construction, not an input to it.

\subsection{The eight-step audit}
\label{app:build-audit}

The eight construction steps form a strict ordered sequence of
input/output type discipline.  An audit of any candidate construction of
\(\mathfrak D_X^{\mathrm{DI}}\) reduces to verifying, for every
\(R\in\mathcal R_{\mathrm{HN}}\):
\begin{enumerate}
\item the lattice combinatorics of
Construction~\ref{app:build-charge-window};
\item the finite-type quasi-smooth derived stacks of
Construction~\ref{app:build-finite-moduli};
\item the hybrid local/wrapped factorization datum and the eight-word
two-step flag atlas of Construction~\ref{app:build-hybrid-wrapped};
\item the four Cech--Borel orientation classes and the
\((O1)^+\)~Coxeter coherence of
Construction~\ref{app:build-orientation};
\item the finite Hall bialgebra and primitive matrices of
Construction~\ref{app:build-hall};
\item the rank-one Pfaffian crossing and the half-Hilbert orbit
transport of Construction~\ref{app:build-dirac-pfaffian};
\item the bar coalgebra, conilpotent collision coproduct, and the
Francis--Gaitsgory cobar quasi-isomorphism of
Construction~\ref{app:build-koszul};
\item the cofinal datum
\(\textup{(R0)}\)--\(\textup{(R8)}\) of
Construction~\ref{app:build-recognition};
\end{enumerate}
together with the compatible vanishing of the obstruction tuples
\((\mathfrak O^{E_3-\mathrm{src}}_R,
\mathfrak O_{\mathrm{hyb},R},
\mathfrak O_{\mathrm{Kos},R},
\mathfrak O^{\mathrm{O1}}_R,
\mathfrak O^{\mathrm{O1}^+}_R,
\mathfrak O^{\mathrm{O2}}_R,
\mathfrak O_R^{D_0})\) at every height \(R\), and the
Mittag--Leffler exactness condition \textup{[ML]} of
Definition~\ref{def:O1-strong-reduced-orientation} on the cofinal HN inverse
system.  Each obstruction is attributable to the unique step
in which its components are constructed.

\paragraph{Order rigidity.}  Each of the seven adjacent input/output
dependencies fails type-checking under exchange:
\begin{enumerate}
\item charge window $\to$ finite moduli: $\Lambda_{\mathrm{ample}}$
  parametrises the HN strata; without the window, the moduli stack
  has no finite cofinal subsystem.
\item finite moduli $\to$ hybrid carrier: $\mathrm{Ran}^{\mathrm{hyb}}(E)$
  needs the finite Hall stack as factor on which to graft the
  wrapped/local atlas.
\item hybrid carrier $\to$ reduced orientation: the cosection-twisted
  Joyce--Upmeier line bundle is defined on the carrier, not on the
  raw moduli (the cosection lives on $\mathrm{Ran}^{\mathrm{hyb}}$).
\item reduced orientation $\to$ Hall correspondences: the Pfaffian
  line bundle on the Hall stack is the orientation-twisted sign data;
  without orientation, the Hall product is unsigned.
\item Hall correspondences $\to$ Pfaffian line + Dirac operator: the
  primitive Hall bracket $[-,-]_X$ defines the Hall--Drinfeld
  correspondence whose Pfaffian section is the Dirac--Igusa Pfaffian.
\item Pfaffian line $\to$ chiral Koszul source coalgebra: the source
  $C_X$ is the bar of the augmented hybrid factorization algebra
  pulled back along the Pfaffian comparison
  $\iota_{\mathrm{aut}}$.
\item chiral Koszul source $\to$ primitive recognition: the primitive
  part $P_X^{\Pi,\mathrm{red}}$ is the Frobenius--Serre primitive
  filtration on the source, which the recognition theorem matches to
  $\mathfrak g_{\Delta_5}$.
\end{enumerate}
Two diagnostic exchanges: orientation before hybrid leaves the middle
type-II wall \(\delta_2\leftrightarrow(0,1,1)\) without a
representative, by
Lemma~\ref{app:pfaffian-wall-three-walls}; recognition before the
Koszul cobar leaves the identification
\(\Omega_E^{\mathrm{ch}}C_X\simeq\mathsf{Den}_{\Delta_5,E}\)
unconstructed.  All other adjacent exchanges fail by the same
input/output type criterion.

\paragraph{Conditional theorem.}  Granted the eight constructions and
the four obstructions \((D_0),(O1),(O1)^+,(O2)\),
Theorem~\ref{thm:pfaffian-dirac-finite-stage} reads
\[
\operatorname{Pf}_{\mathrm{prot}}(\mathfrak D_X^{\mathrm{DI}})=\Delta_5,\qquad
\epsilon_o=\nu_{\Delta_5}\big|_{W^{(2)}(\La^{2,1}_{II})},\qquad
P_X^{\Pi,\mathrm{red}}\cong\fg_{\Delta_5}.
\]
The scalar square
\(Z_{\mathrm{BPS}}^{K3\times E}=
(\Phi_{10}^{\mathrm{un}})^{-1}=\Delta_5^{-2}\) is the level-4 trace
of the constructed protected algebra; promotion to the 3d
gravitational path integral is not claimed and requires the
Hall--Borcherds residual of
\texttt{chiral-bar-cobar-vol2}.

%%% End of Appendix F.
 after
% \appendix has been opened at main.tex line 17081, so that
% \thesection = \Alph{section}.  No \chapter, no second \appendix.

\section{Eight build algorithms for the Dirac--Igusa pro-object}
\label{app:eight-build-algorithms}

The Dirac--Igusa pro-object
\[
\mathfrak D_X^{\mathrm{DI}}
=\varprojlim_R
\bigl(
\mathcal A^{E_3}_{X,R},
\mathcal F_{X,R}^{\mathrm{hyb}},
\Gamma_{X,R},
\Pi_X,
\Phi_R,
o_R,
\mathcal H_R,
P_R^\Pi,
C_{X,R},
\Theta_{\mathrm{Kos},R},
\mathscr L_{\mathrm{Pf},R},
\operatorname{pf}_{X,R},
\epsilon_{o,R},
\mathrm{Rec}_R
\bigr)
\]
of Definition~\ref{thm:finite-stage-dirac-igusa-pro-object} and
Definition~\ref{def:O1-strong-reduced-orientation} is the conditional source
operator-level package whose protected Pfaffian, after the four
obstructions \((D_0)\), \((O1)\), \((O1)^+\), \((O2)\) are discharged,
realizes the Igusa cusp form,
\[
\operatorname{Pf}_{\mathrm{prot}}(\mathfrak D_X^{\mathrm{DI}})=\Delta_5,\qquad
\epsilon_o=\nu_{\Delta_5}\big|_{W^{(2)}(\La^{2,1}_{II})},\qquad
P_X^{\Pi,\mathrm{red}}\cong\fg_{\Delta_5}.
\]
Construction at finite Harder--Narasimhan height \(R\) proceeds through
eight construction steps.  Each step is a procedure on input data
already supplied by earlier steps and feeds named output data into the
next.  The eight steps stratify the realization datum
\(\mathfrak K_X=(L_X,H_X,O_X,D_X,R_X)\) of
Definition~\ref{thm:finite-stage-dirac-igusa-pro-object}: steps 1--3
realize \((L_X)\) and \((H_X)\); step 4 realizes the local content of
\((O_X)\); step 5 produces the Hall correspondence object; step 6 fixes
the Pfaffian--Dirac line and feeds the orientation character; steps 7--8
realize \((D_X)\) and \((R_X)\).

Throughout, \(X=S\times E\) with \(S\) a smooth complex K3 surface and
\(E\) a smooth elliptic curve.  \(\sigma=\sigma_{s,t}\) is Liu's product
stability condition on the Abramovich--Polishchuk heart in
\(\mathsf{D}^b(\mathrm{Coh}(X))\)
\cite{AbramovichPolishchukTStructures,LiuProductStability}.  The HN
sector \((\sigma,S)\) is fixed and the cofinal subsystem
\(\mathcal R_{\mathrm{HN}}=\{R_0<R_1<\cdots\}\) is the one of
Definition~\ref{def:O1-strong-reduced-orientation}, governed by
\textup{[K3$\times$E-fin]}, \textup{[K3$\times$E-atlas]}, and
\textup{[ML]}.

Each construction step carries a licensing tag pulled from the
five-type universal stage chain
\(\mathsf P\to\mathsf C\to\mathsf S\to\mathsf Z\to\mathsf A\) of
\texttt{chiral-bar-cobar-vol2:CLAUDE.md~§9}: \(\alpha\) primitive,
\(\beta\) functorial passage, \(\gamma\) shadow specialization,
\(\delta\) endpoint convergence, \(\epsilon\) terminal scalar.  The
ordering charge\(\to\)moduli\(\to\)hybrid\(\to\)orientation\(\to\)Hall
\(\to\)Dirac+Pfaffian\(\to\)Koszul\(\to\)recognition is the unique
ordering compatible with input/output type discipline; reordering any
two adjacent steps fails an input check.

The four obstructions are attributed to specific steps:
\((D_0)\) is the cofinal-trivialization datum of the entire eight-step
construction at every finite \(R\) (steps 1--6 supply its components);
\((O1)\) lives at step 4; \((O1)^+\) lives at step 4 with a Coxeter
coherence check feeding step 6; \((O2)\) lives at step 6.  Step 8
inherits the obstruction tower discharged by steps 1--7 and adds the
exact-completion residual \textup{(R8)} of
Definition~\ref{def:cofinal-primitive-recognition-datum}.

\subsection{Construction~1 — Charge window}
\label{app:build-charge-window}

\paragraph{Tag.} \(\alpha\)-primitive (lattice datum).

\paragraph{Input.}
HN height \(R\in\mathcal R_{\mathrm{HN}}\); the formal dyonic Mukai
lattice \(\Gamma_X^{\mathrm{form}}\) of
Definition~\ref{def:gram-index-lattice}; the Gram map
\(\Pi_X:\Gamma_X^{\mathrm{form}}\to\Gamma_{\mathrm{gram}}\) and the
cocycle \(B\) of Lemma~\ref{lem:gram-additivity-defect}; the active
support \(\Gamma_{\mathrm{act}}=\{\gamma\in\Gamma_{\mathrm{eff}}\mid
f(nm,l)\ne0\}\) of the Igusa Jacobi form \(\phi_{0,1}\); the cofinal
subsystem \(\mathcal R_{\mathrm{HN}}\).

\paragraph{Output.}
A finite cusp-positive active window
\[
A_R\subset\Gamma_{\mathrm{act}}\subset\Gamma_{\mathrm{gram}},
\qquad
A_R\hbox{ downward saturated for the cofinal target-window order},
\]
together with the normal-ordered extension
\[
\widehat\Gamma_R=
\{(c,T)\mid c\in\Gamma_R^{HN},\ T\in\mathcal T_R(c)\}
\subset\widehat\Gamma_X,
\qquad
\Gamma_R^\Pi=\overline\Pi_X(\widehat\Gamma_R)\subset\Gamma_{\mathrm{gram}},
\]
the lift \(L:\Gamma_{\mathrm{gram}}\to\Gamma_X^{\mathrm{form}}\) of
Lemma~\ref{lem:primitive-lift-every-gram-triple}, and the finite test
window \(\Gamma_R^{\mathrm{test}}\) of
Definition~\ref{def:O1-strong-reduced-orientation}.

\paragraph{Procedure.}
\begin{enumerate}
\item Compute the active support cone
\[
\Gamma_{\mathrm{act}}=\{(n,l,m)\in\Gamma_{\mathrm{eff}}\mid f(nm,l)\ne 0\}
\]
from the Fourier expansion of the index-one weak Jacobi form
\(\phi_{0,1}(\tau,z)=\sum_{n,l}c(n,l)q^nr^l\) of
\cite{EZ:1}, with the index substitution
\(f(nm,l):=c(nm,l)\) tabulated in
Appendix~\ref{sec:n1-boundary-compatibility-conditions} for the
\(N=1\) Igusa boundary.

\item Set
\[
A_R^{\circ}
=\{\gamma=(n,l,m)\in\Gamma_{\mathrm{act}}\mid n+|l|+m\le R\}.
\]
Form its downward saturation
\[
A_R
=\{\gamma'\in\Gamma_{\mathrm{act}}\mid
\exists\gamma\in A_R^{\circ},\ 0<\gamma'\le\gamma,\
\gamma-\gamma'\in Q_+\},
\]
in the sense of Definition~\ref{def:cofinal-primitive-recognition-datum}.

\item Normal-ordered lift.  For each \(c\in\Gamma_R^{HN}\) compute the
\(B\)-translate orbit
\[
\mathcal T_R(c)
=\Bigl\{\sum_{i<j}B(c_i,c_j)\Bigm|c=\sum_i c_i,\
c_i\in F_{\sigma,S}(R),\ \sum_i h_S(c_i)\le R\Bigr\}.
\]
Set \(\widehat\Gamma_R=\{(c,T):c\in\Gamma_R^{HN},\ T\in\mathcal T_R(c)\}\).

\item Finite target test window:
\(\Gamma_R^{\mathrm{test}}=\overline\Pi_X(\widehat\Gamma_R)\), compatibly
with \(R\to R'\) by inclusion.

\item Output \((A_R,\widehat\Gamma_R,\Gamma_R^\Pi,\Gamma_R^{\mathrm{test}},L)\).
\end{enumerate}

\paragraph{Type checks.}
\(A_R\) is finite and downward saturated;
\(\overline\Pi_X(\widehat\Gamma_R)=\Gamma_R^\Pi\) lies in
\(\Gamma_{\mathrm{gram}}\); \(\Gamma_R^{HN}\subset
F_{\sigma,S}(R)\) is finite by \textup{[K3$\times$E-fin]}; the cocycle
identity
\(\Pi_X(c+c')=\Pi_X(c)+\Pi_X(c')+B(c,c')\) is an axiom of the cocycle.

\paragraph{Obstructions discharged or invoked.}
None; this step is purely lattice combinatorics from the imported
Gritsenko--Nikulin Borcherds product
\cite{GN,GNPartI,GNII,BorcherdsGKM88,Borcherds95}.
The downward saturation is the input requirement of \textup{(R0)} in
Definition~\ref{def:cofinal-primitive-recognition-datum}.

\subsection{Construction~2 — Finite moduli}
\label{app:build-finite-moduli}

\paragraph{Tag.} \(\alpha\)-primitive (geometric source).

\paragraph{Input.}
The active window \(A_R\) and lift \(L\) of
Construction~\ref{app:build-charge-window}; Liu's product stability
condition \(\sigma_{s,t}\); the standing hypotheses
\textup{[K3$\times$E-fin]} and \textup{[K3$\times$E-atlas]}.

\paragraph{Output.}
A finite class set \(C_R\) and lift
\(\ell_R:\Gamma_R^\Pi\to C_R\) with
\(\Pi_X(\ell_R(\gamma))=\gamma\); the family of finite-type quasi-smooth
derived Artin substacks
\[
\mathfrak M^{ss}_{R,c}
=\mathfrak M^{ss}_{\sigma_{s,t}}(X,c)_{\le R}
\subset\operatorname{Perf}(X),
\qquad c\in C_R,
\]
of \(\sigma_{s,t}\)-semistable objects of full Liu numerical class
\(c\) and HN height \(\le R\); the scalar-rigidified truncation
\(\mathfrak M^{\mathrm{sc}}_{R,c}\); the universal perfect complex
\(\mathcal A_{R,c}\) on \(X\times\mathfrak M^{\mathrm{sc}}_{R,c}\); the
finite stratification by retained closed substacks of finite residual
inertia.

\paragraph{Procedure.}
\begin{enumerate}
\item Form \(C_R=\{c=L\gamma:\gamma\in A_R\}\) closed under retained
extension sums (taken in \(\Gamma_X^{\mathrm{form}}\)).

\item For each \(c\in C_R\) cut the substack of
\(\sigma_{s,t}\)-semistable objects of class \(c\) by HN height \(\le
R\):
\[
\mathfrak M^{ss}_{R,c}=
\bigl\{[A]\in\operatorname{Perf}(X)\mid
\sigma_{s,t}\hbox{-semistable},\
\hn(A)\le R,\
\nu(A)=c\bigr\}.
\]

\item Verify finite-typeness: \textup{[K3$\times$E-fin]} (Liu product
stability + product-boundedness theorem at fixed full Liu charge with
finite fibres over the active charge) gives the finite-type assertion.
PTVV \cite{PTVV2013} on the Calabi--Yau threefold \(X\) gives the
\((-1)\)-shifted symplectic structure on
\(\operatorname{Perf}(X)\), restricted to the retained semistable
sector.

\item Quasi-smoothness: the universal perfect complex
\(\mathcal A_{R,c}\) supplies a relative cotangent complex of bounded
Tor-amplitude on the retained charts.  Tangent at \(A\) is
\(T_A\mathfrak M^{ss}_{R,c}\simeq\operatorname{RHom}_X(A,A)[1]\).

\item Scalar rigidification: form \(\mathfrak M^{\mathrm{sc}}_{R,c}\) by
killing the scalar \(\mathbb G_m\)-automorphism of
\(c\)-stable factors.

\item Two-step flag stack
\(\mathfrak F^{(2)}_{R,c,c',c''}\) and extension stack
\(\mathfrak E_{R,c,c'}\) constructed as iterated derived vector stacks
of \(R\pi_*\mathcal RHom_X(\mathcal B,\mathcal A[1])\); both finite type
and quasi-smooth in the retained sector by
Proposition~\ref{prop:finite-dcritical-hall-product}.

\item Output \((C_R,\ell_R,\{\mathfrak M^{ss}_{R,c}\},
\{\mathfrak E_{R,c,c'}\},\{\mathfrak F^{(2)}_{R,c,c',c''}\})\) with
universal perfect complexes and the finite stratification.
\end{enumerate}

\paragraph{Type checks.}
\(\mathfrak M^{ss}_{R,c}\) is a finite-type quasi-smooth derived Artin
substack of \(\operatorname{Perf}(X)\); the scalar rigidification leaves
only finite residual inertia
\(H_S\subset E_{\mathrm{tors}}\) on each connected stabilizer
component, by
Proposition~\ref{prop:finite-dcritical-hall-product};
PTVV \((-1)\)-shifted symplectic structure restricts.

\paragraph{Obstructions discharged or invoked.}
\textup{[K3$\times$E-fin]} is invoked here in its product-boundedness
form and supplies the finite-type assertion not in the cited literature
for the K3\(\times\)E threefold; it is part of the standing datum.
\textup{[K3$\times$E-atlas]} supplies the quasi-smooth d-critical
charts.  These enter \((D_0)\) at the level of the finite Hall atlas
of Theorem~\ref{thm:finite-stage-dirac-igusa-pro-object}.

\subsection{Construction~3 — Hybrid wrapped factorization carrier}
\label{app:build-hybrid-wrapped}

\paragraph{Tag.} \(\beta\)-functorial passage (Stage-1 carrier).

\paragraph{Input.}
The finite moduli stacks
\(\{\mathfrak M^{\mathrm{sc}}_{R,c}\}\) of
Construction~\ref{app:build-finite-moduli}; the
projection-finite/wrapped split by elliptic degree
\(b_R^{\mathrm{geom}}:\Gamma_{\sigma,S}^{HN}(R)\to\Z_{\ge0}\) of
Definition~\ref{def:hybrid-wrapped-factorization-datum}; a fixed
origin \(0_E\in E\); the Abramovich--Polishchuk/Liu wrapped anchor
candidate.

\paragraph{Output.}
The finite hybrid local/wrapped factorization datum
\[
\mathcal F_{X,\sigma,S,\le R}^{\mathrm{hyb}}
=
(\mathcal C_{\sigma,S,\le R},
b_R^{\mathrm{geom}},
\mathcal F_R^{\mathrm{loc}},
\mathcal G_R^{\mathrm{wr}},
\mathfrak E_R^{\mathrm{loc}},
\mathfrak E_R^{\mathrm{hyb}},
\mathfrak E_R^{\mathrm{wr}},
\mathfrak{Flag}_R^{\mathrm{hyb}},
\mathsf{Corr}_R^{E,\mathrm{hyb}},
\lambda_R,
Q_{E,R},
\theta^Q_R,
o_R,
I_R^{\mathrm{prot}})
\]
of Definition~\ref{def:hybrid-wrapped-factorization-datum},
together with the hybrid Ran base
\(\operatorname{Ran}^{\mathrm{hyb}}(E)\), the eight-word two-step
binary flag atlas \(\{w\in\{\mathrm L,\mathrm W\}^3\}\), and the
quotient-after-correspondence functor \(Q_{E,R}\).

\paragraph{Procedure.}
\begin{enumerate}
\item \emph{Elliptic-degree split.}  Decompose
\(\Gamma_{\sigma,S}^{HN}(R)
=\Gamma_R^{\mathrm{loc}}\sqcup\Gamma_R^{\mathrm{wr}}\) where
\(\Gamma_R^{\mathrm{loc}}=\{b_R^{\mathrm{geom}}=0\}\) and
\(\Gamma_R^{\mathrm{wr}}=\{b_R^{\mathrm{geom}}>0\}\).  Verify additivity
\(b_R^{\mathrm{geom}}(\gamma+\gamma')
=b_R^{\mathrm{geom}}(\gamma)+b_R^{\mathrm{geom}}(\gamma')\) on retained
extensions remaining in the quotient.

\item \emph{Projection-finite local sector.}  For
\(\alpha\in\Gamma_R^{\mathrm{loc}}\) and finite index set \(I\), build
the closed-support incidence stack over \(E^I\):
\[
\mathfrak M^{\mathrm{loc,cl}}_{\alpha,R,I}
=
\Bigl\{(A,\mathbf e)\in\mathfrak M^{\mathrm{sc}}_{R,L\alpha}\times E^I
\Bigm|p_E(\mathrm{Supp}\,A)\subset\bigcup_{i\in I}S\times\{e_i\}\Bigr\}.
\]
Form
\(\mathscr H^{\mathrm{loc}}_{\alpha,R,I}=
(\pi_{\alpha,R,I})_!(\Phi_{\alpha,R,I}^{\mathrm{loc}}\otimes
o_{\alpha,R,I}^{\mathrm{loc}})\) and the open-support sheaf
\(\mathcal F_{\alpha,R}^{\mathrm{loc}}(U)\) over \(U\subset E\) by
restricted Borel--Moore configuration sums.

\item \emph{Wrapped prequotient sector.}  For
\(\eta\in\Gamma_R^{\mathrm{wr}}\), form the rigidified prequotient
\(\mathcal M_{\eta,R}^{\mathrm{wr,rig}}\to\mathcal M_{\eta,R}^{\mathrm{wr}}\)
by the Abramovich--Polishchuk / determinant anchor
\[
\lambda_\eta:\mathcal M_{\eta,R}^{\mathrm{wr,rig}}\longrightarrow E,
\qquad
\lambda_\eta(A)=\det Rp_{E,*}A
\otimes\mathcal O_E(-\chi(A)\cdot 0_E),
\]
divided by a finite cover when \(\chi(A)\ne 1\) so that
\(\lambda_\eta(t\cdot A)=\lambda_\eta(A)+t\) holds.  Verify the
losslessness residual \(\mathfrak o^{\lambda,\mathrm{loss}}_R=0\) by
the Ext-distinguishing test of
Definition~\ref{def:hybrid-wrapped-factorization-datum}, separating
\(W_0=i_{s,*}(\mathcal O_E\oplus\mathcal O_E)\) from
\(W_L=i_{s,*}(L\oplus L^{-1})\).

\item \emph{Mixed local/wrapped correspondences.}  For an ordered
hybrid input
\(\alpha_-\in\Gamma_R^{\mathrm{loc}}\),
\(\eta\in\Gamma_R^{\mathrm{wr}}\),
\(\beta_+\in\Gamma_R^{\mathrm{loc}}\), form
\(\mathfrak E^{\mathrm{mix}}_{\alpha_-;\eta;\beta_+,R}\) before the
reduced \(E\)-quotient.

\item \emph{Eight-word two-step flag atlas.}  Build flag stacks
\(\mathfrak F^{(2),w}_R\) for each
\(w\in\{\mathrm L,\mathrm W\}^3=
\{\mathrm{LLL},\mathrm{LLW},\mathrm{LWL},\mathrm{WLL},\mathrm{LWW},
\mathrm{WLW},\mathrm{WWL},\mathrm{WWW}\}\), classifying filtrations
with three associated graded factors of those types.  Verify the
four-input pentagon coherence and the quotient-after-correspondence
descent.

\item \emph{Hybrid correspondence-module functor.}  Assemble the
correspondence category
\(\mathsf{Corr}^{E,\mathrm{hyb}}_R\) and the functor
\(\mathcal B_R^{E,\mathrm{hyb}}:\mathsf{Corr}^{E,\mathrm{hyb}}_R
\to\mathsf{Ch}_{\C}\) by
\(\mathcal B_R^{E,\mathrm{hyb}}(e)=q_!\,\mathrm{TS}^{\mathrm{red}}_e\,p^*\)
for every generating arrow
\(e=(Y\xleftarrow{p}\mathfrak E_e\xrightarrow{q}Z)\).

\item \emph{Quotient functor \(Q_{E,R}\).}  Apply
\(E\)-equivariant compactly supported vanishing-cycle chains and the
orientation descent of step~4 to produce the reduced object
\[
\mathcal F^{\mathrm{geom}}_{X,\sigma,S,\le R}
=Q_{E,R}\circ\mathcal B_R^{E,\mathrm{hyb}}.
\]

\item Output the hybrid datum with all coherences and the protected
integration map \(I_R^{\mathrm{prot}}\), with homogeneous
\(s\)-degree the Borcherds trace exponent
\(m_R=\operatorname{pr}_3\overline\Pi_X\).
\end{enumerate}

\paragraph{Type checks.}
The carrier
\(\mathcal F_{X,\sigma,S,\le R}^{\mathrm{hyb}}\) is a hybrid
correspondence-module functor on
\(\mathsf{Corr}^{E,\mathrm{hyb}}_R\) with values in
\(\mathsf{Ch}_{\C}\); the four-input pentagon is supplied by
the eight-word two-step flag atlas; the anchor residual tuple
\(\mathfrak o^\lambda_R=
(\mathfrak o^{\lambda,\mathrm{ex}}_R,
\mathfrak o^{\lambda,\mathrm{unit}}_R,
\mathfrak o^{\lambda,\mathrm{loss}}_R,
\mathfrak o^{\lambda,\mathrm{multi}}_R,
\mathfrak o^{\lambda,\mathrm{tr}}_{R'R})\) vanishes; the hybrid
residual \(\mathfrak O_{\mathrm{hyb},R}=0\) at every height \(R\),
compatibly with transition maps.

\paragraph{Obstructions discharged or invoked.}
The middle type-II wall lies in
\(\delta_2\leftrightarrow(0,1,1)\), elliptic degree
\(d=m+1=2\) by Lemma~\ref{app:pfaffian-wall-three-walls}; it
cannot be represented by a projection-finite local stack and demands
the wrapped prequotient.  This step's correctness — the existence of
the wrapped representative
\(x_{\delta_2,R}\in\mathcal M_{\eta,R}^{\mathrm{wr,rig}}\) with
\(\overline\Pi_X(x_{\delta_2,R})=(0,1,1)\) and its compatible anchor
\(\lambda_{\eta,R}\) — is precisely the hybrid input to \((O2)\) at
the middle wall.

\subsection{Construction~4 — Reduced d-critical orientation
and Joyce--Upmeier transport}
\label{app:build-orientation}

\paragraph{Tag.} \(\beta\)-functorial passage (oriented carrier).

\paragraph{Input.}
The hybrid carrier
\(\mathcal F_{X,\sigma,S,\le R}^{\mathrm{hyb}}\) of
Construction~\ref{app:build-hybrid-wrapped}; the K3 holomorphic
\(2\)-form \(\sigma_S\) on \(S\); BBDJS d-critical chart machinery
\cite{BBDJS2015}; the Joyce--Upmeier multiplicative orientation
formalism \cite{JoyceUpmeier,JoyceTanakaUpmeier}; the cosection
localization machinery of Kiem--Li and the cosection-reduced
obstruction theory; the type-II Weyl group
\(W^{(2)}(\La^{2,1}_{II})\) generated by
\(s_{\delta_1},s_{\delta_2},s_{\delta_3}\).

\paragraph{Output.}
The cosection
\(\sigma=p_S^*\sigma_S\); the reduced tangent complex
\[
\operatorname{RHom}_{\mathrm{red}}(A,A)
=
\operatorname{Cone}\!\Bigl(
\operatorname{RHom}_X(A,A)
\xrightarrow{(\mathrm{tr},\mathrm{cs})}
R\Gamma(X,\mathcal O_X)\oplus H^2(S,\mathcal O_S)[-1]
\Bigr)[-1];
\]
the BBDJS vanishing-cycle complex \(\Phi_{R,c}\); the reduced
orientation line
\[
o_{R,c}
=\det\operatorname{Ext}^1_{\mathrm{red}}(A,A)^{-1}
\quad\text{on Darboux charts},\qquad
o_{R,c}^{\otimes 2}\simeq\det\operatorname{RHom}_{\mathrm{red}}(A,A);
\]
the Joyce--Upmeier multiplicative transport
\(o_{c\oplus c'}\simeq o_c\otimes o_{c'}\) compatible with the reduced
extension correspondences; the quotient orientation
\(\mathfrak Q^{or}_R\) of
Theorem~\ref{prop:generic-local-sign-orientation-character}; the
Weyl-equivariant lifts
\(\tau_{i,R,S}:o_{R,S}\to s_{\delta_i}^*o_{R,S}\) for \(i=1,2,3\) of
\((O1)^+\); the local generic wall-sign data \(\epsilon_{o,R}^{\mathrm{loc}}\).

\paragraph{Procedure.}
\begin{enumerate}
\item \emph{Cosection.}  Pull back the K3 semiregularity cosection
\(\mathrm{cs}\) along \(p_S:X\to S\) to obtain the 2-step cosection
\(\sigma=p_S^*\sigma_S\) on the K3 factor; verify additivity in exact
triangles by \textup{(RedOr)} of
Theorem~\ref{thm:finite-stage-dirac-igusa-pro-object}.

\item \emph{Reduced obstruction.}  Form the reduced tangent complex by
the displayed cone above (Kiem--Li cosection localization).  PTVV
\cite{PTVV2013} \((-1)\)-shifted symplectic on
\(\operatorname{RHom}_X(A,A)\) descends, after the cosection
contraction, to the reduced complex.

\item \emph{BBDJS d-critical chart.}  On each Darboux chart of
\(\mathfrak M^{\mathrm{sc}}_{R,c}\), build the BBDJS vanishing-cycle
complex \(\Phi_{R,c}\) \cite{BBDJS2015}; verify reduced
Thom--Sebastiani transport coherent on two-step flag stacks of
Construction~\ref{app:build-hybrid-wrapped}.

\item \emph{Reduced orientation line.}  Define
\(o_{R,c}=\det\operatorname{Ext}^1_{\mathrm{red}}(A,A)^{-1}\); verify
the squared isomorphism
\(o_{R,c}^{\otimes 2}\simeq\det\operatorname{RHom}_{\mathrm{red}}(A,A)\)
by \(3\)-Calabi--Yau Serre duality
\(\operatorname{Ext}^2_{\mathrm{red}}(A,A)
\simeq\operatorname{Ext}^1_{\mathrm{red}}(A,A)^\vee\).

\item \emph{Quotient-orientation descent on the K3\(\times\)E-quotient
self-Ext frame bundle.}  Compute the four Cech--Borel obstruction
classes for every retained stratum \(S\):
\[
\alpha^{\mathrm{red}}_{R,S}=w_2(\mathscr D_{R,S}),
\quad
\alpha^{E,\mathrm{free}}_{R,S}=
\mathrm{edge}_{2,0}(w_2^E(\mathscr D_{R,S})),
\quad
\widetilde\beta^H_{R,S}=w_2^G(\mathscr D_{R,S}),
\quad
\lambda^H_{R,S}\in H^1_G(S;\mathbb F_2),
\]
and verify simultaneous vanishing with compatible
null-trivializations on overlaps, on Hall correspondence pullbacks,
on two-step flag stacks, and under successor maps.  The full
calculation proceeds primary stabilizer by stabilizer; for cyclic
\(\Z/2^a\) the absent mixed coefficient requires the rank-two
two-primary check.  Output the descended quotient orientation system
\(\mathfrak Q^{or}_R\).

\item \emph{Multiplicative Joyce--Upmeier transport.}  Transport
\(o_{c\oplus c'}\simeq o_c\otimes o_{c'}\) coherently across two-step
flag stacks
\cite{JoyceUpmeier,JoyceTanakaUpmeier}; verify the reduced
Thom--Sebastiani identification on three-input flags.

\item \emph{Weyl-equivariant lifts \((O1)^+\).}  For each
\(i\in\{1,2,3\}\) construct the lift
\(\tau_{i,R,S}:o_{R,S}\to s_{\delta_i}^*o_{R,S}\); verify the
projective Weyl cocycle
\(c_{o,R}\in Z^2(\mathcal G_R;\underline{\mathbb F}_2)\) on the
retained action groupoid; null-trivialize \([c_{o,R}]\) and choose a
\(1\)-cochain killing it; verify \(\tau_{i,R,S}^2=1\); verify the
Coxeter relations
\((s_{\delta_i}s_{\delta_j})^3=1\) for \(i\ne j\) (the
\((3,3,3)\)-hyperbolic triangle relations of
Appendix~\ref{app:pfaffian-wall-three-walls}); verify that the torsor
defects \(\omega_{i,\mathcal C}\) of
Definition~\ref{thm:finite-stage-dirac-igusa-pro-object}~$(O_X)$ vanish.

\item \emph{Local wall-sign datum.}  At each generic type-II wall point
\(\Hpl_{\delta_i}\), with the rank-one local model
\(K_{\delta_i}^{\mathrm{cross}}=[\R\xrightarrow{u_i}\R]\), record the
local generic sign \(\epsilon_{o,R}^{\mathrm{loc}}(s_{\delta_i})=-1\)
of Proposition~\ref{prop:rank-one-crossing-o2-sign}.

\item Output
\((o_R,\mathfrak Q^{or}_R,\{\tau_{i,R,S}\},\epsilon_{o,R}^{\mathrm{loc}})\).
\end{enumerate}

\paragraph{Type checks.}
\(o_R\) is a Picard-line section squaring to the reduced determinant;
the four Cech--Borel classes vanish with compatible
null-trivializations; the lifts \(\tau_{i,R}\) satisfy
\(\tau_{i,R}^2=1\) and the \((3,3,3)\)-Coxeter relations; multiplicative
transport agrees on flag stacks.

\paragraph{Obstructions discharged or invoked.}
This step is the entire site of \((O1)\) (strong reduced orientation
on the K3\(\times\)E-quotient self-Ext frame bundle) and of
\((O1)^+\) (Weyl-equivariant transport along
\(W^{(2)}(\La^{2,1}_{II})\) reflections).  Discharging \((O1)\)
requires the simultaneous vanishing and compatible
null-trivialization of the four Cech--Borel classes on every retained
stratum; discharging \((O1)^+\) requires \([c_{o,R}]=0\), the
Coxeter coherence, and the vanishing of all torsor defects
\(\omega_{i,\mathcal C}\).  The local wall-sign datum is the input to
\((O2)\) at step~6.

\subsection{Construction~5 — Hall correspondences and primitive part}
\label{app:build-hall}

\paragraph{Tag.} \(\gamma\)-shadow specialization
(Stage-2 chiral Hall shadow).

\paragraph{Input.}
The oriented hybrid carrier from
Construction~\ref{app:build-orientation}; the extension stack
\(\mathfrak E^{\mathrm{geom}}_{R,\widehat c,\widehat c'}\); the
two-step flag stack
\(\mathfrak F^{(2),\mathrm{geom}}_{R,\widehat c,\widehat c',\widehat c''}\);
the BBDJS reduced Thom--Sebastiani transport; six-functor admissibility
of \(p^*,q_!,p_!\) (Construction~\ref{app:build-finite-moduli}~step
\textup{(Six)}).

\paragraph{Output.}
The finite Hall object
\[
\mathcal H_R^{\mathrm{geom}}
=\bigoplus_{\widehat c\in\widehat\Gamma_R}
H_c^\bullet\!\bigl(
\mathfrak M^{\mathrm{geom}}_{R,\widehat c},
\Phi_{R,c}\otimes o_{R,c}\bigr);
\]
the Hall product \(m_R=q_!\,\mathrm{TS}^{\mathrm{red}}\,p^*\) and Hall
coproduct \(\Delta_R=p_!\,(\mathrm{TS}^{\mathrm{red}})^{-1}\,q^*\); the
finite Hall bialgebra
\((\mathcal H_R^{\mathrm{geom}},m_R,\Delta_R,\eta_R,\epsilon_R)\); the
primitive part
\[
P_R^{\mathrm{geom}}
=\ker(\Delta_R-1\otimes\mathrm{id}-\mathrm{id}\otimes 1)\cap\ker\epsilon_R;
\]
the normal-ordered Gram descent
\(P_R^{\Pi,\mathrm{geom}}
=\overline\Pi_{X,*}^{\Theta_R}P_R^{\mathrm{geom}}\); the
\(\Gamma_R^\Pi\)-graded super-Lie bracket
\([\,\bar x,\bar y\,]_{\mathrm{red}}\) of
Proposition~\ref{thm:hall-primitives-formal};
the Hopf pairing matrix \(G_{\gamma,ab}\) and radical \(K_\gamma\).

\paragraph{Procedure.}
\begin{enumerate}
\item \emph{Pullback.}  For
\(\widehat c,\widehat c'\in\widehat\Gamma_R\), pull back along
\(p:\mathfrak E^{\mathrm{geom}}_{R,\widehat c,\widehat c'}
\to\mathfrak M^{\mathrm{geom}}_{R,\widehat c}
\times\mathfrak M^{\mathrm{geom}}_{R,\widehat c'}\); apply reduced
Thom--Sebastiani \(\mathrm{TS}^{\mathrm{red}}\) to combine vanishing
cycle data \cite{BBDJS2015}; pushforward exceptionally along
\(q:\mathfrak E^{\mathrm{geom}}_{R,\widehat c,\widehat c'}
\to\mathfrak M^{\mathrm{geom}}_{R,\widehat c\star\widehat c'}\) to
obtain
\(m_R(x\otimes y)=q_!\,\mathrm{TS}^{\mathrm{red}}\,p^*(x\boxtimes y)\).

\item \emph{Coproduct.}  Dually pull back along \(q\), apply
inverse reduced Thom--Sebastiani, pushforward along \(p\) to obtain
\(\Delta_R(z)=p_!\,(\mathrm{TS}^{\mathrm{red}})^{-1}\,q^*(z)\).
The compact geometric coproduct correspondence
\(\mathfrak C^{\mathrm{geom}}_{\rho;\mu,\nu}\) of
Proposition~\ref{thm:hall-primitives-formal} is the
support data for splitting.

\item \emph{Bialgebra check.}  Verify associativity of \(m_R\) and
coassociativity of \(\Delta_R\) on the two-step flag stack
\(\mathfrak F^{(2),\mathrm{geom}}_{R,\widehat c,\widehat c',\widehat c''}\)
by base change and Thom--Sebastiani coherence.  Verify the bialgebra
identity (Hopf compatibility) via the Massey--Saito Joyce--Upmeier
multiplicative orientation transport.

\item \emph{Primitive extraction.}  Compute
\(P_R^{\mathrm{geom}}=\ker(\Delta_R-1\otimes\mathrm{id}
-\mathrm{id}\otimes 1)\cap\ker\epsilon_R\); since
\(\Delta_R-1\otimes\mathrm{id}-\mathrm{id}\otimes 1\) is degree-graded
on \(\widehat\Gamma_R\), this reduces to a finite linear computation
in each charge.

\item \emph{Normal-ordered descent.}  Apply
\(\overline\Pi_{X,*}^{\Theta_R}\) of
Theorem~\ref{thm:hall-product} to obtain the
\(\Gamma_R^\Pi\)-graded
\(P_R^{\Pi,\mathrm{geom}}\); verify additivity
\(\overline\Pi_X(\widehat c\star\widehat c')
=\overline\Pi_X(\widehat c)+\overline\Pi_X(\widehat c')\).

\item \emph{Bracket and pairing matrices.}  Compute the matrices
\[
M^{\alpha+\beta,c}_{\alpha,a;\beta,b}
=\int_{\mathfrak E^{\mathrm{geom}}_{\alpha,\beta;\alpha+\beta}}
q^!(e_{\alpha+\beta,c}^\vee)\cap
\mathrm{TS}^{\mathrm{red}}(p^*(e_{\alpha,a}\boxtimes e_{\beta,b})),
\quad
B_{\alpha,\beta}(x\otimes y)
=M_{\alpha,\beta}(x\otimes y)
-(-1)^{|x||y|}M_{\beta,\alpha}\mathsf{sw}(x\otimes y),
\]
\[
G_{\gamma,ab}=
\mathrm{tr}_0(q_!\,\mathrm{TS}^{\mathrm{red}}\,p^*
(e_{\gamma,a}\boxtimes e_{-\gamma,b}))
\quad\hbox{on}\quad
\mathfrak P^{\mathrm{geom}}_{\gamma,-\gamma;0}.
\]
Form the radical
\(K_{\gamma,\bar p}
=\ker G_{\gamma,\bar p}\cap\ker G_{-\gamma,\bar p}^t\) and the
quotient \(L_{\gamma,\bar p}\); fix splittings.

\item Output
\((\mathcal H_R^{\mathrm{geom}},m_R,\Delta_R,P_R^{\mathrm{geom}},
P_R^{\Pi,\mathrm{geom}},
\{B_{\alpha,\beta}\},\{G_{\gamma,ab}\},\{K_\gamma\},\{L_\gamma\})\).
\end{enumerate}

\paragraph{Type checks.}
\(\mathcal H_R^{\mathrm{geom}}\) is a finite \(\widehat\Gamma_R\)-graded
Hall bialgebra; \(m_R\) is associative on the two-step flag stack;
\(\Delta_R\) is coassociative; the supercommutator of primitives is
primitive (the bialgebra identity); the descended bracket is
\(\Gamma_R^\Pi\)-graded after \(\overline\Pi_{X,*}^{\Theta_R}\);
finite-dimensional in each homogeneous parity piece by
finite-typeness.

\paragraph{Obstructions discharged or invoked.}
\((D_0)\) is partially discharged: the finite Hall package
\((\mathcal H_R^{\mathrm{geom}},m_R,\Delta_R,P_R)\) is constructed.
The seven obstruction classes
\(\mathfrak o_\Theta^{\mathrm{top}},
\mathfrak o_\Theta^{\mathrm{Hoch}},
\mathfrak o_\Theta^{\mathrm{tri}},
\mathfrak o_\Theta^{\mathrm{Jac}},
\mathfrak o_F,
\mathfrak o_\Theta^{\mathrm{cyc}},
\mathfrak o_\Theta^{\mathrm{rad}}\) of
Theorem~\ref{thm:hall-product} are required to vanish on
this finite Hall stage and compatibly in the inverse limit.

\paragraph{Cross-volume anchor.}  The output
\((\mathcal H_R^{\mathrm{geom}},m_R)\) is the finite-stage realization
of the K3-side BKM Hall--Drinfeld double
\(\mathcal D_\hbar(\mathrm{CoHA}_{K3\times E})\) of
\texttt{calabi-yau-quantum-groups:CLAUDE.md~§6.7}.  In the universal
stage chain \(\mathsf P\to\mathsf C\to\mathsf S\) the Hall positive half
\(Y^+(X)\) is precisely \(P_R^{\mathrm{geom}}\) at finite \(R\); the
Drinfeld double \(G(X)=D(Y^+(X))\) requires Hopf pairing, completion,
integral form, and stable-envelope transport which are supplied by
\(G_{\gamma,ab}\) and \((\mathrm{ML})\).  The κ-tuple coordinate
\(\kappa_{\mathrm{BKM}}=5\) is the weight of \(\Delta_5\) and the
output of the level-4 scalar trace; it is not a level-2 invariant of
the carrier itself.

\subsection{Construction~6 — Dirac operator and Pfaffian line}
\label{app:build-dirac-pfaffian}

\paragraph{Tag.} \(\gamma\)-shadow specialization
(orientation gerbe + Pfaffian).

\paragraph{Input.}
The oriented Hall package
\((\mathcal H_R^{\mathrm{geom}},o_R,\{\tau_{i,R,S}\})\) of
Constructions~\ref{app:build-orientation}--\ref{app:build-hall}; the
local wall-sign datum
\(\epsilon_{o,R}^{\mathrm{loc}}\) of
Construction~\ref{app:build-orientation}; the rank-one type-II
Pfaffian crossing datum of
Definition~\ref{def:rank-one-type-ii-crossing} and
Appendix~\ref{app:pfaffian-wall-rank-one}; the half-Hilbert orbit
under \(W^{(2)}(\La^{2,1}_{II})\) of size \(2^{12}=4096\).

\paragraph{Output.}
The Dirac operator on the orientation gerbe
\(D_R\); the Pfaffian line
\(\mathscr L_{\mathrm{Pf},R}=
\det\operatorname{Ext}^1_{\mathrm{red}}(A,A)^{-1}\) (squaring to the
reduced determinant); the protected Pfaffian section
\(\operatorname{pf}_{X,R}\); the global wall-sign character
\(\epsilon_{o,R}:W^{(2)}(\La^{2,1}_{II})\to\{\pm1\}\); the Pfaffian
section asymptotic
\[
\operatorname{pf}_{X,R}
=
64q^{1/2}r^{1/2}s^{1/2}
\prod_{\gamma=(n,l,m)\in A_R}
(1-q^nr^ls^m)^{f(nm,l)}+\hbox{higher-window correction},
\]
which exhausts to \(\Delta_5\) as \(R\to\infty\).

\paragraph{Procedure.}
\begin{enumerate}
\item \emph{Dirac operator.}  On each Darboux chart of
\(\mathfrak M^{\mathrm{sc}}_{R,c}\), with reduced normal self-Ext
splitting \(K^{\mathrm{nor}}_{R,c}\), build the odd Dirac operator
\(D_{R,c}\) acting on \(\Omega^{\mathrm{sp}}\otimes o_{R,c}\)
(spinors on the orientation gerbe); verify that, at a generic
type-II wall point, the local rank-one normal block is
\(D_{\delta_i}=\begin{psmallmatrix}0&u\\-u&0\end{psmallmatrix}\) with
\(\operatorname{Pf}(D_{\delta_i})=u\)
(Definition~\ref{def:rank-one-type-ii-crossing}).

\item \emph{Pfaffian line.}  Set
\(\mathscr L_{\mathrm{Pf},R}=
\det\operatorname{Ext}^1_{\mathrm{red}}(A,A)^{-1}\); verify
\(\mathscr L_{\mathrm{Pf},R}^{\otimes 2}\simeq
\det\operatorname{RHom}_{\mathrm{red}}(A,A)\) by
\(3\)-Calabi--Yau Serre duality.

\item \emph{Local wall sign.}  At each generic point of
\(\Hpl_{\delta_i}\) with normal model
\(K^{\mathrm{nor}}_{\delta_i,R}=[\R\xrightarrow{u_i}\R]\), compute
the monodromy of the Pfaffian section
\(\operatorname{pf}_{\delta_i,R}=\upsilon_iu_i\) along the
wall-reflection loop \(\ell_{\delta_i,R,S}\) closed by the
\((O1)^+\) lift \(\tau_{i,R,S}\):
\[
M_{\ell_{\delta_i,R,S}}(\upsilon_iu_i)
=\upsilon_i(-u_i)=-\upsilon_iu_i,
\qquad
\epsilon_{o,R}^{\mathrm{loc}}(s_{\delta_i})
=(-1)^{N^{\mathrm{Pf}}_{\delta_i,R}}=-1.
\]

\item \emph{Half-Hilbert orbit transport.}  The reflection
\(s_{\delta_i}\) acts on the half-Hilbert orbit (compact source labels
of size \(2^{12}=4096\) parametrized by 12 binary parities tied to the
\(2^{12}\) odd-spin structures on the lattice
\(\La^{2,1}_{II}/2\La^{2,1}_{II}\)).  Sum the local rank-one
contributions over the orbit:
\[
\epsilon_{o,R}(s_{\delta_i})
=\prod_{x\in\hbox{orbit}/2^{12}}
\epsilon_{o,R}^{\mathrm{loc}}(s_{\delta_i};x).
\]
The output is the global character
\(\epsilon_{o,R}:W^{(2)}(\La^{2,1}_{II})\to\{\pm 1\}\).

\item \emph{Compatibility package.}  Verify the component, boundary,
overlap, quotient, and protected integration compatibility package of
Proposition~\ref{prop:appE-local-pfaffian-sign} on every
retained type-II wall stratum.

\item \emph{Pfaffian section.}  Define
\(\operatorname{pf}_{X,R}\) as the protected integration of the
oriented Hall product over
\(\mathcal F_{X,\sigma,S,\le R}^{\mathrm{hyb}}\):
\[
\operatorname{pf}_{X,R}
=
64q^{1/2}r^{1/2}s^{1/2}
I_R^{\mathrm{prot}}\!\left(
\bigotimes_{\gamma\in A_R}\operatorname{pf}_{\gamma,R}\right),
\qquad
\operatorname{pf}_{\gamma,R}=(1-q^nr^ls^m)^{f(nm,l)}.
\]

\item Output
\((D_R,\mathscr L_{\mathrm{Pf},R},\operatorname{pf}_{X,R},
\epsilon_{o,R})\), with the protected integration compatibility data.
\end{enumerate}

\paragraph{Type checks.}
\(\mathscr L_{\mathrm{Pf},R}\) is a Picard-line whose square
isomorphism is supplied by reduced Serre duality;
\(\operatorname{pf}_{X,R}\) is a section of
\(\mathscr L_{\mathrm{Pf},R}\) with the displayed cusp expansion; the
character \(\epsilon_{o,R}\) takes values in \(\{\pm 1\}\) and
satisfies the Coxeter relations of
\(W^{(2)}(\La^{2,1}_{II})\); on each simple type-II reflection
\(\epsilon_{o,R}(s_{\delta_i})=-1\) by the rank-one local model.

\paragraph{Obstructions discharged or invoked.}
This step is the entire site of \((O2)\) (local Pfaffian wall normal
form on type-II walls; the \(2^{12}=4096\) sign sum on the
half-Hilbert orbit transport).  The local rank-one Pfaffian crossing
of Appendix~\ref{app:pfaffian-wall-rank-one} discharges \((O2)\)
generically; the global half-Hilbert orbit transport remains the open
\((O2)\) component of
Open~Problem~\ref{def:O1-plus-weyl-transport}.  The character
identity
\(\epsilon_o=\nu_{\Delta_5}\big|_{W^{(2)}(\La^{2,1}_{II})}\) becomes a
theorem only after \((D_0)\), \((O1)\), \((O1)^+\), and \((O2)\) are
discharged compatibly in \(R\).

\subsection{Construction~7 — Chiral Koszul source coalgebra}
\label{app:build-koszul}

\paragraph{Tag.} \(\gamma\)-shadow specialization
(\emph{Bar = twisting; not bulk}; the bulk is
\(Z^{\mathrm{der}}_{\mathrm{ch}}(\mathcal F_X^{\mathrm{hyb}})\),
constructed elsewhere).

\paragraph{Input.}
The oriented Hall primitives
\(P_R^{\Pi,\mathrm{geom}}\) of
Construction~\ref{app:build-hall}; the hybrid carrier
\(\mathcal F_{X,\sigma,S,\le R}^{\mathrm{hyb}}\) of
Construction~\ref{app:build-hybrid-wrapped}; the augmentation
\(\varepsilon_{X,R}^{\mathrm{hyb}}:
\mathcal F_{X,\sigma,S,\le R}^{\mathrm{hyb}}\to k\mathbf 1\); the
hybrid residual \(\mathfrak O_{\mathrm{hyb},R}=0\); the bounded
length \(N_R\) (positivity of HN height on every non-vacuum colour);
the Beilinson--Drinfeld/Francis--Gaitsgory bar-cobar formalism
\cite{BD,FrancisGaitsgoryCKD}; the conilpotent completed domain on
the elliptic curve \(E\).

\paragraph{Output.}
The conilpotent chiral coalgebra on \(E\):
\[
C_{X,R}^{\mathrm{geom}}
=\bar B_E^{\mathrm{ch}}
(\mathcal F_{X,\sigma,S,\le R}^{\mathrm{hyb}},
\varepsilon_{X,R}^{\mathrm{hyb}})
=k\mathbf 1\oplus\bigoplus_{1\le p\le N_R}
\bar B_E^{\mathrm{ch},p}(\overline{\mathcal F}_{X,R}^{\mathrm{geom}});
\]
the bar-length filtration
\(F^{\mathrm{co}}_{R,p}\); the chiral collision coproduct
\(\Delta_R^{\mathrm{ch}}\); the bar comparison
\(b_{X,R}=\mathrm{id}_{\bar B_E^{\mathrm{ch}}(\cdot)}\); the source
cobar counit
\(\varepsilon^{\mathrm{src}}_{\mathrm{bar},R}:
\Omega_E^{\mathrm{ch}}C_{X,R}^{\mathrm{geom}}
\xrightarrow{\simeq}
\mathcal F_{X,\sigma,S,\le R}^{\mathrm{hyb}}\); after a separate
source-to-target quasi-isomorphism
\(\vartheta_R:\mathcal F_{X,\sigma,S,\le R}^{\mathrm{hyb}}
\xrightarrow{\simeq}\mathsf{Den}_{\Delta_5,E,\le R}\), the cobar
comparison
\(\Theta_{\mathrm{Kos},R}=\vartheta_R\circ\varepsilon^{\mathrm{src}}_{\mathrm{bar},R}\).

\paragraph{Procedure.}
\begin{enumerate}
\item \emph{Augmentation kernel.}  Set
\(\overline{\mathcal F}_{X,R}^{\mathrm{geom}}
=\ker(\varepsilon_{X,R}^{\mathrm{hyb}})\); verify the bounded-length
hypothesis (positivity of HN height on every non-vacuum surviving
colour, hence at most \(N_R\) factors in any non-zero word).

\item \emph{Bar coalgebra.}  Form the reduced chiral bar
\[
C_{X,R}^{\mathrm{geom}}
=k\mathbf 1\oplus\bigoplus_{1\le p\le N_R}
\bar B_E^{\mathrm{ch},p}
(\overline{\mathcal F}_{X,R}^{\mathrm{geom}}),
\]
length \(p\) summand generated by ordered \(p\)-fold inputs from the
augmentation ideal of total HN height \(\le R\) (height \(>R\) outputs
zero); coaugmentation by the inclusion of the vacuum summand; counit by
projection.

\item \emph{Bar-length filtration.}
\[
F^{\mathrm{co}}_{R,p}C_{X,R}^{\mathrm{geom}}
=k\mathbf 1\oplus\bigoplus_{1\le j\le p}
\bar B_E^{\mathrm{ch},j}
(\overline{\mathcal F}_{X,R}^{\mathrm{geom}}),\qquad
0\le p\le N_R.
\]

\item \emph{Collision coproduct.}  For
\(x=[x_1|\cdots|x_p]\), define the cut-map coproduct
\[
\Delta_R^{\mathrm{ch}}(x)
=\mathbf1\otimes x+x\otimes\mathbf1
+\sum_{i=1}^{p-1}[x_1|\cdots|x_i]\otimes[x_{i+1}|\cdots|x_p],
\]
with each cut realized by the corresponding finite collision/flag
pull-push map in
\(\mathsf{Corr}^{E,\mathrm{hyb}}_{R,\mathrm{geom}}\) of
Construction~\ref{app:build-hybrid-wrapped}.  Verify
coassociativity by the four-input pentagon and the eight-word flag
coherences; verify counit identities via vacuum collision flags.

\item \emph{Conilpotence.}  Verify
\((\overline\Delta_R^{\mathrm{bar}})^{N_R+1}=0\): each reduced cut
strictly lowers bar length on each branch.

\item \emph{Bar comparison.}  Take \(b_{X,R}\) to be the identity of
\(\bar B_E^{\mathrm{ch}}
(\mathcal F_{X,\sigma,S,\le R}^{\mathrm{hyb}},
\varepsilon_{X,R}^{\mathrm{hyb}})\), eliminating the coalgebra-side
finite Koszul residuals
\(\mathfrak o^C_R, \mathfrak o^{\mathrm{fil}}_R,
\mathfrak o^{\mathrm{conil}}_R, \mathfrak o^{\Delta,\mathrm{ch}}_R\)
of Definition~\ref{def:chiral-koszul-source-datum}.

\item \emph{Source cobar counit.}  Apply Francis--Gaitsgory chiral
Koszul duality
\cite[Theorem~6.2.2]{FrancisGaitsgoryCKD} in the conilpotent completed
domain to obtain
\(\varepsilon^{\mathrm{src}}_{\mathrm{bar},R}:
\Omega_E^{\mathrm{ch}}C_{X,R}^{\mathrm{geom}}
\xrightarrow{\simeq}
\mathcal F_{X,\sigma,S,\le R}^{\mathrm{hyb}}\); verify it is a
source-internal quasi-isomorphism of chiral algebras on \(E\), not a
homotopy inverse of the target counit.

\item \emph{Source-to-target identification.}  Construct
\(\vartheta_R\) from the geometric primitive representatives, Hall
products, coproducts, pairings, and HN transition compatibility;
verify the residual
\(\mathfrak o^\vartheta_R\) of
Corollary~\ref{cor:source-bar-remaining-cobar-obstruction} vanishes;
set \(\Theta_{\mathrm{Kos},R}=\vartheta_R\circ
\varepsilon^{\mathrm{src}}_{\mathrm{bar},R}\).

\item \emph{Inverse limit.}  Verify the Koszul Mittag--Leffler
condition \(R^1\varprojlim=0\) for the inverse systems of source
coalgebras, completed chiral bar/cobar constructions, target
truncations, and the cones of \(b_{X,R}\) and
\(\Theta_{\mathrm{Kos},R}\) \cite[\S3.5]{WeibelHA}.  Form
\[
C_X^{\mathrm{geom}}
=\varprojlim_R C_{X,R}^{\mathrm{geom}},\qquad
\Theta_{\mathrm{Kos}}:
\Omega_E^{\mathrm{ch}}C_X^{\mathrm{geom}}
\xrightarrow{\simeq}\mathsf{Den}_{\Delta_5,E}.
\]

\item Output
\((C_{X,R}^{\mathrm{geom}},F^{\mathrm{co}}_{R,\bullet},
\Delta_R^{\mathrm{ch}},b_{X,R},\Theta_{\mathrm{Kos},R})\) at every
\(R\), with successor compatibility.
\end{enumerate}

\paragraph{Type checks.}
\(C_{X,R}^{\mathrm{geom}}\) is a coaugmented conilpotent chiral
coalgebra on \(E\) with bounded bar-length filtration of length at
most \(N_R\); \(\Delta_R^{\mathrm{ch}}\) is coassociative and counital
on the chosen flag stacks; the bar-cobar inversion is supplied by
Francis--Gaitsgory chiral Koszul; the cobar comparison
\(\Theta_{\mathrm{Kos},R}\) is source-internal, not the target
bar-cobar counit (the source/target separation of
Lemma~\ref{lem:source-target-koszul-separation}).

\paragraph{Obstructions discharged or invoked.}
The Koszul Mittag--Leffler exactness hypothesis is the
\textup{(D0)}-component for the chiral Koszul lane.  The total finite
Koszul residual
\(\mathfrak O_{\mathrm{Kos},R}=
(\mathfrak o^C_R,\mathfrak o^{\mathrm{fil}}_R,
\mathfrak o^{\mathrm{conil}}_R,\mathfrak o^{\Delta,\mathrm{ch}}_R,
\{\mathfrak o^{C,\mathrm{tr}}_{R'R}\},\mathfrak o^b_R,
\mathfrak o^\Omega_R,\mathfrak o^{\mathrm{hyb}}_R,
\mathfrak o^\Pi_R,\mathfrak o^{\mathrm{sep}}_R,
\mathfrak o^{\Prim}_R)\)
vanishes after this step.

\paragraph{Cross-volume firewall.}  The bar coalgebra
\(\bar B_E^{\mathrm{ch}}
(\mathcal F_{X,R}^{\mathrm{hyb}})\) is a \emph{twisting / coupling}
coalgebra (Maurer--Cartan classification), \emph{not} the bulk
\(Z^{\mathrm{der}}_{\mathrm{ch}}(\mathcal F_X^{\mathrm{hyb}})
=\mathrm{ChirHoch}^\bullet(\mathcal F_X^{\mathrm{hyb}},
\mathcal F_X^{\mathrm{hyb}})\); cf.\
\texttt{chiral-bar-cobar-vol2:CLAUDE.md~§6} (universal stage chain
\(\mathsf P\to\mathsf C\to\mathsf S\to\mathsf Z\)).
The chiral Koszul comparison
\(\Theta_{\mathrm{Kos}}:\Omega_E^{\mathrm{ch}}C_X^{\mathrm{geom}}
\xrightarrow{\simeq}\mathsf{Den}_{\Delta_5,E}\) is a comparison
between two algebras of the same level; the homotopy inverse to the
target-internal counit
\(\Omega_E^{\mathrm{ch}}\bar B_E^{\mathrm{ch}}
(\mathsf{Den}_{\Delta_5,E})\to\mathsf{Den}_{\Delta_5,E}\) is not used
in either direction.

\subsection{Construction~8 — Primitive recognition}
\label{app:build-recognition}

\paragraph{Tag.} \(\delta\)-endpoint convergence
(source-to-target identification).

\paragraph{Input.}
The descended source primitive Lie superalgebra
\(P_X^{\Pi,\mathrm{red}}\) of
Construction~\ref{app:build-hall} (after the radical quotient
\(\overline{\mathrm{Rad}\langle\,\cdot,\cdot\rangle_X}\)); the imported
target denominator algebra
\(\fg_{\Delta_5}=G(\mathscr B_{\Delta_5})\) on
\(\La^{2,1}_{II}\) of Definition~\ref{def:bkm-imaginary-simple-roots} and
Definition~\ref{def:bkm-relations}; the increasing cofinal
sequence of finite downward-saturated windows
\(W_1\subset W_2\subset\cdots\) with
\(\bigcup_\nu W_\nu=\mathcal R_+\); the Borcherds--Kac--Moody data
\((H,I,p,\alpha_i,(\,,\,))\) of
Theorem~\ref{thm:primitive-recognition}.

\paragraph{Output.}
The cofinal finite-window primitive-recognition datum
\(\{\textup{(R0),(R1),(R2),(R3),(R4),(R5),(R6),(R7),(R8)}\}\) of
Definition~\ref{def:cofinal-primitive-recognition-datum} for every
\(W_\nu\), compatible under \(W_\nu\subset W_{\nu+1}\); the source
representatives
\(e_i^X,f_i^X,h_i^X\); the source presentation map
\(\pi_W:\widehat{\mathfrak F}(\mathscr B_{\Delta_5})_W
\to P_X^{\Pi,\mathrm{red}}\big|_W\); the kernel equality
\(\ker\pi_W=
(\overline{J_{\mathrm{BK}}+\mathrm{Rad}_{\mathrm{GN}}})_W\); the
isomorphism
\[
\boxed{\quad
P_X^{\Pi,\mathrm{red}}\cong\fg_{\Delta_5},\quad
\Omega_E^{\mathrm{ch}}C_X\simeq\mathsf{Den}_{\Delta_5,E}.
\quad}
\]

\paragraph{Procedure.}
For every cofinal \(W_\nu\):
\begin{enumerate}
\item[\textup{(R0)}] \emph{Compact source labels.}  Verify every
finite source basis vector in \(W_\nu\cup(-W_\nu)\cup\{0\}\) satisfies
the compact source-admissibility condition of
Theorem~\ref{thm:finite-stage-dirac-igusa-pro-object}.

\item[\textup{(R1)}] \emph{Representatives and parities.}  Supply
parity-homogeneous primitive representatives
\(e_i^X\in P^{\Pi,\mathrm{red}}_{X,\alpha_i}\),
\(f_i^X\in P^{\Pi,\mathrm{red}}_{X,-\alpha_i}\),
\(h_i^X=\iota_H^{-1}(\alpha_i)\) for every simple-root index \(i\) with
\(\alpha_i\in W_\nu\), with the Gritsenko--Nikulin parity fibres
\(\mu_{\overline 0}(a),\mu_{\overline 1}(a)\) for imaginary simple
roots and even real representatives for \(\delta_1,\delta_2,\delta_3\).

\item[\textup{(R2)}] \emph{Relations.}  Verify in the protected Hall
super-bracket the Cartan, Chevalley, real Serre, Borcherds
orthogonality, and super sign relations for every relation whose
displayed source and target degrees lie in
\(W_\nu\cup(-W_\nu)\cup\{0\}\).  In the real rank-three subsystem the
Cartan matrix is
\(((\delta_i,\delta_j))=\mathrm{diag}(2)-2(J-\mathrm{diag}(1))\) with
real Serre exponent~3.

\item[\textup{(R3)}] \emph{Full finite parity dimensions.}  For every
\(\beta\in W_\nu\) verify
\(\dim P^{\Pi,\mathrm{red}}_{X,\beta,\overline p}
=\rootmult_{\overline p}^{\mathrm{GN}}(\beta)\) for both parities
\(p=0,1\).  In the first timelike channel
\(\beta=\delta_{123}\),
Proposition~\ref{prop:bkm-delta123-presentation-count} requires
\((d^{\mathrm{GN}}_{(1,1,1),\overline 0},
d^{\mathrm{GN}}_{(1,1,1),\overline 1})=(29,93)\).  Determinant
information \(29-93=-64=f(1,1)\) is consistent but not the
recognition.

\item[\textup{(R4)}] \emph{Hopf pairing and radical.}  Compute the
positive-negative pairing matrices in parity blocks; verify the
kernel agrees with the
Gritsenko--Nikulin/Kac invariant-pairing kernel; verify Lie-ideal and
coideal stability under all finite brackets; verify carrying under
HN transitions to the next window.

\item[\textup{(R5)}] \emph{No extra relations.}  Verify the kernel
equality
\(\ker\pi_\nu=(\overline{J_{\mathrm{BK}}+\mathrm{Rad}_{\mathrm{GN}}})_{\le W_\nu}\) on the finite triangular span generated by
\(W_\nu\cup(-W_\nu)\cup\{0\}\).

\item[\textup{(R6)}] \emph{Generation.}  Verify that every source
primitive root space in \(W_\nu\cup(-W_\nu)\), after the finite
radical quotient, is generated by the Cartan and the supplied
simple-root representatives through brackets whose intermediate
positive degrees stay in \(W_\nu\).

\item[\textup{(R7)}] \emph{Completed PBW compatibility.}  Verify the
PBW-graded isomorphism
\(\mathrm{gr}_{\mathrm{PBW}}U(P_X^{\Pi,\mathrm{red}})_{\le W_\nu}
\xrightarrow{\sim}
\mathrm{gr}_{\mathrm{PBW}}U(\fg_{\Delta_5})_{\le W_\nu}\); verify
strict preservation under \(W_{\nu+1}\to W_\nu\) transition maps.

\item[\textup{(R8)}] \emph{Exact triangular completion.}  Verify
strictness of the transition maps and closedness of images on every
finite window, so that passage to the inverse limit commutes with the
finite kernels, cokernels, radicals, PBW associated gradeds, and
parity pieces.  Equivalently the relevant \(\lim^1\) terms vanish for
these triangular finite-window systems.

\item Output the global isomorphism
\(P_X^{\Pi,\mathrm{red}}\cong\fg_{\Delta_5}\) by passage to the inverse
limit (compatible cofinal datum); together with the chiral Koszul
comparison of Construction~\ref{app:build-koszul} this gives
\(\Omega_E^{\mathrm{ch}}C_X\simeq\mathsf{Den}_{\Delta_5,E}\) and hence
the protected Pfaffian identification
\(\operatorname{Pf}_{\mathrm{prot}}(\mathfrak D_X^{\mathrm{DI}})=\Delta_5\).
\end{enumerate}

\paragraph{Type checks.}
The source presentation map \(\pi:\widehat{\mathfrak F}(\mathscr B_{\Delta_5})\to P_X^{\Pi,\mathrm{red}}\) is
continuous and degree-preserving; \textup{(R5)} gives injectivity
after the completed radical quotient; \textup{(R6)} gives
surjectivity; \textup{(R7)} gives PBW compatibility; \textup{(R8)}
gives exactness in the inverse-limit step.  No claim on signed
superdimensions \(\sdim=f(nm,l)\) substitutes for the two-integer
parity dimensions of \textup{(R3)}.

\paragraph{Obstructions discharged or invoked.}
This step inherits the obstruction tower discharged by
Constructions~1--7 and adds the cofinal exact-completion residual
\textup{(R8)}.  The Beilinson--Drinfeld bar-cobar firewall is
respected throughout: \textup{(R5)} is a source kernel-equality
theorem, not an inversion of the target bar-cobar counit;
\textup{(R7)} is a PBW graded comparison of source and target
enveloping algebras, not a graph-equality argument from the scalar
square.

\paragraph{Cross-volume anchor.}  The output is
\texttt{chiral-bar-cobar-vol2:CLAUDE.md~§9} Construction
Problem~1: the operator-level construction of
\(\mathfrak D_X^{\mathrm{DI}}\) for \(K3\times E\) with
\(\operatorname{Pf}_{\mathrm{prot}}(\mathfrak D_X^{\mathrm{DI}})=\Delta_5\), under the
four obstructions \((D_0),(O1),(O1)^+,(O2)\).  The Drinfeld doubling
\(Y^+(X)\to G(X)=D(Y^+(X))\) of
\texttt{calabi-yau-quantum-groups:CLAUDE.md} requires the Hopf
pairing of \textup{(R4)}, completion by \textup{(R8)}, integral form,
and stable-envelope transport which are supplied by the same
data.  The κ-tuple coordinate \(\kappa_{\mathrm{BKM}}=5\) is the
weight of \(\Delta_5\) and lives at level~4 of the universal stage
chain; it is a consequence of the construction, not an input to it.

\subsection{The eight-step audit}
\label{app:build-audit}

The eight construction steps form a strict ordered sequence of
input/output type discipline.  An audit of any candidate construction of
\(\mathfrak D_X^{\mathrm{DI}}\) reduces to verifying, for every
\(R\in\mathcal R_{\mathrm{HN}}\):
\begin{enumerate}
\item the lattice combinatorics of
Construction~\ref{app:build-charge-window};
\item the finite-type quasi-smooth derived stacks of
Construction~\ref{app:build-finite-moduli};
\item the hybrid local/wrapped factorization datum and the eight-word
two-step flag atlas of Construction~\ref{app:build-hybrid-wrapped};
\item the four Cech--Borel orientation classes and the
\((O1)^+\)~Coxeter coherence of
Construction~\ref{app:build-orientation};
\item the finite Hall bialgebra and primitive matrices of
Construction~\ref{app:build-hall};
\item the rank-one Pfaffian crossing and the half-Hilbert orbit
transport of Construction~\ref{app:build-dirac-pfaffian};
\item the bar coalgebra, conilpotent collision coproduct, and the
Francis--Gaitsgory cobar quasi-isomorphism of
Construction~\ref{app:build-koszul};
\item the cofinal datum
\(\textup{(R0)}\)--\(\textup{(R8)}\) of
Construction~\ref{app:build-recognition};
\end{enumerate}
together with the compatible vanishing of the obstruction tuples
\((\mathfrak O^{E_3-\mathrm{src}}_R,
\mathfrak O_{\mathrm{hyb},R},
\mathfrak O_{\mathrm{Kos},R},
\mathfrak O^{\mathrm{O1}}_R,
\mathfrak O^{\mathrm{O1}^+}_R,
\mathfrak O^{\mathrm{O2}}_R,
\mathfrak O_R^{D_0})\) at every height \(R\), and the
Mittag--Leffler exactness condition \textup{[ML]} of
Definition~\ref{def:O1-strong-reduced-orientation} on the cofinal HN inverse
system.  Each obstruction is attributable to the unique step
in which its components are constructed.

\paragraph{Order rigidity.}  Each of the seven adjacent input/output
dependencies fails type-checking under exchange:
\begin{enumerate}
\item charge window $\to$ finite moduli: $\Lambda_{\mathrm{ample}}$
  parametrises the HN strata; without the window, the moduli stack
  has no finite cofinal subsystem.
\item finite moduli $\to$ hybrid carrier: $\mathrm{Ran}^{\mathrm{hyb}}(E)$
  needs the finite Hall stack as factor on which to graft the
  wrapped/local atlas.
\item hybrid carrier $\to$ reduced orientation: the cosection-twisted
  Joyce--Upmeier line bundle is defined on the carrier, not on the
  raw moduli (the cosection lives on $\mathrm{Ran}^{\mathrm{hyb}}$).
\item reduced orientation $\to$ Hall correspondences: the Pfaffian
  line bundle on the Hall stack is the orientation-twisted sign data;
  without orientation, the Hall product is unsigned.
\item Hall correspondences $\to$ Pfaffian line + Dirac operator: the
  primitive Hall bracket $[-,-]_X$ defines the Hall--Drinfeld
  correspondence whose Pfaffian section is the Dirac--Igusa Pfaffian.
\item Pfaffian line $\to$ chiral Koszul source coalgebra: the source
  $C_X$ is the bar of the augmented hybrid factorization algebra
  pulled back along the Pfaffian comparison
  $\iota_{\mathrm{aut}}$.
\item chiral Koszul source $\to$ primitive recognition: the primitive
  part $P_X^{\Pi,\mathrm{red}}$ is the Frobenius--Serre primitive
  filtration on the source, which the recognition theorem matches to
  $\mathfrak g_{\Delta_5}$.
\end{enumerate}
Two diagnostic exchanges: orientation before hybrid leaves the middle
type-II wall \(\delta_2\leftrightarrow(0,1,1)\) without a
representative, by
Lemma~\ref{app:pfaffian-wall-three-walls}; recognition before the
Koszul cobar leaves the identification
\(\Omega_E^{\mathrm{ch}}C_X\simeq\mathsf{Den}_{\Delta_5,E}\)
unconstructed.  All other adjacent exchanges fail by the same
input/output type criterion.

\paragraph{Conditional theorem.}  Granted the eight constructions and
the four obstructions \((D_0),(O1),(O1)^+,(O2)\),
Theorem~\ref{thm:pfaffian-dirac-finite-stage} reads
\[
\operatorname{Pf}_{\mathrm{prot}}(\mathfrak D_X^{\mathrm{DI}})=\Delta_5,\qquad
\epsilon_o=\nu_{\Delta_5}\big|_{W^{(2)}(\La^{2,1}_{II})},\qquad
P_X^{\Pi,\mathrm{red}}\cong\fg_{\Delta_5}.
\]
The scalar square
\(Z_{\mathrm{BPS}}^{K3\times E}=
(\Phi_{10}^{\mathrm{un}})^{-1}=\Delta_5^{-2}\) is the level-4 trace
of the constructed protected algebra; promotion to the 3d
gravitational path integral is not claimed and requires the
Hall--Borcherds residual of
\texttt{chiral-bar-cobar-vol2}.

%%% End of Appendix F.
 after
% \appendix has been opened at main.tex line 17081, so that
% \thesection = \Alph{section}.  No \chapter, no second \appendix.

\section{Eight build algorithms for the Dirac--Igusa pro-object}
\label{app:eight-build-algorithms}

The Dirac--Igusa pro-object
\[
\mathfrak D_X^{\mathrm{DI}}
=\varprojlim_R
\bigl(
\mathcal A^{E_3}_{X,R},
\mathcal F_{X,R}^{\mathrm{hyb}},
\Gamma_{X,R},
\Pi_X,
\Phi_R,
o_R,
\mathcal H_R,
P_R^\Pi,
C_{X,R},
\Theta_{\mathrm{Kos},R},
\mathscr L_{\mathrm{Pf},R},
\operatorname{pf}_{X,R},
\epsilon_{o,R},
\mathrm{Rec}_R
\bigr)
\]
of Definition~\ref{thm:finite-stage-dirac-igusa-pro-object} and
Definition~\ref{def:O1-strong-reduced-orientation} is the conditional source
operator-level package whose protected Pfaffian, after the four
obstructions \((D_0)\), \((O1)\), \((O1)^+\), \((O2)\) are discharged,
realizes the Igusa cusp form,
\[
\operatorname{Pf}_{\mathrm{prot}}(\mathfrak D_X^{\mathrm{DI}})=\Delta_5,\qquad
\epsilon_o=\nu_{\Delta_5}\big|_{W^{(2)}(\La^{2,1}_{II})},\qquad
P_X^{\Pi,\mathrm{red}}\cong\fg_{\Delta_5}.
\]
Construction at finite Harder--Narasimhan height \(R\) proceeds through
eight construction steps.  Each step is a procedure on input data
already supplied by earlier steps and feeds named output data into the
next.  The eight steps stratify the realization datum
\(\mathfrak K_X=(L_X,H_X,O_X,D_X,R_X)\) of
Definition~\ref{thm:finite-stage-dirac-igusa-pro-object}: steps 1--3
realize \((L_X)\) and \((H_X)\); step 4 realizes the local content of
\((O_X)\); step 5 produces the Hall correspondence object; step 6 fixes
the Pfaffian--Dirac line and feeds the orientation character; steps 7--8
realize \((D_X)\) and \((R_X)\).

Throughout, \(X=S\times E\) with \(S\) a smooth complex K3 surface and
\(E\) a smooth elliptic curve.  \(\sigma=\sigma_{s,t}\) is Liu's product
stability condition on the Abramovich--Polishchuk heart in
\(\mathsf{D}^b(\mathrm{Coh}(X))\)
\cite{AbramovichPolishchukTStructures,LiuProductStability}.  The HN
sector \((\sigma,S)\) is fixed and the cofinal subsystem
\(\mathcal R_{\mathrm{HN}}=\{R_0<R_1<\cdots\}\) is the one of
Definition~\ref{def:O1-strong-reduced-orientation}, governed by
\textup{[K3$\times$E-fin]}, \textup{[K3$\times$E-atlas]}, and
\textup{[ML]}.

Each construction step carries a licensing tag pulled from the
five-type universal stage chain
\(\mathsf P\to\mathsf C\to\mathsf S\to\mathsf Z\to\mathsf A\) of
\texttt{chiral-bar-cobar-vol2:CLAUDE.md~§9}: \(\alpha\) primitive,
\(\beta\) functorial passage, \(\gamma\) shadow specialization,
\(\delta\) endpoint convergence, \(\epsilon\) terminal scalar.  The
ordering charge\(\to\)moduli\(\to\)hybrid\(\to\)orientation\(\to\)Hall
\(\to\)Dirac+Pfaffian\(\to\)Koszul\(\to\)recognition is the unique
ordering compatible with input/output type discipline; reordering any
two adjacent steps fails an input check.

The four obstructions are attributed to specific steps:
\((D_0)\) is the cofinal-trivialization datum of the entire eight-step
construction at every finite \(R\) (steps 1--6 supply its components);
\((O1)\) lives at step 4; \((O1)^+\) lives at step 4 with a Coxeter
coherence check feeding step 6; \((O2)\) lives at step 6.  Step 8
inherits the obstruction tower discharged by steps 1--7 and adds the
exact-completion residual \textup{(R8)} of
Definition~\ref{def:cofinal-primitive-recognition-datum}.

\subsection{Construction~1 — Charge window}
\label{app:build-charge-window}

\paragraph{Tag.} \(\alpha\)-primitive (lattice datum).

\paragraph{Input.}
HN height \(R\in\mathcal R_{\mathrm{HN}}\); the formal dyonic Mukai
lattice \(\Gamma_X^{\mathrm{form}}\) of
Definition~\ref{def:gram-index-lattice}; the Gram map
\(\Pi_X:\Gamma_X^{\mathrm{form}}\to\Gamma_{\mathrm{gram}}\) and the
cocycle \(B\) of Lemma~\ref{lem:gram-additivity-defect}; the active
support \(\Gamma_{\mathrm{act}}=\{\gamma\in\Gamma_{\mathrm{eff}}\mid
f(nm,l)\ne0\}\) of the Igusa Jacobi form \(\phi_{0,1}\); the cofinal
subsystem \(\mathcal R_{\mathrm{HN}}\).

\paragraph{Output.}
A finite cusp-positive active window
\[
A_R\subset\Gamma_{\mathrm{act}}\subset\Gamma_{\mathrm{gram}},
\qquad
A_R\hbox{ downward saturated for the cofinal target-window order},
\]
together with the normal-ordered extension
\[
\widehat\Gamma_R=
\{(c,T)\mid c\in\Gamma_R^{HN},\ T\in\mathcal T_R(c)\}
\subset\widehat\Gamma_X,
\qquad
\Gamma_R^\Pi=\overline\Pi_X(\widehat\Gamma_R)\subset\Gamma_{\mathrm{gram}},
\]
the lift \(L:\Gamma_{\mathrm{gram}}\to\Gamma_X^{\mathrm{form}}\) of
Lemma~\ref{lem:primitive-lift-every-gram-triple}, and the finite test
window \(\Gamma_R^{\mathrm{test}}\) of
Definition~\ref{def:O1-strong-reduced-orientation}.

\paragraph{Procedure.}
\begin{enumerate}
\item Compute the active support cone
\[
\Gamma_{\mathrm{act}}=\{(n,l,m)\in\Gamma_{\mathrm{eff}}\mid f(nm,l)\ne 0\}
\]
from the Fourier expansion of the index-one weak Jacobi form
\(\phi_{0,1}(\tau,z)=\sum_{n,l}c(n,l)q^nr^l\) of
\cite{EZ:1}, with the index substitution
\(f(nm,l):=c(nm,l)\) tabulated in
Appendix~\ref{sec:n1-boundary-compatibility-conditions} for the
\(N=1\) Igusa boundary.

\item Set
\[
A_R^{\circ}
=\{\gamma=(n,l,m)\in\Gamma_{\mathrm{act}}\mid n+|l|+m\le R\}.
\]
Form its downward saturation
\[
A_R
=\{\gamma'\in\Gamma_{\mathrm{act}}\mid
\exists\gamma\in A_R^{\circ},\ 0<\gamma'\le\gamma,\
\gamma-\gamma'\in Q_+\},
\]
in the sense of Definition~\ref{def:cofinal-primitive-recognition-datum}.

\item Normal-ordered lift.  For each \(c\in\Gamma_R^{HN}\) compute the
\(B\)-translate orbit
\[
\mathcal T_R(c)
=\Bigl\{\sum_{i<j}B(c_i,c_j)\Bigm|c=\sum_i c_i,\
c_i\in F_{\sigma,S}(R),\ \sum_i h_S(c_i)\le R\Bigr\}.
\]
Set \(\widehat\Gamma_R=\{(c,T):c\in\Gamma_R^{HN},\ T\in\mathcal T_R(c)\}\).

\item Finite target test window:
\(\Gamma_R^{\mathrm{test}}=\overline\Pi_X(\widehat\Gamma_R)\), compatibly
with \(R\to R'\) by inclusion.

\item Output \((A_R,\widehat\Gamma_R,\Gamma_R^\Pi,\Gamma_R^{\mathrm{test}},L)\).
\end{enumerate}

\paragraph{Type checks.}
\(A_R\) is finite and downward saturated;
\(\overline\Pi_X(\widehat\Gamma_R)=\Gamma_R^\Pi\) lies in
\(\Gamma_{\mathrm{gram}}\); \(\Gamma_R^{HN}\subset
F_{\sigma,S}(R)\) is finite by \textup{[K3$\times$E-fin]}; the cocycle
identity
\(\Pi_X(c+c')=\Pi_X(c)+\Pi_X(c')+B(c,c')\) is an axiom of the cocycle.

\paragraph{Obstructions discharged or invoked.}
None; this step is purely lattice combinatorics from the imported
Gritsenko--Nikulin Borcherds product
\cite{GN,GNPartI,GNII,BorcherdsGKM88,Borcherds95}.
The downward saturation is the input requirement of \textup{(R0)} in
Definition~\ref{def:cofinal-primitive-recognition-datum}.

\subsection{Construction~2 — Finite moduli}
\label{app:build-finite-moduli}

\paragraph{Tag.} \(\alpha\)-primitive (geometric source).

\paragraph{Input.}
The active window \(A_R\) and lift \(L\) of
Construction~\ref{app:build-charge-window}; Liu's product stability
condition \(\sigma_{s,t}\); the standing hypotheses
\textup{[K3$\times$E-fin]} and \textup{[K3$\times$E-atlas]}.

\paragraph{Output.}
A finite class set \(C_R\) and lift
\(\ell_R:\Gamma_R^\Pi\to C_R\) with
\(\Pi_X(\ell_R(\gamma))=\gamma\); the family of finite-type quasi-smooth
derived Artin substacks
\[
\mathfrak M^{ss}_{R,c}
=\mathfrak M^{ss}_{\sigma_{s,t}}(X,c)_{\le R}
\subset\operatorname{Perf}(X),
\qquad c\in C_R,
\]
of \(\sigma_{s,t}\)-semistable objects of full Liu numerical class
\(c\) and HN height \(\le R\); the scalar-rigidified truncation
\(\mathfrak M^{\mathrm{sc}}_{R,c}\); the universal perfect complex
\(\mathcal A_{R,c}\) on \(X\times\mathfrak M^{\mathrm{sc}}_{R,c}\); the
finite stratification by retained closed substacks of finite residual
inertia.

\paragraph{Procedure.}
\begin{enumerate}
\item Form \(C_R=\{c=L\gamma:\gamma\in A_R\}\) closed under retained
extension sums (taken in \(\Gamma_X^{\mathrm{form}}\)).

\item For each \(c\in C_R\) cut the substack of
\(\sigma_{s,t}\)-semistable objects of class \(c\) by HN height \(\le
R\):
\[
\mathfrak M^{ss}_{R,c}=
\bigl\{[A]\in\operatorname{Perf}(X)\mid
\sigma_{s,t}\hbox{-semistable},\
\hn(A)\le R,\
\nu(A)=c\bigr\}.
\]

\item Verify finite-typeness: \textup{[K3$\times$E-fin]} (Liu product
stability + product-boundedness theorem at fixed full Liu charge with
finite fibres over the active charge) gives the finite-type assertion.
PTVV \cite{PTVV2013} on the Calabi--Yau threefold \(X\) gives the
\((-1)\)-shifted symplectic structure on
\(\operatorname{Perf}(X)\), restricted to the retained semistable
sector.

\item Quasi-smoothness: the universal perfect complex
\(\mathcal A_{R,c}\) supplies a relative cotangent complex of bounded
Tor-amplitude on the retained charts.  Tangent at \(A\) is
\(T_A\mathfrak M^{ss}_{R,c}\simeq\operatorname{RHom}_X(A,A)[1]\).

\item Scalar rigidification: form \(\mathfrak M^{\mathrm{sc}}_{R,c}\) by
killing the scalar \(\mathbb G_m\)-automorphism of
\(c\)-stable factors.

\item Two-step flag stack
\(\mathfrak F^{(2)}_{R,c,c',c''}\) and extension stack
\(\mathfrak E_{R,c,c'}\) constructed as iterated derived vector stacks
of \(R\pi_*\mathcal RHom_X(\mathcal B,\mathcal A[1])\); both finite type
and quasi-smooth in the retained sector by
Proposition~\ref{prop:finite-dcritical-hall-product}.

\item Output \((C_R,\ell_R,\{\mathfrak M^{ss}_{R,c}\},
\{\mathfrak E_{R,c,c'}\},\{\mathfrak F^{(2)}_{R,c,c',c''}\})\) with
universal perfect complexes and the finite stratification.
\end{enumerate}

\paragraph{Type checks.}
\(\mathfrak M^{ss}_{R,c}\) is a finite-type quasi-smooth derived Artin
substack of \(\operatorname{Perf}(X)\); the scalar rigidification leaves
only finite residual inertia
\(H_S\subset E_{\mathrm{tors}}\) on each connected stabilizer
component, by
Proposition~\ref{prop:finite-dcritical-hall-product};
PTVV \((-1)\)-shifted symplectic structure restricts.

\paragraph{Obstructions discharged or invoked.}
\textup{[K3$\times$E-fin]} is invoked here in its product-boundedness
form and supplies the finite-type assertion not in the cited literature
for the K3\(\times\)E threefold; it is part of the standing datum.
\textup{[K3$\times$E-atlas]} supplies the quasi-smooth d-critical
charts.  These enter \((D_0)\) at the level of the finite Hall atlas
of Theorem~\ref{thm:finite-stage-dirac-igusa-pro-object}.

\subsection{Construction~3 — Hybrid wrapped factorization carrier}
\label{app:build-hybrid-wrapped}

\paragraph{Tag.} \(\beta\)-functorial passage (Stage-1 carrier).

\paragraph{Input.}
The finite moduli stacks
\(\{\mathfrak M^{\mathrm{sc}}_{R,c}\}\) of
Construction~\ref{app:build-finite-moduli}; the
projection-finite/wrapped split by elliptic degree
\(b_R^{\mathrm{geom}}:\Gamma_{\sigma,S}^{HN}(R)\to\Z_{\ge0}\) of
Definition~\ref{def:hybrid-wrapped-factorization-datum}; a fixed
origin \(0_E\in E\); the Abramovich--Polishchuk/Liu wrapped anchor
candidate.

\paragraph{Output.}
The finite hybrid local/wrapped factorization datum
\[
\mathcal F_{X,\sigma,S,\le R}^{\mathrm{hyb}}
=
(\mathcal C_{\sigma,S,\le R},
b_R^{\mathrm{geom}},
\mathcal F_R^{\mathrm{loc}},
\mathcal G_R^{\mathrm{wr}},
\mathfrak E_R^{\mathrm{loc}},
\mathfrak E_R^{\mathrm{hyb}},
\mathfrak E_R^{\mathrm{wr}},
\mathfrak{Flag}_R^{\mathrm{hyb}},
\mathsf{Corr}_R^{E,\mathrm{hyb}},
\lambda_R,
Q_{E,R},
\theta^Q_R,
o_R,
I_R^{\mathrm{prot}})
\]
of Definition~\ref{def:hybrid-wrapped-factorization-datum},
together with the hybrid Ran base
\(\operatorname{Ran}^{\mathrm{hyb}}(E)\), the eight-word two-step
binary flag atlas \(\{w\in\{\mathrm L,\mathrm W\}^3\}\), and the
quotient-after-correspondence functor \(Q_{E,R}\).

\paragraph{Procedure.}
\begin{enumerate}
\item \emph{Elliptic-degree split.}  Decompose
\(\Gamma_{\sigma,S}^{HN}(R)
=\Gamma_R^{\mathrm{loc}}\sqcup\Gamma_R^{\mathrm{wr}}\) where
\(\Gamma_R^{\mathrm{loc}}=\{b_R^{\mathrm{geom}}=0\}\) and
\(\Gamma_R^{\mathrm{wr}}=\{b_R^{\mathrm{geom}}>0\}\).  Verify additivity
\(b_R^{\mathrm{geom}}(\gamma+\gamma')
=b_R^{\mathrm{geom}}(\gamma)+b_R^{\mathrm{geom}}(\gamma')\) on retained
extensions remaining in the quotient.

\item \emph{Projection-finite local sector.}  For
\(\alpha\in\Gamma_R^{\mathrm{loc}}\) and finite index set \(I\), build
the closed-support incidence stack over \(E^I\):
\[
\mathfrak M^{\mathrm{loc,cl}}_{\alpha,R,I}
=
\Bigl\{(A,\mathbf e)\in\mathfrak M^{\mathrm{sc}}_{R,L\alpha}\times E^I
\Bigm|p_E(\mathrm{Supp}\,A)\subset\bigcup_{i\in I}S\times\{e_i\}\Bigr\}.
\]
Form
\(\mathscr H^{\mathrm{loc}}_{\alpha,R,I}=
(\pi_{\alpha,R,I})_!(\Phi_{\alpha,R,I}^{\mathrm{loc}}\otimes
o_{\alpha,R,I}^{\mathrm{loc}})\) and the open-support sheaf
\(\mathcal F_{\alpha,R}^{\mathrm{loc}}(U)\) over \(U\subset E\) by
restricted Borel--Moore configuration sums.

\item \emph{Wrapped prequotient sector.}  For
\(\eta\in\Gamma_R^{\mathrm{wr}}\), form the rigidified prequotient
\(\mathcal M_{\eta,R}^{\mathrm{wr,rig}}\to\mathcal M_{\eta,R}^{\mathrm{wr}}\)
by the Abramovich--Polishchuk / determinant anchor
\[
\lambda_\eta:\mathcal M_{\eta,R}^{\mathrm{wr,rig}}\longrightarrow E,
\qquad
\lambda_\eta(A)=\det Rp_{E,*}A
\otimes\mathcal O_E(-\chi(A)\cdot 0_E),
\]
divided by a finite cover when \(\chi(A)\ne 1\) so that
\(\lambda_\eta(t\cdot A)=\lambda_\eta(A)+t\) holds.  Verify the
losslessness residual \(\mathfrak o^{\lambda,\mathrm{loss}}_R=0\) by
the Ext-distinguishing test of
Definition~\ref{def:hybrid-wrapped-factorization-datum}, separating
\(W_0=i_{s,*}(\mathcal O_E\oplus\mathcal O_E)\) from
\(W_L=i_{s,*}(L\oplus L^{-1})\).

\item \emph{Mixed local/wrapped correspondences.}  For an ordered
hybrid input
\(\alpha_-\in\Gamma_R^{\mathrm{loc}}\),
\(\eta\in\Gamma_R^{\mathrm{wr}}\),
\(\beta_+\in\Gamma_R^{\mathrm{loc}}\), form
\(\mathfrak E^{\mathrm{mix}}_{\alpha_-;\eta;\beta_+,R}\) before the
reduced \(E\)-quotient.

\item \emph{Eight-word two-step flag atlas.}  Build flag stacks
\(\mathfrak F^{(2),w}_R\) for each
\(w\in\{\mathrm L,\mathrm W\}^3=
\{\mathrm{LLL},\mathrm{LLW},\mathrm{LWL},\mathrm{WLL},\mathrm{LWW},
\mathrm{WLW},\mathrm{WWL},\mathrm{WWW}\}\), classifying filtrations
with three associated graded factors of those types.  Verify the
four-input pentagon coherence and the quotient-after-correspondence
descent.

\item \emph{Hybrid correspondence-module functor.}  Assemble the
correspondence category
\(\mathsf{Corr}^{E,\mathrm{hyb}}_R\) and the functor
\(\mathcal B_R^{E,\mathrm{hyb}}:\mathsf{Corr}^{E,\mathrm{hyb}}_R
\to\mathsf{Ch}_{\C}\) by
\(\mathcal B_R^{E,\mathrm{hyb}}(e)=q_!\,\mathrm{TS}^{\mathrm{red}}_e\,p^*\)
for every generating arrow
\(e=(Y\xleftarrow{p}\mathfrak E_e\xrightarrow{q}Z)\).

\item \emph{Quotient functor \(Q_{E,R}\).}  Apply
\(E\)-equivariant compactly supported vanishing-cycle chains and the
orientation descent of step~4 to produce the reduced object
\[
\mathcal F^{\mathrm{geom}}_{X,\sigma,S,\le R}
=Q_{E,R}\circ\mathcal B_R^{E,\mathrm{hyb}}.
\]

\item Output the hybrid datum with all coherences and the protected
integration map \(I_R^{\mathrm{prot}}\), with homogeneous
\(s\)-degree the Borcherds trace exponent
\(m_R=\operatorname{pr}_3\overline\Pi_X\).
\end{enumerate}

\paragraph{Type checks.}
The carrier
\(\mathcal F_{X,\sigma,S,\le R}^{\mathrm{hyb}}\) is a hybrid
correspondence-module functor on
\(\mathsf{Corr}^{E,\mathrm{hyb}}_R\) with values in
\(\mathsf{Ch}_{\C}\); the four-input pentagon is supplied by
the eight-word two-step flag atlas; the anchor residual tuple
\(\mathfrak o^\lambda_R=
(\mathfrak o^{\lambda,\mathrm{ex}}_R,
\mathfrak o^{\lambda,\mathrm{unit}}_R,
\mathfrak o^{\lambda,\mathrm{loss}}_R,
\mathfrak o^{\lambda,\mathrm{multi}}_R,
\mathfrak o^{\lambda,\mathrm{tr}}_{R'R})\) vanishes; the hybrid
residual \(\mathfrak O_{\mathrm{hyb},R}=0\) at every height \(R\),
compatibly with transition maps.

\paragraph{Obstructions discharged or invoked.}
The middle type-II wall lies in
\(\delta_2\leftrightarrow(0,1,1)\), elliptic degree
\(d=m+1=2\) by Lemma~\ref{app:pfaffian-wall-three-walls}; it
cannot be represented by a projection-finite local stack and demands
the wrapped prequotient.  This step's correctness — the existence of
the wrapped representative
\(x_{\delta_2,R}\in\mathcal M_{\eta,R}^{\mathrm{wr,rig}}\) with
\(\overline\Pi_X(x_{\delta_2,R})=(0,1,1)\) and its compatible anchor
\(\lambda_{\eta,R}\) — is precisely the hybrid input to \((O2)\) at
the middle wall.

\subsection{Construction~4 — Reduced d-critical orientation
and Joyce--Upmeier transport}
\label{app:build-orientation}

\paragraph{Tag.} \(\beta\)-functorial passage (oriented carrier).

\paragraph{Input.}
The hybrid carrier
\(\mathcal F_{X,\sigma,S,\le R}^{\mathrm{hyb}}\) of
Construction~\ref{app:build-hybrid-wrapped}; the K3 holomorphic
\(2\)-form \(\sigma_S\) on \(S\); BBDJS d-critical chart machinery
\cite{BBDJS2015}; the Joyce--Upmeier multiplicative orientation
formalism \cite{JoyceUpmeier,JoyceTanakaUpmeier}; the cosection
localization machinery of Kiem--Li and the cosection-reduced
obstruction theory; the type-II Weyl group
\(W^{(2)}(\La^{2,1}_{II})\) generated by
\(s_{\delta_1},s_{\delta_2},s_{\delta_3}\).

\paragraph{Output.}
The cosection
\(\sigma=p_S^*\sigma_S\); the reduced tangent complex
\[
\operatorname{RHom}_{\mathrm{red}}(A,A)
=
\operatorname{Cone}\!\Bigl(
\operatorname{RHom}_X(A,A)
\xrightarrow{(\mathrm{tr},\mathrm{cs})}
R\Gamma(X,\mathcal O_X)\oplus H^2(S,\mathcal O_S)[-1]
\Bigr)[-1];
\]
the BBDJS vanishing-cycle complex \(\Phi_{R,c}\); the reduced
orientation line
\[
o_{R,c}
=\det\operatorname{Ext}^1_{\mathrm{red}}(A,A)^{-1}
\quad\text{on Darboux charts},\qquad
o_{R,c}^{\otimes 2}\simeq\det\operatorname{RHom}_{\mathrm{red}}(A,A);
\]
the Joyce--Upmeier multiplicative transport
\(o_{c\oplus c'}\simeq o_c\otimes o_{c'}\) compatible with the reduced
extension correspondences; the quotient orientation
\(\mathfrak Q^{or}_R\) of
Theorem~\ref{prop:generic-local-sign-orientation-character}; the
Weyl-equivariant lifts
\(\tau_{i,R,S}:o_{R,S}\to s_{\delta_i}^*o_{R,S}\) for \(i=1,2,3\) of
\((O1)^+\); the local generic wall-sign data \(\epsilon_{o,R}^{\mathrm{loc}}\).

\paragraph{Procedure.}
\begin{enumerate}
\item \emph{Cosection.}  Pull back the K3 semiregularity cosection
\(\mathrm{cs}\) along \(p_S:X\to S\) to obtain the 2-step cosection
\(\sigma=p_S^*\sigma_S\) on the K3 factor; verify additivity in exact
triangles by \textup{(RedOr)} of
Theorem~\ref{thm:finite-stage-dirac-igusa-pro-object}.

\item \emph{Reduced obstruction.}  Form the reduced tangent complex by
the displayed cone above (Kiem--Li cosection localization).  PTVV
\cite{PTVV2013} \((-1)\)-shifted symplectic on
\(\operatorname{RHom}_X(A,A)\) descends, after the cosection
contraction, to the reduced complex.

\item \emph{BBDJS d-critical chart.}  On each Darboux chart of
\(\mathfrak M^{\mathrm{sc}}_{R,c}\), build the BBDJS vanishing-cycle
complex \(\Phi_{R,c}\) \cite{BBDJS2015}; verify reduced
Thom--Sebastiani transport coherent on two-step flag stacks of
Construction~\ref{app:build-hybrid-wrapped}.

\item \emph{Reduced orientation line.}  Define
\(o_{R,c}=\det\operatorname{Ext}^1_{\mathrm{red}}(A,A)^{-1}\); verify
the squared isomorphism
\(o_{R,c}^{\otimes 2}\simeq\det\operatorname{RHom}_{\mathrm{red}}(A,A)\)
by \(3\)-Calabi--Yau Serre duality
\(\operatorname{Ext}^2_{\mathrm{red}}(A,A)
\simeq\operatorname{Ext}^1_{\mathrm{red}}(A,A)^\vee\).

\item \emph{Quotient-orientation descent on the K3\(\times\)E-quotient
self-Ext frame bundle.}  Compute the four Cech--Borel obstruction
classes for every retained stratum \(S\):
\[
\alpha^{\mathrm{red}}_{R,S}=w_2(\mathscr D_{R,S}),
\quad
\alpha^{E,\mathrm{free}}_{R,S}=
\mathrm{edge}_{2,0}(w_2^E(\mathscr D_{R,S})),
\quad
\widetilde\beta^H_{R,S}=w_2^G(\mathscr D_{R,S}),
\quad
\lambda^H_{R,S}\in H^1_G(S;\mathbb F_2),
\]
and verify simultaneous vanishing with compatible
null-trivializations on overlaps, on Hall correspondence pullbacks,
on two-step flag stacks, and under successor maps.  The full
calculation proceeds primary stabilizer by stabilizer; for cyclic
\(\Z/2^a\) the absent mixed coefficient requires the rank-two
two-primary check.  Output the descended quotient orientation system
\(\mathfrak Q^{or}_R\).

\item \emph{Multiplicative Joyce--Upmeier transport.}  Transport
\(o_{c\oplus c'}\simeq o_c\otimes o_{c'}\) coherently across two-step
flag stacks
\cite{JoyceUpmeier,JoyceTanakaUpmeier}; verify the reduced
Thom--Sebastiani identification on three-input flags.

\item \emph{Weyl-equivariant lifts \((O1)^+\).}  For each
\(i\in\{1,2,3\}\) construct the lift
\(\tau_{i,R,S}:o_{R,S}\to s_{\delta_i}^*o_{R,S}\); verify the
projective Weyl cocycle
\(c_{o,R}\in Z^2(\mathcal G_R;\underline{\mathbb F}_2)\) on the
retained action groupoid; null-trivialize \([c_{o,R}]\) and choose a
\(1\)-cochain killing it; verify \(\tau_{i,R,S}^2=1\); verify the
Coxeter relations
\((s_{\delta_i}s_{\delta_j})^3=1\) for \(i\ne j\) (the
\((3,3,3)\)-hyperbolic triangle relations of
Appendix~\ref{app:pfaffian-wall-three-walls}); verify that the torsor
defects \(\omega_{i,\mathcal C}\) of
Definition~\ref{thm:finite-stage-dirac-igusa-pro-object}~$(O_X)$ vanish.

\item \emph{Local wall-sign datum.}  At each generic type-II wall point
\(\Hpl_{\delta_i}\), with the rank-one local model
\(K_{\delta_i}^{\mathrm{cross}}=[\R\xrightarrow{u_i}\R]\), record the
local generic sign \(\epsilon_{o,R}^{\mathrm{loc}}(s_{\delta_i})=-1\)
of Proposition~\ref{prop:rank-one-crossing-o2-sign}.

\item Output
\((o_R,\mathfrak Q^{or}_R,\{\tau_{i,R,S}\},\epsilon_{o,R}^{\mathrm{loc}})\).
\end{enumerate}

\paragraph{Type checks.}
\(o_R\) is a Picard-line section squaring to the reduced determinant;
the four Cech--Borel classes vanish with compatible
null-trivializations; the lifts \(\tau_{i,R}\) satisfy
\(\tau_{i,R}^2=1\) and the \((3,3,3)\)-Coxeter relations; multiplicative
transport agrees on flag stacks.

\paragraph{Obstructions discharged or invoked.}
This step is the entire site of \((O1)\) (strong reduced orientation
on the K3\(\times\)E-quotient self-Ext frame bundle) and of
\((O1)^+\) (Weyl-equivariant transport along
\(W^{(2)}(\La^{2,1}_{II})\) reflections).  Discharging \((O1)\)
requires the simultaneous vanishing and compatible
null-trivialization of the four Cech--Borel classes on every retained
stratum; discharging \((O1)^+\) requires \([c_{o,R}]=0\), the
Coxeter coherence, and the vanishing of all torsor defects
\(\omega_{i,\mathcal C}\).  The local wall-sign datum is the input to
\((O2)\) at step~6.

\subsection{Construction~5 — Hall correspondences and primitive part}
\label{app:build-hall}

\paragraph{Tag.} \(\gamma\)-shadow specialization
(Stage-2 chiral Hall shadow).

\paragraph{Input.}
The oriented hybrid carrier from
Construction~\ref{app:build-orientation}; the extension stack
\(\mathfrak E^{\mathrm{geom}}_{R,\widehat c,\widehat c'}\); the
two-step flag stack
\(\mathfrak F^{(2),\mathrm{geom}}_{R,\widehat c,\widehat c',\widehat c''}\);
the BBDJS reduced Thom--Sebastiani transport; six-functor admissibility
of \(p^*,q_!,p_!\) (Construction~\ref{app:build-finite-moduli}~step
\textup{(Six)}).

\paragraph{Output.}
The finite Hall object
\[
\mathcal H_R^{\mathrm{geom}}
=\bigoplus_{\widehat c\in\widehat\Gamma_R}
H_c^\bullet\!\bigl(
\mathfrak M^{\mathrm{geom}}_{R,\widehat c},
\Phi_{R,c}\otimes o_{R,c}\bigr);
\]
the Hall product \(m_R=q_!\,\mathrm{TS}^{\mathrm{red}}\,p^*\) and Hall
coproduct \(\Delta_R=p_!\,(\mathrm{TS}^{\mathrm{red}})^{-1}\,q^*\); the
finite Hall bialgebra
\((\mathcal H_R^{\mathrm{geom}},m_R,\Delta_R,\eta_R,\epsilon_R)\); the
primitive part
\[
P_R^{\mathrm{geom}}
=\ker(\Delta_R-1\otimes\mathrm{id}-\mathrm{id}\otimes 1)\cap\ker\epsilon_R;
\]
the normal-ordered Gram descent
\(P_R^{\Pi,\mathrm{geom}}
=\overline\Pi_{X,*}^{\Theta_R}P_R^{\mathrm{geom}}\); the
\(\Gamma_R^\Pi\)-graded super-Lie bracket
\([\,\bar x,\bar y\,]_{\mathrm{red}}\) of
Proposition~\ref{thm:hall-primitives-formal};
the Hopf pairing matrix \(G_{\gamma,ab}\) and radical \(K_\gamma\).

\paragraph{Procedure.}
\begin{enumerate}
\item \emph{Pullback.}  For
\(\widehat c,\widehat c'\in\widehat\Gamma_R\), pull back along
\(p:\mathfrak E^{\mathrm{geom}}_{R,\widehat c,\widehat c'}
\to\mathfrak M^{\mathrm{geom}}_{R,\widehat c}
\times\mathfrak M^{\mathrm{geom}}_{R,\widehat c'}\); apply reduced
Thom--Sebastiani \(\mathrm{TS}^{\mathrm{red}}\) to combine vanishing
cycle data \cite{BBDJS2015}; pushforward exceptionally along
\(q:\mathfrak E^{\mathrm{geom}}_{R,\widehat c,\widehat c'}
\to\mathfrak M^{\mathrm{geom}}_{R,\widehat c\star\widehat c'}\) to
obtain
\(m_R(x\otimes y)=q_!\,\mathrm{TS}^{\mathrm{red}}\,p^*(x\boxtimes y)\).

\item \emph{Coproduct.}  Dually pull back along \(q\), apply
inverse reduced Thom--Sebastiani, pushforward along \(p\) to obtain
\(\Delta_R(z)=p_!\,(\mathrm{TS}^{\mathrm{red}})^{-1}\,q^*(z)\).
The compact geometric coproduct correspondence
\(\mathfrak C^{\mathrm{geom}}_{\rho;\mu,\nu}\) of
Proposition~\ref{thm:hall-primitives-formal} is the
support data for splitting.

\item \emph{Bialgebra check.}  Verify associativity of \(m_R\) and
coassociativity of \(\Delta_R\) on the two-step flag stack
\(\mathfrak F^{(2),\mathrm{geom}}_{R,\widehat c,\widehat c',\widehat c''}\)
by base change and Thom--Sebastiani coherence.  Verify the bialgebra
identity (Hopf compatibility) via the Massey--Saito Joyce--Upmeier
multiplicative orientation transport.

\item \emph{Primitive extraction.}  Compute
\(P_R^{\mathrm{geom}}=\ker(\Delta_R-1\otimes\mathrm{id}
-\mathrm{id}\otimes 1)\cap\ker\epsilon_R\); since
\(\Delta_R-1\otimes\mathrm{id}-\mathrm{id}\otimes 1\) is degree-graded
on \(\widehat\Gamma_R\), this reduces to a finite linear computation
in each charge.

\item \emph{Normal-ordered descent.}  Apply
\(\overline\Pi_{X,*}^{\Theta_R}\) of
Theorem~\ref{thm:hall-product} to obtain the
\(\Gamma_R^\Pi\)-graded
\(P_R^{\Pi,\mathrm{geom}}\); verify additivity
\(\overline\Pi_X(\widehat c\star\widehat c')
=\overline\Pi_X(\widehat c)+\overline\Pi_X(\widehat c')\).

\item \emph{Bracket and pairing matrices.}  Compute the matrices
\[
M^{\alpha+\beta,c}_{\alpha,a;\beta,b}
=\int_{\mathfrak E^{\mathrm{geom}}_{\alpha,\beta;\alpha+\beta}}
q^!(e_{\alpha+\beta,c}^\vee)\cap
\mathrm{TS}^{\mathrm{red}}(p^*(e_{\alpha,a}\boxtimes e_{\beta,b})),
\quad
B_{\alpha,\beta}(x\otimes y)
=M_{\alpha,\beta}(x\otimes y)
-(-1)^{|x||y|}M_{\beta,\alpha}\mathsf{sw}(x\otimes y),
\]
\[
G_{\gamma,ab}=
\mathrm{tr}_0(q_!\,\mathrm{TS}^{\mathrm{red}}\,p^*
(e_{\gamma,a}\boxtimes e_{-\gamma,b}))
\quad\hbox{on}\quad
\mathfrak P^{\mathrm{geom}}_{\gamma,-\gamma;0}.
\]
Form the radical
\(K_{\gamma,\bar p}
=\ker G_{\gamma,\bar p}\cap\ker G_{-\gamma,\bar p}^t\) and the
quotient \(L_{\gamma,\bar p}\); fix splittings.

\item Output
\((\mathcal H_R^{\mathrm{geom}},m_R,\Delta_R,P_R^{\mathrm{geom}},
P_R^{\Pi,\mathrm{geom}},
\{B_{\alpha,\beta}\},\{G_{\gamma,ab}\},\{K_\gamma\},\{L_\gamma\})\).
\end{enumerate}

\paragraph{Type checks.}
\(\mathcal H_R^{\mathrm{geom}}\) is a finite \(\widehat\Gamma_R\)-graded
Hall bialgebra; \(m_R\) is associative on the two-step flag stack;
\(\Delta_R\) is coassociative; the supercommutator of primitives is
primitive (the bialgebra identity); the descended bracket is
\(\Gamma_R^\Pi\)-graded after \(\overline\Pi_{X,*}^{\Theta_R}\);
finite-dimensional in each homogeneous parity piece by
finite-typeness.

\paragraph{Obstructions discharged or invoked.}
\((D_0)\) is partially discharged: the finite Hall package
\((\mathcal H_R^{\mathrm{geom}},m_R,\Delta_R,P_R)\) is constructed.
The seven obstruction classes
\(\mathfrak o_\Theta^{\mathrm{top}},
\mathfrak o_\Theta^{\mathrm{Hoch}},
\mathfrak o_\Theta^{\mathrm{tri}},
\mathfrak o_\Theta^{\mathrm{Jac}},
\mathfrak o_F,
\mathfrak o_\Theta^{\mathrm{cyc}},
\mathfrak o_\Theta^{\mathrm{rad}}\) of
Theorem~\ref{thm:hall-product} are required to vanish on
this finite Hall stage and compatibly in the inverse limit.

\paragraph{Cross-volume anchor.}  The output
\((\mathcal H_R^{\mathrm{geom}},m_R)\) is the finite-stage realization
of the K3-side BKM Hall--Drinfeld double
\(\mathcal D_\hbar(\mathrm{CoHA}_{K3\times E})\) of
\texttt{calabi-yau-quantum-groups:CLAUDE.md~§6.7}.  In the universal
stage chain \(\mathsf P\to\mathsf C\to\mathsf S\) the Hall positive half
\(Y^+(X)\) is precisely \(P_R^{\mathrm{geom}}\) at finite \(R\); the
Drinfeld double \(G(X)=D(Y^+(X))\) requires Hopf pairing, completion,
integral form, and stable-envelope transport which are supplied by
\(G_{\gamma,ab}\) and \((\mathrm{ML})\).  The κ-tuple coordinate
\(\kappa_{\mathrm{BKM}}=5\) is the weight of \(\Delta_5\) and the
output of the level-4 scalar trace; it is not a level-2 invariant of
the carrier itself.

\subsection{Construction~6 — Dirac operator and Pfaffian line}
\label{app:build-dirac-pfaffian}

\paragraph{Tag.} \(\gamma\)-shadow specialization
(orientation gerbe + Pfaffian).

\paragraph{Input.}
The oriented Hall package
\((\mathcal H_R^{\mathrm{geom}},o_R,\{\tau_{i,R,S}\})\) of
Constructions~\ref{app:build-orientation}--\ref{app:build-hall}; the
local wall-sign datum
\(\epsilon_{o,R}^{\mathrm{loc}}\) of
Construction~\ref{app:build-orientation}; the rank-one type-II
Pfaffian crossing datum of
Definition~\ref{def:rank-one-type-ii-crossing} and
Appendix~\ref{app:pfaffian-wall-rank-one}; the half-Hilbert orbit
under \(W^{(2)}(\La^{2,1}_{II})\) of size \(2^{12}=4096\).

\paragraph{Output.}
The Dirac operator on the orientation gerbe
\(D_R\); the Pfaffian line
\(\mathscr L_{\mathrm{Pf},R}=
\det\operatorname{Ext}^1_{\mathrm{red}}(A,A)^{-1}\) (squaring to the
reduced determinant); the protected Pfaffian section
\(\operatorname{pf}_{X,R}\); the global wall-sign character
\(\epsilon_{o,R}:W^{(2)}(\La^{2,1}_{II})\to\{\pm1\}\); the Pfaffian
section asymptotic
\[
\operatorname{pf}_{X,R}
=
64q^{1/2}r^{1/2}s^{1/2}
\prod_{\gamma=(n,l,m)\in A_R}
(1-q^nr^ls^m)^{f(nm,l)}+\hbox{higher-window correction},
\]
which exhausts to \(\Delta_5\) as \(R\to\infty\).

\paragraph{Procedure.}
\begin{enumerate}
\item \emph{Dirac operator.}  On each Darboux chart of
\(\mathfrak M^{\mathrm{sc}}_{R,c}\), with reduced normal self-Ext
splitting \(K^{\mathrm{nor}}_{R,c}\), build the odd Dirac operator
\(D_{R,c}\) acting on \(\Omega^{\mathrm{sp}}\otimes o_{R,c}\)
(spinors on the orientation gerbe); verify that, at a generic
type-II wall point, the local rank-one normal block is
\(D_{\delta_i}=\begin{psmallmatrix}0&u\\-u&0\end{psmallmatrix}\) with
\(\operatorname{Pf}(D_{\delta_i})=u\)
(Definition~\ref{def:rank-one-type-ii-crossing}).

\item \emph{Pfaffian line.}  Set
\(\mathscr L_{\mathrm{Pf},R}=
\det\operatorname{Ext}^1_{\mathrm{red}}(A,A)^{-1}\); verify
\(\mathscr L_{\mathrm{Pf},R}^{\otimes 2}\simeq
\det\operatorname{RHom}_{\mathrm{red}}(A,A)\) by
\(3\)-Calabi--Yau Serre duality.

\item \emph{Local wall sign.}  At each generic point of
\(\Hpl_{\delta_i}\) with normal model
\(K^{\mathrm{nor}}_{\delta_i,R}=[\R\xrightarrow{u_i}\R]\), compute
the monodromy of the Pfaffian section
\(\operatorname{pf}_{\delta_i,R}=\upsilon_iu_i\) along the
wall-reflection loop \(\ell_{\delta_i,R,S}\) closed by the
\((O1)^+\) lift \(\tau_{i,R,S}\):
\[
M_{\ell_{\delta_i,R,S}}(\upsilon_iu_i)
=\upsilon_i(-u_i)=-\upsilon_iu_i,
\qquad
\epsilon_{o,R}^{\mathrm{loc}}(s_{\delta_i})
=(-1)^{N^{\mathrm{Pf}}_{\delta_i,R}}=-1.
\]

\item \emph{Half-Hilbert orbit transport.}  The reflection
\(s_{\delta_i}\) acts on the half-Hilbert orbit (compact source labels
of size \(2^{12}=4096\) parametrized by 12 binary parities tied to the
\(2^{12}\) odd-spin structures on the lattice
\(\La^{2,1}_{II}/2\La^{2,1}_{II}\)).  Sum the local rank-one
contributions over the orbit:
\[
\epsilon_{o,R}(s_{\delta_i})
=\prod_{x\in\hbox{orbit}/2^{12}}
\epsilon_{o,R}^{\mathrm{loc}}(s_{\delta_i};x).
\]
The output is the global character
\(\epsilon_{o,R}:W^{(2)}(\La^{2,1}_{II})\to\{\pm 1\}\).

\item \emph{Compatibility package.}  Verify the component, boundary,
overlap, quotient, and protected integration compatibility package of
Proposition~\ref{prop:appE-local-pfaffian-sign} on every
retained type-II wall stratum.

\item \emph{Pfaffian section.}  Define
\(\operatorname{pf}_{X,R}\) as the protected integration of the
oriented Hall product over
\(\mathcal F_{X,\sigma,S,\le R}^{\mathrm{hyb}}\):
\[
\operatorname{pf}_{X,R}
=
64q^{1/2}r^{1/2}s^{1/2}
I_R^{\mathrm{prot}}\!\left(
\bigotimes_{\gamma\in A_R}\operatorname{pf}_{\gamma,R}\right),
\qquad
\operatorname{pf}_{\gamma,R}=(1-q^nr^ls^m)^{f(nm,l)}.
\]

\item Output
\((D_R,\mathscr L_{\mathrm{Pf},R},\operatorname{pf}_{X,R},
\epsilon_{o,R})\), with the protected integration compatibility data.
\end{enumerate}

\paragraph{Type checks.}
\(\mathscr L_{\mathrm{Pf},R}\) is a Picard-line whose square
isomorphism is supplied by reduced Serre duality;
\(\operatorname{pf}_{X,R}\) is a section of
\(\mathscr L_{\mathrm{Pf},R}\) with the displayed cusp expansion; the
character \(\epsilon_{o,R}\) takes values in \(\{\pm 1\}\) and
satisfies the Coxeter relations of
\(W^{(2)}(\La^{2,1}_{II})\); on each simple type-II reflection
\(\epsilon_{o,R}(s_{\delta_i})=-1\) by the rank-one local model.

\paragraph{Obstructions discharged or invoked.}
This step is the entire site of \((O2)\) (local Pfaffian wall normal
form on type-II walls; the \(2^{12}=4096\) sign sum on the
half-Hilbert orbit transport).  The local rank-one Pfaffian crossing
of Appendix~\ref{app:pfaffian-wall-rank-one} discharges \((O2)\)
generically; the global half-Hilbert orbit transport remains the open
\((O2)\) component of
Open~Problem~\ref{def:O1-plus-weyl-transport}.  The character
identity
\(\epsilon_o=\nu_{\Delta_5}\big|_{W^{(2)}(\La^{2,1}_{II})}\) becomes a
theorem only after \((D_0)\), \((O1)\), \((O1)^+\), and \((O2)\) are
discharged compatibly in \(R\).

\subsection{Construction~7 — Chiral Koszul source coalgebra}
\label{app:build-koszul}

\paragraph{Tag.} \(\gamma\)-shadow specialization
(\emph{Bar = twisting; not bulk}; the bulk is
\(Z^{\mathrm{der}}_{\mathrm{ch}}(\mathcal F_X^{\mathrm{hyb}})\),
constructed elsewhere).

\paragraph{Input.}
The oriented Hall primitives
\(P_R^{\Pi,\mathrm{geom}}\) of
Construction~\ref{app:build-hall}; the hybrid carrier
\(\mathcal F_{X,\sigma,S,\le R}^{\mathrm{hyb}}\) of
Construction~\ref{app:build-hybrid-wrapped}; the augmentation
\(\varepsilon_{X,R}^{\mathrm{hyb}}:
\mathcal F_{X,\sigma,S,\le R}^{\mathrm{hyb}}\to k\mathbf 1\); the
hybrid residual \(\mathfrak O_{\mathrm{hyb},R}=0\); the bounded
length \(N_R\) (positivity of HN height on every non-vacuum colour);
the Beilinson--Drinfeld/Francis--Gaitsgory bar-cobar formalism
\cite{BD,FrancisGaitsgoryCKD}; the conilpotent completed domain on
the elliptic curve \(E\).

\paragraph{Output.}
The conilpotent chiral coalgebra on \(E\):
\[
C_{X,R}^{\mathrm{geom}}
=\bar B_E^{\mathrm{ch}}
(\mathcal F_{X,\sigma,S,\le R}^{\mathrm{hyb}},
\varepsilon_{X,R}^{\mathrm{hyb}})
=k\mathbf 1\oplus\bigoplus_{1\le p\le N_R}
\bar B_E^{\mathrm{ch},p}(\overline{\mathcal F}_{X,R}^{\mathrm{geom}});
\]
the bar-length filtration
\(F^{\mathrm{co}}_{R,p}\); the chiral collision coproduct
\(\Delta_R^{\mathrm{ch}}\); the bar comparison
\(b_{X,R}=\mathrm{id}_{\bar B_E^{\mathrm{ch}}(\cdot)}\); the source
cobar counit
\(\varepsilon^{\mathrm{src}}_{\mathrm{bar},R}:
\Omega_E^{\mathrm{ch}}C_{X,R}^{\mathrm{geom}}
\xrightarrow{\simeq}
\mathcal F_{X,\sigma,S,\le R}^{\mathrm{hyb}}\); after a separate
source-to-target quasi-isomorphism
\(\vartheta_R:\mathcal F_{X,\sigma,S,\le R}^{\mathrm{hyb}}
\xrightarrow{\simeq}\mathsf{Den}_{\Delta_5,E,\le R}\), the cobar
comparison
\(\Theta_{\mathrm{Kos},R}=\vartheta_R\circ\varepsilon^{\mathrm{src}}_{\mathrm{bar},R}\).

\paragraph{Procedure.}
\begin{enumerate}
\item \emph{Augmentation kernel.}  Set
\(\overline{\mathcal F}_{X,R}^{\mathrm{geom}}
=\ker(\varepsilon_{X,R}^{\mathrm{hyb}})\); verify the bounded-length
hypothesis (positivity of HN height on every non-vacuum surviving
colour, hence at most \(N_R\) factors in any non-zero word).

\item \emph{Bar coalgebra.}  Form the reduced chiral bar
\[
C_{X,R}^{\mathrm{geom}}
=k\mathbf 1\oplus\bigoplus_{1\le p\le N_R}
\bar B_E^{\mathrm{ch},p}
(\overline{\mathcal F}_{X,R}^{\mathrm{geom}}),
\]
length \(p\) summand generated by ordered \(p\)-fold inputs from the
augmentation ideal of total HN height \(\le R\) (height \(>R\) outputs
zero); coaugmentation by the inclusion of the vacuum summand; counit by
projection.

\item \emph{Bar-length filtration.}
\[
F^{\mathrm{co}}_{R,p}C_{X,R}^{\mathrm{geom}}
=k\mathbf 1\oplus\bigoplus_{1\le j\le p}
\bar B_E^{\mathrm{ch},j}
(\overline{\mathcal F}_{X,R}^{\mathrm{geom}}),\qquad
0\le p\le N_R.
\]

\item \emph{Collision coproduct.}  For
\(x=[x_1|\cdots|x_p]\), define the cut-map coproduct
\[
\Delta_R^{\mathrm{ch}}(x)
=\mathbf1\otimes x+x\otimes\mathbf1
+\sum_{i=1}^{p-1}[x_1|\cdots|x_i]\otimes[x_{i+1}|\cdots|x_p],
\]
with each cut realized by the corresponding finite collision/flag
pull-push map in
\(\mathsf{Corr}^{E,\mathrm{hyb}}_{R,\mathrm{geom}}\) of
Construction~\ref{app:build-hybrid-wrapped}.  Verify
coassociativity by the four-input pentagon and the eight-word flag
coherences; verify counit identities via vacuum collision flags.

\item \emph{Conilpotence.}  Verify
\((\overline\Delta_R^{\mathrm{bar}})^{N_R+1}=0\): each reduced cut
strictly lowers bar length on each branch.

\item \emph{Bar comparison.}  Take \(b_{X,R}\) to be the identity of
\(\bar B_E^{\mathrm{ch}}
(\mathcal F_{X,\sigma,S,\le R}^{\mathrm{hyb}},
\varepsilon_{X,R}^{\mathrm{hyb}})\), eliminating the coalgebra-side
finite Koszul residuals
\(\mathfrak o^C_R, \mathfrak o^{\mathrm{fil}}_R,
\mathfrak o^{\mathrm{conil}}_R, \mathfrak o^{\Delta,\mathrm{ch}}_R\)
of Definition~\ref{def:chiral-koszul-source-datum}.

\item \emph{Source cobar counit.}  Apply Francis--Gaitsgory chiral
Koszul duality
\cite[Theorem~6.2.2]{FrancisGaitsgoryCKD} in the conilpotent completed
domain to obtain
\(\varepsilon^{\mathrm{src}}_{\mathrm{bar},R}:
\Omega_E^{\mathrm{ch}}C_{X,R}^{\mathrm{geom}}
\xrightarrow{\simeq}
\mathcal F_{X,\sigma,S,\le R}^{\mathrm{hyb}}\); verify it is a
source-internal quasi-isomorphism of chiral algebras on \(E\), not a
homotopy inverse of the target counit.

\item \emph{Source-to-target identification.}  Construct
\(\vartheta_R\) from the geometric primitive representatives, Hall
products, coproducts, pairings, and HN transition compatibility;
verify the residual
\(\mathfrak o^\vartheta_R\) of
Corollary~\ref{cor:source-bar-remaining-cobar-obstruction} vanishes;
set \(\Theta_{\mathrm{Kos},R}=\vartheta_R\circ
\varepsilon^{\mathrm{src}}_{\mathrm{bar},R}\).

\item \emph{Inverse limit.}  Verify the Koszul Mittag--Leffler
condition \(R^1\varprojlim=0\) for the inverse systems of source
coalgebras, completed chiral bar/cobar constructions, target
truncations, and the cones of \(b_{X,R}\) and
\(\Theta_{\mathrm{Kos},R}\) \cite[\S3.5]{WeibelHA}.  Form
\[
C_X^{\mathrm{geom}}
=\varprojlim_R C_{X,R}^{\mathrm{geom}},\qquad
\Theta_{\mathrm{Kos}}:
\Omega_E^{\mathrm{ch}}C_X^{\mathrm{geom}}
\xrightarrow{\simeq}\mathsf{Den}_{\Delta_5,E}.
\]

\item Output
\((C_{X,R}^{\mathrm{geom}},F^{\mathrm{co}}_{R,\bullet},
\Delta_R^{\mathrm{ch}},b_{X,R},\Theta_{\mathrm{Kos},R})\) at every
\(R\), with successor compatibility.
\end{enumerate}

\paragraph{Type checks.}
\(C_{X,R}^{\mathrm{geom}}\) is a coaugmented conilpotent chiral
coalgebra on \(E\) with bounded bar-length filtration of length at
most \(N_R\); \(\Delta_R^{\mathrm{ch}}\) is coassociative and counital
on the chosen flag stacks; the bar-cobar inversion is supplied by
Francis--Gaitsgory chiral Koszul; the cobar comparison
\(\Theta_{\mathrm{Kos},R}\) is source-internal, not the target
bar-cobar counit (the source/target separation of
Lemma~\ref{lem:source-target-koszul-separation}).

\paragraph{Obstructions discharged or invoked.}
The Koszul Mittag--Leffler exactness hypothesis is the
\textup{(D0)}-component for the chiral Koszul lane.  The total finite
Koszul residual
\(\mathfrak O_{\mathrm{Kos},R}=
(\mathfrak o^C_R,\mathfrak o^{\mathrm{fil}}_R,
\mathfrak o^{\mathrm{conil}}_R,\mathfrak o^{\Delta,\mathrm{ch}}_R,
\{\mathfrak o^{C,\mathrm{tr}}_{R'R}\},\mathfrak o^b_R,
\mathfrak o^\Omega_R,\mathfrak o^{\mathrm{hyb}}_R,
\mathfrak o^\Pi_R,\mathfrak o^{\mathrm{sep}}_R,
\mathfrak o^{\Prim}_R)\)
vanishes after this step.

\paragraph{Cross-volume firewall.}  The bar coalgebra
\(\bar B_E^{\mathrm{ch}}
(\mathcal F_{X,R}^{\mathrm{hyb}})\) is a \emph{twisting / coupling}
coalgebra (Maurer--Cartan classification), \emph{not} the bulk
\(Z^{\mathrm{der}}_{\mathrm{ch}}(\mathcal F_X^{\mathrm{hyb}})
=\mathrm{ChirHoch}^\bullet(\mathcal F_X^{\mathrm{hyb}},
\mathcal F_X^{\mathrm{hyb}})\); cf.\
\texttt{chiral-bar-cobar-vol2:CLAUDE.md~§6} (universal stage chain
\(\mathsf P\to\mathsf C\to\mathsf S\to\mathsf Z\)).
The chiral Koszul comparison
\(\Theta_{\mathrm{Kos}}:\Omega_E^{\mathrm{ch}}C_X^{\mathrm{geom}}
\xrightarrow{\simeq}\mathsf{Den}_{\Delta_5,E}\) is a comparison
between two algebras of the same level; the homotopy inverse to the
target-internal counit
\(\Omega_E^{\mathrm{ch}}\bar B_E^{\mathrm{ch}}
(\mathsf{Den}_{\Delta_5,E})\to\mathsf{Den}_{\Delta_5,E}\) is not used
in either direction.

\subsection{Construction~8 — Primitive recognition}
\label{app:build-recognition}

\paragraph{Tag.} \(\delta\)-endpoint convergence
(source-to-target identification).

\paragraph{Input.}
The descended source primitive Lie superalgebra
\(P_X^{\Pi,\mathrm{red}}\) of
Construction~\ref{app:build-hall} (after the radical quotient
\(\overline{\mathrm{Rad}\langle\,\cdot,\cdot\rangle_X}\)); the imported
target denominator algebra
\(\fg_{\Delta_5}=G(\mathscr B_{\Delta_5})\) on
\(\La^{2,1}_{II}\) of Definition~\ref{def:bkm-imaginary-simple-roots} and
Definition~\ref{def:bkm-relations}; the increasing cofinal
sequence of finite downward-saturated windows
\(W_1\subset W_2\subset\cdots\) with
\(\bigcup_\nu W_\nu=\mathcal R_+\); the Borcherds--Kac--Moody data
\((H,I,p,\alpha_i,(\,,\,))\) of
Theorem~\ref{thm:primitive-recognition}.

\paragraph{Output.}
The cofinal finite-window primitive-recognition datum
\(\{\textup{(R0),(R1),(R2),(R3),(R4),(R5),(R6),(R7),(R8)}\}\) of
Definition~\ref{def:cofinal-primitive-recognition-datum} for every
\(W_\nu\), compatible under \(W_\nu\subset W_{\nu+1}\); the source
representatives
\(e_i^X,f_i^X,h_i^X\); the source presentation map
\(\pi_W:\widehat{\mathfrak F}(\mathscr B_{\Delta_5})_W
\to P_X^{\Pi,\mathrm{red}}\big|_W\); the kernel equality
\(\ker\pi_W=
(\overline{J_{\mathrm{BK}}+\mathrm{Rad}_{\mathrm{GN}}})_W\); the
isomorphism
\[
\boxed{\quad
P_X^{\Pi,\mathrm{red}}\cong\fg_{\Delta_5},\quad
\Omega_E^{\mathrm{ch}}C_X\simeq\mathsf{Den}_{\Delta_5,E}.
\quad}
\]

\paragraph{Procedure.}
For every cofinal \(W_\nu\):
\begin{enumerate}
\item[\textup{(R0)}] \emph{Compact source labels.}  Verify every
finite source basis vector in \(W_\nu\cup(-W_\nu)\cup\{0\}\) satisfies
the compact source-admissibility condition of
Theorem~\ref{thm:finite-stage-dirac-igusa-pro-object}.

\item[\textup{(R1)}] \emph{Representatives and parities.}  Supply
parity-homogeneous primitive representatives
\(e_i^X\in P^{\Pi,\mathrm{red}}_{X,\alpha_i}\),
\(f_i^X\in P^{\Pi,\mathrm{red}}_{X,-\alpha_i}\),
\(h_i^X=\iota_H^{-1}(\alpha_i)\) for every simple-root index \(i\) with
\(\alpha_i\in W_\nu\), with the Gritsenko--Nikulin parity fibres
\(\mu_{\overline 0}(a),\mu_{\overline 1}(a)\) for imaginary simple
roots and even real representatives for \(\delta_1,\delta_2,\delta_3\).

\item[\textup{(R2)}] \emph{Relations.}  Verify in the protected Hall
super-bracket the Cartan, Chevalley, real Serre, Borcherds
orthogonality, and super sign relations for every relation whose
displayed source and target degrees lie in
\(W_\nu\cup(-W_\nu)\cup\{0\}\).  In the real rank-three subsystem the
Cartan matrix is
\(((\delta_i,\delta_j))=\mathrm{diag}(2)-2(J-\mathrm{diag}(1))\) with
real Serre exponent~3.

\item[\textup{(R3)}] \emph{Full finite parity dimensions.}  For every
\(\beta\in W_\nu\) verify
\(\dim P^{\Pi,\mathrm{red}}_{X,\beta,\overline p}
=\rootmult_{\overline p}^{\mathrm{GN}}(\beta)\) for both parities
\(p=0,1\).  In the first timelike channel
\(\beta=\delta_{123}\),
Proposition~\ref{prop:bkm-delta123-presentation-count} requires
\((d^{\mathrm{GN}}_{(1,1,1),\overline 0},
d^{\mathrm{GN}}_{(1,1,1),\overline 1})=(29,93)\).  Determinant
information \(29-93=-64=f(1,1)\) is consistent but not the
recognition.

\item[\textup{(R4)}] \emph{Hopf pairing and radical.}  Compute the
positive-negative pairing matrices in parity blocks; verify the
kernel agrees with the
Gritsenko--Nikulin/Kac invariant-pairing kernel; verify Lie-ideal and
coideal stability under all finite brackets; verify carrying under
HN transitions to the next window.

\item[\textup{(R5)}] \emph{No extra relations.}  Verify the kernel
equality
\(\ker\pi_\nu=(\overline{J_{\mathrm{BK}}+\mathrm{Rad}_{\mathrm{GN}}})_{\le W_\nu}\) on the finite triangular span generated by
\(W_\nu\cup(-W_\nu)\cup\{0\}\).

\item[\textup{(R6)}] \emph{Generation.}  Verify that every source
primitive root space in \(W_\nu\cup(-W_\nu)\), after the finite
radical quotient, is generated by the Cartan and the supplied
simple-root representatives through brackets whose intermediate
positive degrees stay in \(W_\nu\).

\item[\textup{(R7)}] \emph{Completed PBW compatibility.}  Verify the
PBW-graded isomorphism
\(\mathrm{gr}_{\mathrm{PBW}}U(P_X^{\Pi,\mathrm{red}})_{\le W_\nu}
\xrightarrow{\sim}
\mathrm{gr}_{\mathrm{PBW}}U(\fg_{\Delta_5})_{\le W_\nu}\); verify
strict preservation under \(W_{\nu+1}\to W_\nu\) transition maps.

\item[\textup{(R8)}] \emph{Exact triangular completion.}  Verify
strictness of the transition maps and closedness of images on every
finite window, so that passage to the inverse limit commutes with the
finite kernels, cokernels, radicals, PBW associated gradeds, and
parity pieces.  Equivalently the relevant \(\lim^1\) terms vanish for
these triangular finite-window systems.

\item Output the global isomorphism
\(P_X^{\Pi,\mathrm{red}}\cong\fg_{\Delta_5}\) by passage to the inverse
limit (compatible cofinal datum); together with the chiral Koszul
comparison of Construction~\ref{app:build-koszul} this gives
\(\Omega_E^{\mathrm{ch}}C_X\simeq\mathsf{Den}_{\Delta_5,E}\) and hence
the protected Pfaffian identification
\(\operatorname{Pf}_{\mathrm{prot}}(\mathfrak D_X^{\mathrm{DI}})=\Delta_5\).
\end{enumerate}

\paragraph{Type checks.}
The source presentation map \(\pi:\widehat{\mathfrak F}(\mathscr B_{\Delta_5})\to P_X^{\Pi,\mathrm{red}}\) is
continuous and degree-preserving; \textup{(R5)} gives injectivity
after the completed radical quotient; \textup{(R6)} gives
surjectivity; \textup{(R7)} gives PBW compatibility; \textup{(R8)}
gives exactness in the inverse-limit step.  No claim on signed
superdimensions \(\sdim=f(nm,l)\) substitutes for the two-integer
parity dimensions of \textup{(R3)}.

\paragraph{Obstructions discharged or invoked.}
This step inherits the obstruction tower discharged by
Constructions~1--7 and adds the cofinal exact-completion residual
\textup{(R8)}.  The Beilinson--Drinfeld bar-cobar firewall is
respected throughout: \textup{(R5)} is a source kernel-equality
theorem, not an inversion of the target bar-cobar counit;
\textup{(R7)} is a PBW graded comparison of source and target
enveloping algebras, not a graph-equality argument from the scalar
square.

\paragraph{Cross-volume anchor.}  The output is
\texttt{chiral-bar-cobar-vol2:CLAUDE.md~§9} Construction
Problem~1: the operator-level construction of
\(\mathfrak D_X^{\mathrm{DI}}\) for \(K3\times E\) with
\(\operatorname{Pf}_{\mathrm{prot}}(\mathfrak D_X^{\mathrm{DI}})=\Delta_5\), under the
four obstructions \((D_0),(O1),(O1)^+,(O2)\).  The Drinfeld doubling
\(Y^+(X)\to G(X)=D(Y^+(X))\) of
\texttt{calabi-yau-quantum-groups:CLAUDE.md} requires the Hopf
pairing of \textup{(R4)}, completion by \textup{(R8)}, integral form,
and stable-envelope transport which are supplied by the same
data.  The κ-tuple coordinate \(\kappa_{\mathrm{BKM}}=5\) is the
weight of \(\Delta_5\) and lives at level~4 of the universal stage
chain; it is a consequence of the construction, not an input to it.

\subsection{The eight-step audit}
\label{app:build-audit}

The eight construction steps form a strict ordered sequence of
input/output type discipline.  An audit of any candidate construction of
\(\mathfrak D_X^{\mathrm{DI}}\) reduces to verifying, for every
\(R\in\mathcal R_{\mathrm{HN}}\):
\begin{enumerate}
\item the lattice combinatorics of
Construction~\ref{app:build-charge-window};
\item the finite-type quasi-smooth derived stacks of
Construction~\ref{app:build-finite-moduli};
\item the hybrid local/wrapped factorization datum and the eight-word
two-step flag atlas of Construction~\ref{app:build-hybrid-wrapped};
\item the four Cech--Borel orientation classes and the
\((O1)^+\)~Coxeter coherence of
Construction~\ref{app:build-orientation};
\item the finite Hall bialgebra and primitive matrices of
Construction~\ref{app:build-hall};
\item the rank-one Pfaffian crossing and the half-Hilbert orbit
transport of Construction~\ref{app:build-dirac-pfaffian};
\item the bar coalgebra, conilpotent collision coproduct, and the
Francis--Gaitsgory cobar quasi-isomorphism of
Construction~\ref{app:build-koszul};
\item the cofinal datum
\(\textup{(R0)}\)--\(\textup{(R8)}\) of
Construction~\ref{app:build-recognition};
\end{enumerate}
together with the compatible vanishing of the obstruction tuples
\((\mathfrak O^{E_3-\mathrm{src}}_R,
\mathfrak O_{\mathrm{hyb},R},
\mathfrak O_{\mathrm{Kos},R},
\mathfrak O^{\mathrm{O1}}_R,
\mathfrak O^{\mathrm{O1}^+}_R,
\mathfrak O^{\mathrm{O2}}_R,
\mathfrak O_R^{D_0})\) at every height \(R\), and the
Mittag--Leffler exactness condition \textup{[ML]} of
Definition~\ref{def:O1-strong-reduced-orientation} on the cofinal HN inverse
system.  Each obstruction is attributable to the unique step
in which its components are constructed.

\paragraph{Order rigidity.}  Each of the seven adjacent input/output
dependencies fails type-checking under exchange:
\begin{enumerate}
\item charge window $\to$ finite moduli: $\Lambda_{\mathrm{ample}}$
  parametrises the HN strata; without the window, the moduli stack
  has no finite cofinal subsystem.
\item finite moduli $\to$ hybrid carrier: $\mathrm{Ran}^{\mathrm{hyb}}(E)$
  needs the finite Hall stack as factor on which to graft the
  wrapped/local atlas.
\item hybrid carrier $\to$ reduced orientation: the cosection-twisted
  Joyce--Upmeier line bundle is defined on the carrier, not on the
  raw moduli (the cosection lives on $\mathrm{Ran}^{\mathrm{hyb}}$).
\item reduced orientation $\to$ Hall correspondences: the Pfaffian
  line bundle on the Hall stack is the orientation-twisted sign data;
  without orientation, the Hall product is unsigned.
\item Hall correspondences $\to$ Pfaffian line + Dirac operator: the
  primitive Hall bracket $[-,-]_X$ defines the Hall--Drinfeld
  correspondence whose Pfaffian section is the Dirac--Igusa Pfaffian.
\item Pfaffian line $\to$ chiral Koszul source coalgebra: the source
  $C_X$ is the bar of the augmented hybrid factorization algebra
  pulled back along the Pfaffian comparison
  $\iota_{\mathrm{aut}}$.
\item chiral Koszul source $\to$ primitive recognition: the primitive
  part $P_X^{\Pi,\mathrm{red}}$ is the Frobenius--Serre primitive
  filtration on the source, which the recognition theorem matches to
  $\mathfrak g_{\Delta_5}$.
\end{enumerate}
Two diagnostic exchanges: orientation before hybrid leaves the middle
type-II wall \(\delta_2\leftrightarrow(0,1,1)\) without a
representative, by
Lemma~\ref{app:pfaffian-wall-three-walls}; recognition before the
Koszul cobar leaves the identification
\(\Omega_E^{\mathrm{ch}}C_X\simeq\mathsf{Den}_{\Delta_5,E}\)
unconstructed.  All other adjacent exchanges fail by the same
input/output type criterion.

\paragraph{Conditional theorem.}  Granted the eight constructions and
the four obstructions \((D_0),(O1),(O1)^+,(O2)\),
Theorem~\ref{thm:pfaffian-dirac-finite-stage} reads
\[
\operatorname{Pf}_{\mathrm{prot}}(\mathfrak D_X^{\mathrm{DI}})=\Delta_5,\qquad
\epsilon_o=\nu_{\Delta_5}\big|_{W^{(2)}(\La^{2,1}_{II})},\qquad
P_X^{\Pi,\mathrm{red}}\cong\fg_{\Delta_5}.
\]
The scalar square
\(Z_{\mathrm{BPS}}^{K3\times E}=
(\Phi_{10}^{\mathrm{un}})^{-1}=\Delta_5^{-2}\) is the level-4 trace
of the constructed protected algebra; promotion to the 3d
gravitational path integral is not claimed and requires the
Hall--Borcherds residual of
\texttt{chiral-bar-cobar-vol2}.

%%% End of Appendix F.
 after
% \appendix has been opened at main.tex line 17081, so that
% \thesection = \Alph{section}.  No \chapter, no second \appendix.

\section{Eight build algorithms for the Dirac--Igusa pro-object}
\label{app:eight-build-algorithms}

The Dirac--Igusa pro-object
\[
\mathfrak D_X^{\mathrm{DI}}
=\varprojlim_R
\bigl(
\mathcal A^{E_3}_{X,R},
\mathcal F_{X,R}^{\mathrm{hyb}},
\Gamma_{X,R},
\Pi_X,
\Phi_R,
o_R,
\mathcal H_R,
P_R^\Pi,
C_{X,R},
\Theta_{\mathrm{Kos},R},
\mathscr L_{\mathrm{Pf},R},
\operatorname{pf}_{X,R},
\epsilon_{o,R},
\mathrm{Rec}_R
\bigr)
\]
of Definition~\ref{thm:finite-stage-dirac-igusa-pro-object} and
Definition~\ref{def:O1-strong-reduced-orientation} is the conditional source
operator-level package whose protected Pfaffian, after the four
obstructions \((D_0)\), \((O1)\), \((O1)^+\), \((O2)\) are discharged,
realizes the Igusa cusp form,
\[
\operatorname{Pf}_{\mathrm{prot}}(\mathfrak D_X^{\mathrm{DI}})=\Delta_5,\qquad
\epsilon_o=\nu_{\Delta_5}\big|_{W^{(2)}(\La^{2,1}_{II})},\qquad
P_X^{\Pi,\mathrm{red}}\cong\fg_{\Delta_5}.
\]
Construction at finite Harder--Narasimhan height \(R\) proceeds through
eight construction steps.  Each step is a procedure on input data
already supplied by earlier steps and feeds named output data into the
next.  The eight steps stratify the realization datum
\(\mathfrak K_X=(L_X,H_X,O_X,D_X,R_X)\) of
Definition~\ref{thm:finite-stage-dirac-igusa-pro-object}: steps 1--3
realize \((L_X)\) and \((H_X)\); step 4 realizes the local content of
\((O_X)\); step 5 produces the Hall correspondence object; step 6 fixes
the Pfaffian--Dirac line and feeds the orientation character; steps 7--8
realize \((D_X)\) and \((R_X)\).

Throughout, \(X=S\times E\) with \(S\) a smooth complex K3 surface and
\(E\) a smooth elliptic curve.  \(\sigma=\sigma_{s,t}\) is Liu's product
stability condition on the Abramovich--Polishchuk heart in
\(\mathsf{D}^b(\mathrm{Coh}(X))\)
\cite{AbramovichPolishchukTStructures,LiuProductStability}.  The HN
sector \((\sigma,S)\) is fixed and the cofinal subsystem
\(\mathcal R_{\mathrm{HN}}=\{R_0<R_1<\cdots\}\) is the one of
Definition~\ref{def:O1-strong-reduced-orientation}, governed by
\textup{[K3$\times$E-fin]}, \textup{[K3$\times$E-atlas]}, and
\textup{[ML]}.

Each construction step carries a licensing tag pulled from the
five-type universal stage chain
\(\mathsf P\to\mathsf C\to\mathsf S\to\mathsf Z\to\mathsf A\) of
\texttt{chiral-bar-cobar-vol2:CLAUDE.md~§9}: \(\alpha\) primitive,
\(\beta\) functorial passage, \(\gamma\) shadow specialization,
\(\delta\) endpoint convergence, \(\epsilon\) terminal scalar.  The
ordering charge\(\to\)moduli\(\to\)hybrid\(\to\)orientation\(\to\)Hall
\(\to\)Dirac+Pfaffian\(\to\)Koszul\(\to\)recognition is the unique
ordering compatible with input/output type discipline; reordering any
two adjacent steps fails an input check.

The four obstructions are attributed to specific steps:
\((D_0)\) is the cofinal-trivialization datum of the entire eight-step
construction at every finite \(R\) (steps 1--6 supply its components);
\((O1)\) lives at step 4; \((O1)^+\) lives at step 4 with a Coxeter
coherence check feeding step 6; \((O2)\) lives at step 6.  Step 8
inherits the obstruction tower discharged by steps 1--7 and adds the
exact-completion residual \textup{(R8)} of
Definition~\ref{def:cofinal-primitive-recognition-datum}.

\subsection{Construction~1 — Charge window}
\label{app:build-charge-window}

\paragraph{Tag.} \(\alpha\)-primitive (lattice datum).

\paragraph{Input.}
HN height \(R\in\mathcal R_{\mathrm{HN}}\); the formal dyonic Mukai
lattice \(\Gamma_X^{\mathrm{form}}\) of
Definition~\ref{def:gram-index-lattice}; the Gram map
\(\Pi_X:\Gamma_X^{\mathrm{form}}\to\Gamma_{\mathrm{gram}}\) and the
cocycle \(B\) of Lemma~\ref{lem:gram-additivity-defect}; the active
support \(\Gamma_{\mathrm{act}}=\{\gamma\in\Gamma_{\mathrm{eff}}\mid
f(nm,l)\ne0\}\) of the Igusa Jacobi form \(\phi_{0,1}\); the cofinal
subsystem \(\mathcal R_{\mathrm{HN}}\).

\paragraph{Output.}
A finite cusp-positive active window
\[
A_R\subset\Gamma_{\mathrm{act}}\subset\Gamma_{\mathrm{gram}},
\qquad
A_R\hbox{ downward saturated for the cofinal target-window order},
\]
together with the normal-ordered extension
\[
\widehat\Gamma_R=
\{(c,T)\mid c\in\Gamma_R^{HN},\ T\in\mathcal T_R(c)\}
\subset\widehat\Gamma_X,
\qquad
\Gamma_R^\Pi=\overline\Pi_X(\widehat\Gamma_R)\subset\Gamma_{\mathrm{gram}},
\]
the lift \(L:\Gamma_{\mathrm{gram}}\to\Gamma_X^{\mathrm{form}}\) of
Lemma~\ref{lem:primitive-lift-every-gram-triple}, and the finite test
window \(\Gamma_R^{\mathrm{test}}\) of
Definition~\ref{def:O1-strong-reduced-orientation}.

\paragraph{Procedure.}
\begin{enumerate}
\item Compute the active support cone
\[
\Gamma_{\mathrm{act}}=\{(n,l,m)\in\Gamma_{\mathrm{eff}}\mid f(nm,l)\ne 0\}
\]
from the Fourier expansion of the index-one weak Jacobi form
\(\phi_{0,1}(\tau,z)=\sum_{n,l}c(n,l)q^nr^l\) of
\cite{EZ:1}, with the index substitution
\(f(nm,l):=c(nm,l)\) tabulated in
Appendix~\ref{sec:n1-boundary-compatibility-conditions} for the
\(N=1\) Igusa boundary.

\item Set
\[
A_R^{\circ}
=\{\gamma=(n,l,m)\in\Gamma_{\mathrm{act}}\mid n+|l|+m\le R\}.
\]
Form its downward saturation
\[
A_R
=\{\gamma'\in\Gamma_{\mathrm{act}}\mid
\exists\gamma\in A_R^{\circ},\ 0<\gamma'\le\gamma,\
\gamma-\gamma'\in Q_+\},
\]
in the sense of Definition~\ref{def:cofinal-primitive-recognition-datum}.

\item Normal-ordered lift.  For each \(c\in\Gamma_R^{HN}\) compute the
\(B\)-translate orbit
\[
\mathcal T_R(c)
=\Bigl\{\sum_{i<j}B(c_i,c_j)\Bigm|c=\sum_i c_i,\
c_i\in F_{\sigma,S}(R),\ \sum_i h_S(c_i)\le R\Bigr\}.
\]
Set \(\widehat\Gamma_R=\{(c,T):c\in\Gamma_R^{HN},\ T\in\mathcal T_R(c)\}\).

\item Finite target test window:
\(\Gamma_R^{\mathrm{test}}=\overline\Pi_X(\widehat\Gamma_R)\), compatibly
with \(R\to R'\) by inclusion.

\item Output \((A_R,\widehat\Gamma_R,\Gamma_R^\Pi,\Gamma_R^{\mathrm{test}},L)\).
\end{enumerate}

\paragraph{Type checks.}
\(A_R\) is finite and downward saturated;
\(\overline\Pi_X(\widehat\Gamma_R)=\Gamma_R^\Pi\) lies in
\(\Gamma_{\mathrm{gram}}\); \(\Gamma_R^{HN}\subset
F_{\sigma,S}(R)\) is finite by \textup{[K3$\times$E-fin]}; the cocycle
identity
\(\Pi_X(c+c')=\Pi_X(c)+\Pi_X(c')+B(c,c')\) is an axiom of the cocycle.

\paragraph{Obstructions discharged or invoked.}
None; this step is purely lattice combinatorics from the imported
Gritsenko--Nikulin Borcherds product
\cite{GN,GNPartI,GNII,BorcherdsGKM88,Borcherds95}.
The downward saturation is the input requirement of \textup{(R0)} in
Definition~\ref{def:cofinal-primitive-recognition-datum}.

\subsection{Construction~2 — Finite moduli}
\label{app:build-finite-moduli}

\paragraph{Tag.} \(\alpha\)-primitive (geometric source).

\paragraph{Input.}
The active window \(A_R\) and lift \(L\) of
Construction~\ref{app:build-charge-window}; Liu's product stability
condition \(\sigma_{s,t}\); the standing hypotheses
\textup{[K3$\times$E-fin]} and \textup{[K3$\times$E-atlas]}.

\paragraph{Output.}
A finite class set \(C_R\) and lift
\(\ell_R:\Gamma_R^\Pi\to C_R\) with
\(\Pi_X(\ell_R(\gamma))=\gamma\); the family of finite-type quasi-smooth
derived Artin substacks
\[
\mathfrak M^{ss}_{R,c}
=\mathfrak M^{ss}_{\sigma_{s,t}}(X,c)_{\le R}
\subset\operatorname{Perf}(X),
\qquad c\in C_R,
\]
of \(\sigma_{s,t}\)-semistable objects of full Liu numerical class
\(c\) and HN height \(\le R\); the scalar-rigidified truncation
\(\mathfrak M^{\mathrm{sc}}_{R,c}\); the universal perfect complex
\(\mathcal A_{R,c}\) on \(X\times\mathfrak M^{\mathrm{sc}}_{R,c}\); the
finite stratification by retained closed substacks of finite residual
inertia.

\paragraph{Procedure.}
\begin{enumerate}
\item Form \(C_R=\{c=L\gamma:\gamma\in A_R\}\) closed under retained
extension sums (taken in \(\Gamma_X^{\mathrm{form}}\)).

\item For each \(c\in C_R\) cut the substack of
\(\sigma_{s,t}\)-semistable objects of class \(c\) by HN height \(\le
R\):
\[
\mathfrak M^{ss}_{R,c}=
\bigl\{[A]\in\operatorname{Perf}(X)\mid
\sigma_{s,t}\hbox{-semistable},\
\hn(A)\le R,\
\nu(A)=c\bigr\}.
\]

\item Verify finite-typeness: \textup{[K3$\times$E-fin]} (Liu product
stability + product-boundedness theorem at fixed full Liu charge with
finite fibres over the active charge) gives the finite-type assertion.
PTVV \cite{PTVV2013} on the Calabi--Yau threefold \(X\) gives the
\((-1)\)-shifted symplectic structure on
\(\operatorname{Perf}(X)\), restricted to the retained semistable
sector.

\item Quasi-smoothness: the universal perfect complex
\(\mathcal A_{R,c}\) supplies a relative cotangent complex of bounded
Tor-amplitude on the retained charts.  Tangent at \(A\) is
\(T_A\mathfrak M^{ss}_{R,c}\simeq\operatorname{RHom}_X(A,A)[1]\).

\item Scalar rigidification: form \(\mathfrak M^{\mathrm{sc}}_{R,c}\) by
killing the scalar \(\mathbb G_m\)-automorphism of
\(c\)-stable factors.

\item Two-step flag stack
\(\mathfrak F^{(2)}_{R,c,c',c''}\) and extension stack
\(\mathfrak E_{R,c,c'}\) constructed as iterated derived vector stacks
of \(R\pi_*\mathcal RHom_X(\mathcal B,\mathcal A[1])\); both finite type
and quasi-smooth in the retained sector by
Proposition~\ref{prop:finite-dcritical-hall-product}.

\item Output \((C_R,\ell_R,\{\mathfrak M^{ss}_{R,c}\},
\{\mathfrak E_{R,c,c'}\},\{\mathfrak F^{(2)}_{R,c,c',c''}\})\) with
universal perfect complexes and the finite stratification.
\end{enumerate}

\paragraph{Type checks.}
\(\mathfrak M^{ss}_{R,c}\) is a finite-type quasi-smooth derived Artin
substack of \(\operatorname{Perf}(X)\); the scalar rigidification leaves
only finite residual inertia
\(H_S\subset E_{\mathrm{tors}}\) on each connected stabilizer
component, by
Proposition~\ref{prop:finite-dcritical-hall-product};
PTVV \((-1)\)-shifted symplectic structure restricts.

\paragraph{Obstructions discharged or invoked.}
\textup{[K3$\times$E-fin]} is invoked here in its product-boundedness
form and supplies the finite-type assertion not in the cited literature
for the K3\(\times\)E threefold; it is part of the standing datum.
\textup{[K3$\times$E-atlas]} supplies the quasi-smooth d-critical
charts.  These enter \((D_0)\) at the level of the finite Hall atlas
of Theorem~\ref{thm:finite-stage-dirac-igusa-pro-object}.

\subsection{Construction~3 — Hybrid wrapped factorization carrier}
\label{app:build-hybrid-wrapped}

\paragraph{Tag.} \(\beta\)-functorial passage (Stage-1 carrier).

\paragraph{Input.}
The finite moduli stacks
\(\{\mathfrak M^{\mathrm{sc}}_{R,c}\}\) of
Construction~\ref{app:build-finite-moduli}; the
projection-finite/wrapped split by elliptic degree
\(b_R^{\mathrm{geom}}:\Gamma_{\sigma,S}^{HN}(R)\to\Z_{\ge0}\) of
Definition~\ref{def:hybrid-wrapped-factorization-datum}; a fixed
origin \(0_E\in E\); the Abramovich--Polishchuk/Liu wrapped anchor
candidate.

\paragraph{Output.}
The finite hybrid local/wrapped factorization datum
\[
\mathcal F_{X,\sigma,S,\le R}^{\mathrm{hyb}}
=
(\mathcal C_{\sigma,S,\le R},
b_R^{\mathrm{geom}},
\mathcal F_R^{\mathrm{loc}},
\mathcal G_R^{\mathrm{wr}},
\mathfrak E_R^{\mathrm{loc}},
\mathfrak E_R^{\mathrm{hyb}},
\mathfrak E_R^{\mathrm{wr}},
\mathfrak{Flag}_R^{\mathrm{hyb}},
\mathsf{Corr}_R^{E,\mathrm{hyb}},
\lambda_R,
Q_{E,R},
\theta^Q_R,
o_R,
I_R^{\mathrm{prot}})
\]
of Definition~\ref{def:hybrid-wrapped-factorization-datum},
together with the hybrid Ran base
\(\operatorname{Ran}^{\mathrm{hyb}}(E)\), the eight-word two-step
binary flag atlas \(\{w\in\{\mathrm L,\mathrm W\}^3\}\), and the
quotient-after-correspondence functor \(Q_{E,R}\).

\paragraph{Procedure.}
\begin{enumerate}
\item \emph{Elliptic-degree split.}  Decompose
\(\Gamma_{\sigma,S}^{HN}(R)
=\Gamma_R^{\mathrm{loc}}\sqcup\Gamma_R^{\mathrm{wr}}\) where
\(\Gamma_R^{\mathrm{loc}}=\{b_R^{\mathrm{geom}}=0\}\) and
\(\Gamma_R^{\mathrm{wr}}=\{b_R^{\mathrm{geom}}>0\}\).  Verify additivity
\(b_R^{\mathrm{geom}}(\gamma+\gamma')
=b_R^{\mathrm{geom}}(\gamma)+b_R^{\mathrm{geom}}(\gamma')\) on retained
extensions remaining in the quotient.

\item \emph{Projection-finite local sector.}  For
\(\alpha\in\Gamma_R^{\mathrm{loc}}\) and finite index set \(I\), build
the closed-support incidence stack over \(E^I\):
\[
\mathfrak M^{\mathrm{loc,cl}}_{\alpha,R,I}
=
\Bigl\{(A,\mathbf e)\in\mathfrak M^{\mathrm{sc}}_{R,L\alpha}\times E^I
\Bigm|p_E(\mathrm{Supp}\,A)\subset\bigcup_{i\in I}S\times\{e_i\}\Bigr\}.
\]
Form
\(\mathscr H^{\mathrm{loc}}_{\alpha,R,I}=
(\pi_{\alpha,R,I})_!(\Phi_{\alpha,R,I}^{\mathrm{loc}}\otimes
o_{\alpha,R,I}^{\mathrm{loc}})\) and the open-support sheaf
\(\mathcal F_{\alpha,R}^{\mathrm{loc}}(U)\) over \(U\subset E\) by
restricted Borel--Moore configuration sums.

\item \emph{Wrapped prequotient sector.}  For
\(\eta\in\Gamma_R^{\mathrm{wr}}\), form the rigidified prequotient
\(\mathcal M_{\eta,R}^{\mathrm{wr,rig}}\to\mathcal M_{\eta,R}^{\mathrm{wr}}\)
by the Abramovich--Polishchuk / determinant anchor
\[
\lambda_\eta:\mathcal M_{\eta,R}^{\mathrm{wr,rig}}\longrightarrow E,
\qquad
\lambda_\eta(A)=\det Rp_{E,*}A
\otimes\mathcal O_E(-\chi(A)\cdot 0_E),
\]
divided by a finite cover when \(\chi(A)\ne 1\) so that
\(\lambda_\eta(t\cdot A)=\lambda_\eta(A)+t\) holds.  Verify the
losslessness residual \(\mathfrak o^{\lambda,\mathrm{loss}}_R=0\) by
the Ext-distinguishing test of
Definition~\ref{def:hybrid-wrapped-factorization-datum}, separating
\(W_0=i_{s,*}(\mathcal O_E\oplus\mathcal O_E)\) from
\(W_L=i_{s,*}(L\oplus L^{-1})\).

\item \emph{Mixed local/wrapped correspondences.}  For an ordered
hybrid input
\(\alpha_-\in\Gamma_R^{\mathrm{loc}}\),
\(\eta\in\Gamma_R^{\mathrm{wr}}\),
\(\beta_+\in\Gamma_R^{\mathrm{loc}}\), form
\(\mathfrak E^{\mathrm{mix}}_{\alpha_-;\eta;\beta_+,R}\) before the
reduced \(E\)-quotient.

\item \emph{Eight-word two-step flag atlas.}  Build flag stacks
\(\mathfrak F^{(2),w}_R\) for each
\(w\in\{\mathrm L,\mathrm W\}^3=
\{\mathrm{LLL},\mathrm{LLW},\mathrm{LWL},\mathrm{WLL},\mathrm{LWW},
\mathrm{WLW},\mathrm{WWL},\mathrm{WWW}\}\), classifying filtrations
with three associated graded factors of those types.  Verify the
four-input pentagon coherence and the quotient-after-correspondence
descent.

\item \emph{Hybrid correspondence-module functor.}  Assemble the
correspondence category
\(\mathsf{Corr}^{E,\mathrm{hyb}}_R\) and the functor
\(\mathcal B_R^{E,\mathrm{hyb}}:\mathsf{Corr}^{E,\mathrm{hyb}}_R
\to\mathsf{Ch}_{\C}\) by
\(\mathcal B_R^{E,\mathrm{hyb}}(e)=q_!\,\mathrm{TS}^{\mathrm{red}}_e\,p^*\)
for every generating arrow
\(e=(Y\xleftarrow{p}\mathfrak E_e\xrightarrow{q}Z)\).

\item \emph{Quotient functor \(Q_{E,R}\).}  Apply
\(E\)-equivariant compactly supported vanishing-cycle chains and the
orientation descent of step~4 to produce the reduced object
\[
\mathcal F^{\mathrm{geom}}_{X,\sigma,S,\le R}
=Q_{E,R}\circ\mathcal B_R^{E,\mathrm{hyb}}.
\]

\item Output the hybrid datum with all coherences and the protected
integration map \(I_R^{\mathrm{prot}}\), with homogeneous
\(s\)-degree the Borcherds trace exponent
\(m_R=\operatorname{pr}_3\overline\Pi_X\).
\end{enumerate}

\paragraph{Type checks.}
The carrier
\(\mathcal F_{X,\sigma,S,\le R}^{\mathrm{hyb}}\) is a hybrid
correspondence-module functor on
\(\mathsf{Corr}^{E,\mathrm{hyb}}_R\) with values in
\(\mathsf{Ch}_{\C}\); the four-input pentagon is supplied by
the eight-word two-step flag atlas; the anchor residual tuple
\(\mathfrak o^\lambda_R=
(\mathfrak o^{\lambda,\mathrm{ex}}_R,
\mathfrak o^{\lambda,\mathrm{unit}}_R,
\mathfrak o^{\lambda,\mathrm{loss}}_R,
\mathfrak o^{\lambda,\mathrm{multi}}_R,
\mathfrak o^{\lambda,\mathrm{tr}}_{R'R})\) vanishes; the hybrid
residual \(\mathfrak O_{\mathrm{hyb},R}=0\) at every height \(R\),
compatibly with transition maps.

\paragraph{Obstructions discharged or invoked.}
The middle type-II wall lies in
\(\delta_2\leftrightarrow(0,1,1)\), elliptic degree
\(d=m+1=2\) by Lemma~\ref{app:pfaffian-wall-three-walls}; it
cannot be represented by a projection-finite local stack and demands
the wrapped prequotient.  This step's correctness — the existence of
the wrapped representative
\(x_{\delta_2,R}\in\mathcal M_{\eta,R}^{\mathrm{wr,rig}}\) with
\(\overline\Pi_X(x_{\delta_2,R})=(0,1,1)\) and its compatible anchor
\(\lambda_{\eta,R}\) — is precisely the hybrid input to \((O2)\) at
the middle wall.

\subsection{Construction~4 — Reduced d-critical orientation
and Joyce--Upmeier transport}
\label{app:build-orientation}

\paragraph{Tag.} \(\beta\)-functorial passage (oriented carrier).

\paragraph{Input.}
The hybrid carrier
\(\mathcal F_{X,\sigma,S,\le R}^{\mathrm{hyb}}\) of
Construction~\ref{app:build-hybrid-wrapped}; the K3 holomorphic
\(2\)-form \(\sigma_S\) on \(S\); BBDJS d-critical chart machinery
\cite{BBDJS2015}; the Joyce--Upmeier multiplicative orientation
formalism \cite{JoyceUpmeier,JoyceTanakaUpmeier}; the cosection
localization machinery of Kiem--Li and the cosection-reduced
obstruction theory; the type-II Weyl group
\(W^{(2)}(\La^{2,1}_{II})\) generated by
\(s_{\delta_1},s_{\delta_2},s_{\delta_3}\).

\paragraph{Output.}
The cosection
\(\sigma=p_S^*\sigma_S\); the reduced tangent complex
\[
\operatorname{RHom}_{\mathrm{red}}(A,A)
=
\operatorname{Cone}\!\Bigl(
\operatorname{RHom}_X(A,A)
\xrightarrow{(\mathrm{tr},\mathrm{cs})}
R\Gamma(X,\mathcal O_X)\oplus H^2(S,\mathcal O_S)[-1]
\Bigr)[-1];
\]
the BBDJS vanishing-cycle complex \(\Phi_{R,c}\); the reduced
orientation line
\[
o_{R,c}
=\det\operatorname{Ext}^1_{\mathrm{red}}(A,A)^{-1}
\quad\text{on Darboux charts},\qquad
o_{R,c}^{\otimes 2}\simeq\det\operatorname{RHom}_{\mathrm{red}}(A,A);
\]
the Joyce--Upmeier multiplicative transport
\(o_{c\oplus c'}\simeq o_c\otimes o_{c'}\) compatible with the reduced
extension correspondences; the quotient orientation
\(\mathfrak Q^{or}_R\) of
Theorem~\ref{prop:generic-local-sign-orientation-character}; the
Weyl-equivariant lifts
\(\tau_{i,R,S}:o_{R,S}\to s_{\delta_i}^*o_{R,S}\) for \(i=1,2,3\) of
\((O1)^+\); the local generic wall-sign data \(\epsilon_{o,R}^{\mathrm{loc}}\).

\paragraph{Procedure.}
\begin{enumerate}
\item \emph{Cosection.}  Pull back the K3 semiregularity cosection
\(\mathrm{cs}\) along \(p_S:X\to S\) to obtain the 2-step cosection
\(\sigma=p_S^*\sigma_S\) on the K3 factor; verify additivity in exact
triangles by \textup{(RedOr)} of
Theorem~\ref{thm:finite-stage-dirac-igusa-pro-object}.

\item \emph{Reduced obstruction.}  Form the reduced tangent complex by
the displayed cone above (Kiem--Li cosection localization).  PTVV
\cite{PTVV2013} \((-1)\)-shifted symplectic on
\(\operatorname{RHom}_X(A,A)\) descends, after the cosection
contraction, to the reduced complex.

\item \emph{BBDJS d-critical chart.}  On each Darboux chart of
\(\mathfrak M^{\mathrm{sc}}_{R,c}\), build the BBDJS vanishing-cycle
complex \(\Phi_{R,c}\) \cite{BBDJS2015}; verify reduced
Thom--Sebastiani transport coherent on two-step flag stacks of
Construction~\ref{app:build-hybrid-wrapped}.

\item \emph{Reduced orientation line.}  Define
\(o_{R,c}=\det\operatorname{Ext}^1_{\mathrm{red}}(A,A)^{-1}\); verify
the squared isomorphism
\(o_{R,c}^{\otimes 2}\simeq\det\operatorname{RHom}_{\mathrm{red}}(A,A)\)
by \(3\)-Calabi--Yau Serre duality
\(\operatorname{Ext}^2_{\mathrm{red}}(A,A)
\simeq\operatorname{Ext}^1_{\mathrm{red}}(A,A)^\vee\).

\item \emph{Quotient-orientation descent on the K3\(\times\)E-quotient
self-Ext frame bundle.}  Compute the four Cech--Borel obstruction
classes for every retained stratum \(S\):
\[
\alpha^{\mathrm{red}}_{R,S}=w_2(\mathscr D_{R,S}),
\quad
\alpha^{E,\mathrm{free}}_{R,S}=
\mathrm{edge}_{2,0}(w_2^E(\mathscr D_{R,S})),
\quad
\widetilde\beta^H_{R,S}=w_2^G(\mathscr D_{R,S}),
\quad
\lambda^H_{R,S}\in H^1_G(S;\mathbb F_2),
\]
and verify simultaneous vanishing with compatible
null-trivializations on overlaps, on Hall correspondence pullbacks,
on two-step flag stacks, and under successor maps.  The full
calculation proceeds primary stabilizer by stabilizer; for cyclic
\(\Z/2^a\) the absent mixed coefficient requires the rank-two
two-primary check.  Output the descended quotient orientation system
\(\mathfrak Q^{or}_R\).

\item \emph{Multiplicative Joyce--Upmeier transport.}  Transport
\(o_{c\oplus c'}\simeq o_c\otimes o_{c'}\) coherently across two-step
flag stacks
\cite{JoyceUpmeier,JoyceTanakaUpmeier}; verify the reduced
Thom--Sebastiani identification on three-input flags.

\item \emph{Weyl-equivariant lifts \((O1)^+\).}  For each
\(i\in\{1,2,3\}\) construct the lift
\(\tau_{i,R,S}:o_{R,S}\to s_{\delta_i}^*o_{R,S}\); verify the
projective Weyl cocycle
\(c_{o,R}\in Z^2(\mathcal G_R;\underline{\mathbb F}_2)\) on the
retained action groupoid; null-trivialize \([c_{o,R}]\) and choose a
\(1\)-cochain killing it; verify \(\tau_{i,R,S}^2=1\); verify the
Coxeter relations
\((s_{\delta_i}s_{\delta_j})^3=1\) for \(i\ne j\) (the
\((3,3,3)\)-hyperbolic triangle relations of
Appendix~\ref{app:pfaffian-wall-three-walls}); verify that the torsor
defects \(\omega_{i,\mathcal C}\) of
Definition~\ref{thm:finite-stage-dirac-igusa-pro-object}~$(O_X)$ vanish.

\item \emph{Local wall-sign datum.}  At each generic type-II wall point
\(\Hpl_{\delta_i}\), with the rank-one local model
\(K_{\delta_i}^{\mathrm{cross}}=[\R\xrightarrow{u_i}\R]\), record the
local generic sign \(\epsilon_{o,R}^{\mathrm{loc}}(s_{\delta_i})=-1\)
of Proposition~\ref{prop:rank-one-crossing-o2-sign}.

\item Output
\((o_R,\mathfrak Q^{or}_R,\{\tau_{i,R,S}\},\epsilon_{o,R}^{\mathrm{loc}})\).
\end{enumerate}

\paragraph{Type checks.}
\(o_R\) is a Picard-line section squaring to the reduced determinant;
the four Cech--Borel classes vanish with compatible
null-trivializations; the lifts \(\tau_{i,R}\) satisfy
\(\tau_{i,R}^2=1\) and the \((3,3,3)\)-Coxeter relations; multiplicative
transport agrees on flag stacks.

\paragraph{Obstructions discharged or invoked.}
This step is the entire site of \((O1)\) (strong reduced orientation
on the K3\(\times\)E-quotient self-Ext frame bundle) and of
\((O1)^+\) (Weyl-equivariant transport along
\(W^{(2)}(\La^{2,1}_{II})\) reflections).  Discharging \((O1)\)
requires the simultaneous vanishing and compatible
null-trivialization of the four Cech--Borel classes on every retained
stratum; discharging \((O1)^+\) requires \([c_{o,R}]=0\), the
Coxeter coherence, and the vanishing of all torsor defects
\(\omega_{i,\mathcal C}\).  The local wall-sign datum is the input to
\((O2)\) at step~6.

\subsection{Construction~5 — Hall correspondences and primitive part}
\label{app:build-hall}

\paragraph{Tag.} \(\gamma\)-shadow specialization
(Stage-2 chiral Hall shadow).

\paragraph{Input.}
The oriented hybrid carrier from
Construction~\ref{app:build-orientation}; the extension stack
\(\mathfrak E^{\mathrm{geom}}_{R,\widehat c,\widehat c'}\); the
two-step flag stack
\(\mathfrak F^{(2),\mathrm{geom}}_{R,\widehat c,\widehat c',\widehat c''}\);
the BBDJS reduced Thom--Sebastiani transport; six-functor admissibility
of \(p^*,q_!,p_!\) (Construction~\ref{app:build-finite-moduli}~step
\textup{(Six)}).

\paragraph{Output.}
The finite Hall object
\[
\mathcal H_R^{\mathrm{geom}}
=\bigoplus_{\widehat c\in\widehat\Gamma_R}
H_c^\bullet\!\bigl(
\mathfrak M^{\mathrm{geom}}_{R,\widehat c},
\Phi_{R,c}\otimes o_{R,c}\bigr);
\]
the Hall product \(m_R=q_!\,\mathrm{TS}^{\mathrm{red}}\,p^*\) and Hall
coproduct \(\Delta_R=p_!\,(\mathrm{TS}^{\mathrm{red}})^{-1}\,q^*\); the
finite Hall bialgebra
\((\mathcal H_R^{\mathrm{geom}},m_R,\Delta_R,\eta_R,\epsilon_R)\); the
primitive part
\[
P_R^{\mathrm{geom}}
=\ker(\Delta_R-1\otimes\mathrm{id}-\mathrm{id}\otimes 1)\cap\ker\epsilon_R;
\]
the normal-ordered Gram descent
\(P_R^{\Pi,\mathrm{geom}}
=\overline\Pi_{X,*}^{\Theta_R}P_R^{\mathrm{geom}}\); the
\(\Gamma_R^\Pi\)-graded super-Lie bracket
\([\,\bar x,\bar y\,]_{\mathrm{red}}\) of
Proposition~\ref{thm:hall-primitives-formal};
the Hopf pairing matrix \(G_{\gamma,ab}\) and radical \(K_\gamma\).

\paragraph{Procedure.}
\begin{enumerate}
\item \emph{Pullback.}  For
\(\widehat c,\widehat c'\in\widehat\Gamma_R\), pull back along
\(p:\mathfrak E^{\mathrm{geom}}_{R,\widehat c,\widehat c'}
\to\mathfrak M^{\mathrm{geom}}_{R,\widehat c}
\times\mathfrak M^{\mathrm{geom}}_{R,\widehat c'}\); apply reduced
Thom--Sebastiani \(\mathrm{TS}^{\mathrm{red}}\) to combine vanishing
cycle data \cite{BBDJS2015}; pushforward exceptionally along
\(q:\mathfrak E^{\mathrm{geom}}_{R,\widehat c,\widehat c'}
\to\mathfrak M^{\mathrm{geom}}_{R,\widehat c\star\widehat c'}\) to
obtain
\(m_R(x\otimes y)=q_!\,\mathrm{TS}^{\mathrm{red}}\,p^*(x\boxtimes y)\).

\item \emph{Coproduct.}  Dually pull back along \(q\), apply
inverse reduced Thom--Sebastiani, pushforward along \(p\) to obtain
\(\Delta_R(z)=p_!\,(\mathrm{TS}^{\mathrm{red}})^{-1}\,q^*(z)\).
The compact geometric coproduct correspondence
\(\mathfrak C^{\mathrm{geom}}_{\rho;\mu,\nu}\) of
Proposition~\ref{thm:hall-primitives-formal} is the
support data for splitting.

\item \emph{Bialgebra check.}  Verify associativity of \(m_R\) and
coassociativity of \(\Delta_R\) on the two-step flag stack
\(\mathfrak F^{(2),\mathrm{geom}}_{R,\widehat c,\widehat c',\widehat c''}\)
by base change and Thom--Sebastiani coherence.  Verify the bialgebra
identity (Hopf compatibility) via the Massey--Saito Joyce--Upmeier
multiplicative orientation transport.

\item \emph{Primitive extraction.}  Compute
\(P_R^{\mathrm{geom}}=\ker(\Delta_R-1\otimes\mathrm{id}
-\mathrm{id}\otimes 1)\cap\ker\epsilon_R\); since
\(\Delta_R-1\otimes\mathrm{id}-\mathrm{id}\otimes 1\) is degree-graded
on \(\widehat\Gamma_R\), this reduces to a finite linear computation
in each charge.

\item \emph{Normal-ordered descent.}  Apply
\(\overline\Pi_{X,*}^{\Theta_R}\) of
Theorem~\ref{thm:hall-product} to obtain the
\(\Gamma_R^\Pi\)-graded
\(P_R^{\Pi,\mathrm{geom}}\); verify additivity
\(\overline\Pi_X(\widehat c\star\widehat c')
=\overline\Pi_X(\widehat c)+\overline\Pi_X(\widehat c')\).

\item \emph{Bracket and pairing matrices.}  Compute the matrices
\[
M^{\alpha+\beta,c}_{\alpha,a;\beta,b}
=\int_{\mathfrak E^{\mathrm{geom}}_{\alpha,\beta;\alpha+\beta}}
q^!(e_{\alpha+\beta,c}^\vee)\cap
\mathrm{TS}^{\mathrm{red}}(p^*(e_{\alpha,a}\boxtimes e_{\beta,b})),
\quad
B_{\alpha,\beta}(x\otimes y)
=M_{\alpha,\beta}(x\otimes y)
-(-1)^{|x||y|}M_{\beta,\alpha}\mathsf{sw}(x\otimes y),
\]
\[
G_{\gamma,ab}=
\mathrm{tr}_0(q_!\,\mathrm{TS}^{\mathrm{red}}\,p^*
(e_{\gamma,a}\boxtimes e_{-\gamma,b}))
\quad\hbox{on}\quad
\mathfrak P^{\mathrm{geom}}_{\gamma,-\gamma;0}.
\]
Form the radical
\(K_{\gamma,\bar p}
=\ker G_{\gamma,\bar p}\cap\ker G_{-\gamma,\bar p}^t\) and the
quotient \(L_{\gamma,\bar p}\); fix splittings.

\item Output
\((\mathcal H_R^{\mathrm{geom}},m_R,\Delta_R,P_R^{\mathrm{geom}},
P_R^{\Pi,\mathrm{geom}},
\{B_{\alpha,\beta}\},\{G_{\gamma,ab}\},\{K_\gamma\},\{L_\gamma\})\).
\end{enumerate}

\paragraph{Type checks.}
\(\mathcal H_R^{\mathrm{geom}}\) is a finite \(\widehat\Gamma_R\)-graded
Hall bialgebra; \(m_R\) is associative on the two-step flag stack;
\(\Delta_R\) is coassociative; the supercommutator of primitives is
primitive (the bialgebra identity); the descended bracket is
\(\Gamma_R^\Pi\)-graded after \(\overline\Pi_{X,*}^{\Theta_R}\);
finite-dimensional in each homogeneous parity piece by
finite-typeness.

\paragraph{Obstructions discharged or invoked.}
\((D_0)\) is partially discharged: the finite Hall package
\((\mathcal H_R^{\mathrm{geom}},m_R,\Delta_R,P_R)\) is constructed.
The seven obstruction classes
\(\mathfrak o_\Theta^{\mathrm{top}},
\mathfrak o_\Theta^{\mathrm{Hoch}},
\mathfrak o_\Theta^{\mathrm{tri}},
\mathfrak o_\Theta^{\mathrm{Jac}},
\mathfrak o_F,
\mathfrak o_\Theta^{\mathrm{cyc}},
\mathfrak o_\Theta^{\mathrm{rad}}\) of
Theorem~\ref{thm:hall-product} are required to vanish on
this finite Hall stage and compatibly in the inverse limit.

\paragraph{Cross-volume anchor.}  The output
\((\mathcal H_R^{\mathrm{geom}},m_R)\) is the finite-stage realization
of the K3-side BKM Hall--Drinfeld double
\(\mathcal D_\hbar(\mathrm{CoHA}_{K3\times E})\) of
\texttt{calabi-yau-quantum-groups:CLAUDE.md~§6.7}.  In the universal
stage chain \(\mathsf P\to\mathsf C\to\mathsf S\) the Hall positive half
\(Y^+(X)\) is precisely \(P_R^{\mathrm{geom}}\) at finite \(R\); the
Drinfeld double \(G(X)=D(Y^+(X))\) requires Hopf pairing, completion,
integral form, and stable-envelope transport which are supplied by
\(G_{\gamma,ab}\) and \((\mathrm{ML})\).  The κ-tuple coordinate
\(\kappa_{\mathrm{BKM}}=5\) is the weight of \(\Delta_5\) and the
output of the level-4 scalar trace; it is not a level-2 invariant of
the carrier itself.

\subsection{Construction~6 — Dirac operator and Pfaffian line}
\label{app:build-dirac-pfaffian}

\paragraph{Tag.} \(\gamma\)-shadow specialization
(orientation gerbe + Pfaffian).

\paragraph{Input.}
The oriented Hall package
\((\mathcal H_R^{\mathrm{geom}},o_R,\{\tau_{i,R,S}\})\) of
Constructions~\ref{app:build-orientation}--\ref{app:build-hall}; the
local wall-sign datum
\(\epsilon_{o,R}^{\mathrm{loc}}\) of
Construction~\ref{app:build-orientation}; the rank-one type-II
Pfaffian crossing datum of
Definition~\ref{def:rank-one-type-ii-crossing} and
Appendix~\ref{app:pfaffian-wall-rank-one}; the half-Hilbert orbit
under \(W^{(2)}(\La^{2,1}_{II})\) of size \(2^{12}=4096\).

\paragraph{Output.}
The Dirac operator on the orientation gerbe
\(D_R\); the Pfaffian line
\(\mathscr L_{\mathrm{Pf},R}=
\det\operatorname{Ext}^1_{\mathrm{red}}(A,A)^{-1}\) (squaring to the
reduced determinant); the protected Pfaffian section
\(\operatorname{pf}_{X,R}\); the global wall-sign character
\(\epsilon_{o,R}:W^{(2)}(\La^{2,1}_{II})\to\{\pm1\}\); the Pfaffian
section asymptotic
\[
\operatorname{pf}_{X,R}
=
64q^{1/2}r^{1/2}s^{1/2}
\prod_{\gamma=(n,l,m)\in A_R}
(1-q^nr^ls^m)^{f(nm,l)}+\hbox{higher-window correction},
\]
which exhausts to \(\Delta_5\) as \(R\to\infty\).

\paragraph{Procedure.}
\begin{enumerate}
\item \emph{Dirac operator.}  On each Darboux chart of
\(\mathfrak M^{\mathrm{sc}}_{R,c}\), with reduced normal self-Ext
splitting \(K^{\mathrm{nor}}_{R,c}\), build the odd Dirac operator
\(D_{R,c}\) acting on \(\Omega^{\mathrm{sp}}\otimes o_{R,c}\)
(spinors on the orientation gerbe); verify that, at a generic
type-II wall point, the local rank-one normal block is
\(D_{\delta_i}=\begin{psmallmatrix}0&u\\-u&0\end{psmallmatrix}\) with
\(\operatorname{Pf}(D_{\delta_i})=u\)
(Definition~\ref{def:rank-one-type-ii-crossing}).

\item \emph{Pfaffian line.}  Set
\(\mathscr L_{\mathrm{Pf},R}=
\det\operatorname{Ext}^1_{\mathrm{red}}(A,A)^{-1}\); verify
\(\mathscr L_{\mathrm{Pf},R}^{\otimes 2}\simeq
\det\operatorname{RHom}_{\mathrm{red}}(A,A)\) by
\(3\)-Calabi--Yau Serre duality.

\item \emph{Local wall sign.}  At each generic point of
\(\Hpl_{\delta_i}\) with normal model
\(K^{\mathrm{nor}}_{\delta_i,R}=[\R\xrightarrow{u_i}\R]\), compute
the monodromy of the Pfaffian section
\(\operatorname{pf}_{\delta_i,R}=\upsilon_iu_i\) along the
wall-reflection loop \(\ell_{\delta_i,R,S}\) closed by the
\((O1)^+\) lift \(\tau_{i,R,S}\):
\[
M_{\ell_{\delta_i,R,S}}(\upsilon_iu_i)
=\upsilon_i(-u_i)=-\upsilon_iu_i,
\qquad
\epsilon_{o,R}^{\mathrm{loc}}(s_{\delta_i})
=(-1)^{N^{\mathrm{Pf}}_{\delta_i,R}}=-1.
\]

\item \emph{Half-Hilbert orbit transport.}  The reflection
\(s_{\delta_i}\) acts on the half-Hilbert orbit (compact source labels
of size \(2^{12}=4096\) parametrized by 12 binary parities tied to the
\(2^{12}\) odd-spin structures on the lattice
\(\La^{2,1}_{II}/2\La^{2,1}_{II}\)).  Sum the local rank-one
contributions over the orbit:
\[
\epsilon_{o,R}(s_{\delta_i})
=\prod_{x\in\hbox{orbit}/2^{12}}
\epsilon_{o,R}^{\mathrm{loc}}(s_{\delta_i};x).
\]
The output is the global character
\(\epsilon_{o,R}:W^{(2)}(\La^{2,1}_{II})\to\{\pm 1\}\).

\item \emph{Compatibility package.}  Verify the component, boundary,
overlap, quotient, and protected integration compatibility package of
Proposition~\ref{prop:appE-local-pfaffian-sign} on every
retained type-II wall stratum.

\item \emph{Pfaffian section.}  Define
\(\operatorname{pf}_{X,R}\) as the protected integration of the
oriented Hall product over
\(\mathcal F_{X,\sigma,S,\le R}^{\mathrm{hyb}}\):
\[
\operatorname{pf}_{X,R}
=
64q^{1/2}r^{1/2}s^{1/2}
I_R^{\mathrm{prot}}\!\left(
\bigotimes_{\gamma\in A_R}\operatorname{pf}_{\gamma,R}\right),
\qquad
\operatorname{pf}_{\gamma,R}=(1-q^nr^ls^m)^{f(nm,l)}.
\]

\item Output
\((D_R,\mathscr L_{\mathrm{Pf},R},\operatorname{pf}_{X,R},
\epsilon_{o,R})\), with the protected integration compatibility data.
\end{enumerate}

\paragraph{Type checks.}
\(\mathscr L_{\mathrm{Pf},R}\) is a Picard-line whose square
isomorphism is supplied by reduced Serre duality;
\(\operatorname{pf}_{X,R}\) is a section of
\(\mathscr L_{\mathrm{Pf},R}\) with the displayed cusp expansion; the
character \(\epsilon_{o,R}\) takes values in \(\{\pm 1\}\) and
satisfies the Coxeter relations of
\(W^{(2)}(\La^{2,1}_{II})\); on each simple type-II reflection
\(\epsilon_{o,R}(s_{\delta_i})=-1\) by the rank-one local model.

\paragraph{Obstructions discharged or invoked.}
This step is the entire site of \((O2)\) (local Pfaffian wall normal
form on type-II walls; the \(2^{12}=4096\) sign sum on the
half-Hilbert orbit transport).  The local rank-one Pfaffian crossing
of Appendix~\ref{app:pfaffian-wall-rank-one} discharges \((O2)\)
generically; the global half-Hilbert orbit transport remains the open
\((O2)\) component of
Open~Problem~\ref{def:O1-plus-weyl-transport}.  The character
identity
\(\epsilon_o=\nu_{\Delta_5}\big|_{W^{(2)}(\La^{2,1}_{II})}\) becomes a
theorem only after \((D_0)\), \((O1)\), \((O1)^+\), and \((O2)\) are
discharged compatibly in \(R\).

\subsection{Construction~7 — Chiral Koszul source coalgebra}
\label{app:build-koszul}

\paragraph{Tag.} \(\gamma\)-shadow specialization
(\emph{Bar = twisting; not bulk}; the bulk is
\(Z^{\mathrm{der}}_{\mathrm{ch}}(\mathcal F_X^{\mathrm{hyb}})\),
constructed elsewhere).

\paragraph{Input.}
The oriented Hall primitives
\(P_R^{\Pi,\mathrm{geom}}\) of
Construction~\ref{app:build-hall}; the hybrid carrier
\(\mathcal F_{X,\sigma,S,\le R}^{\mathrm{hyb}}\) of
Construction~\ref{app:build-hybrid-wrapped}; the augmentation
\(\varepsilon_{X,R}^{\mathrm{hyb}}:
\mathcal F_{X,\sigma,S,\le R}^{\mathrm{hyb}}\to k\mathbf 1\); the
hybrid residual \(\mathfrak O_{\mathrm{hyb},R}=0\); the bounded
length \(N_R\) (positivity of HN height on every non-vacuum colour);
the Beilinson--Drinfeld/Francis--Gaitsgory bar-cobar formalism
\cite{BD,FrancisGaitsgoryCKD}; the conilpotent completed domain on
the elliptic curve \(E\).

\paragraph{Output.}
The conilpotent chiral coalgebra on \(E\):
\[
C_{X,R}^{\mathrm{geom}}
=\bar B_E^{\mathrm{ch}}
(\mathcal F_{X,\sigma,S,\le R}^{\mathrm{hyb}},
\varepsilon_{X,R}^{\mathrm{hyb}})
=k\mathbf 1\oplus\bigoplus_{1\le p\le N_R}
\bar B_E^{\mathrm{ch},p}(\overline{\mathcal F}_{X,R}^{\mathrm{geom}});
\]
the bar-length filtration
\(F^{\mathrm{co}}_{R,p}\); the chiral collision coproduct
\(\Delta_R^{\mathrm{ch}}\); the bar comparison
\(b_{X,R}=\mathrm{id}_{\bar B_E^{\mathrm{ch}}(\cdot)}\); the source
cobar counit
\(\varepsilon^{\mathrm{src}}_{\mathrm{bar},R}:
\Omega_E^{\mathrm{ch}}C_{X,R}^{\mathrm{geom}}
\xrightarrow{\simeq}
\mathcal F_{X,\sigma,S,\le R}^{\mathrm{hyb}}\); after a separate
source-to-target quasi-isomorphism
\(\vartheta_R:\mathcal F_{X,\sigma,S,\le R}^{\mathrm{hyb}}
\xrightarrow{\simeq}\mathsf{Den}_{\Delta_5,E,\le R}\), the cobar
comparison
\(\Theta_{\mathrm{Kos},R}=\vartheta_R\circ\varepsilon^{\mathrm{src}}_{\mathrm{bar},R}\).

\paragraph{Procedure.}
\begin{enumerate}
\item \emph{Augmentation kernel.}  Set
\(\overline{\mathcal F}_{X,R}^{\mathrm{geom}}
=\ker(\varepsilon_{X,R}^{\mathrm{hyb}})\); verify the bounded-length
hypothesis (positivity of HN height on every non-vacuum surviving
colour, hence at most \(N_R\) factors in any non-zero word).

\item \emph{Bar coalgebra.}  Form the reduced chiral bar
\[
C_{X,R}^{\mathrm{geom}}
=k\mathbf 1\oplus\bigoplus_{1\le p\le N_R}
\bar B_E^{\mathrm{ch},p}
(\overline{\mathcal F}_{X,R}^{\mathrm{geom}}),
\]
length \(p\) summand generated by ordered \(p\)-fold inputs from the
augmentation ideal of total HN height \(\le R\) (height \(>R\) outputs
zero); coaugmentation by the inclusion of the vacuum summand; counit by
projection.

\item \emph{Bar-length filtration.}
\[
F^{\mathrm{co}}_{R,p}C_{X,R}^{\mathrm{geom}}
=k\mathbf 1\oplus\bigoplus_{1\le j\le p}
\bar B_E^{\mathrm{ch},j}
(\overline{\mathcal F}_{X,R}^{\mathrm{geom}}),\qquad
0\le p\le N_R.
\]

\item \emph{Collision coproduct.}  For
\(x=[x_1|\cdots|x_p]\), define the cut-map coproduct
\[
\Delta_R^{\mathrm{ch}}(x)
=\mathbf1\otimes x+x\otimes\mathbf1
+\sum_{i=1}^{p-1}[x_1|\cdots|x_i]\otimes[x_{i+1}|\cdots|x_p],
\]
with each cut realized by the corresponding finite collision/flag
pull-push map in
\(\mathsf{Corr}^{E,\mathrm{hyb}}_{R,\mathrm{geom}}\) of
Construction~\ref{app:build-hybrid-wrapped}.  Verify
coassociativity by the four-input pentagon and the eight-word flag
coherences; verify counit identities via vacuum collision flags.

\item \emph{Conilpotence.}  Verify
\((\overline\Delta_R^{\mathrm{bar}})^{N_R+1}=0\): each reduced cut
strictly lowers bar length on each branch.

\item \emph{Bar comparison.}  Take \(b_{X,R}\) to be the identity of
\(\bar B_E^{\mathrm{ch}}
(\mathcal F_{X,\sigma,S,\le R}^{\mathrm{hyb}},
\varepsilon_{X,R}^{\mathrm{hyb}})\), eliminating the coalgebra-side
finite Koszul residuals
\(\mathfrak o^C_R, \mathfrak o^{\mathrm{fil}}_R,
\mathfrak o^{\mathrm{conil}}_R, \mathfrak o^{\Delta,\mathrm{ch}}_R\)
of Definition~\ref{def:chiral-koszul-source-datum}.

\item \emph{Source cobar counit.}  Apply Francis--Gaitsgory chiral
Koszul duality
\cite[Theorem~6.2.2]{FrancisGaitsgoryCKD} in the conilpotent completed
domain to obtain
\(\varepsilon^{\mathrm{src}}_{\mathrm{bar},R}:
\Omega_E^{\mathrm{ch}}C_{X,R}^{\mathrm{geom}}
\xrightarrow{\simeq}
\mathcal F_{X,\sigma,S,\le R}^{\mathrm{hyb}}\); verify it is a
source-internal quasi-isomorphism of chiral algebras on \(E\), not a
homotopy inverse of the target counit.

\item \emph{Source-to-target identification.}  Construct
\(\vartheta_R\) from the geometric primitive representatives, Hall
products, coproducts, pairings, and HN transition compatibility;
verify the residual
\(\mathfrak o^\vartheta_R\) of
Corollary~\ref{cor:source-bar-remaining-cobar-obstruction} vanishes;
set \(\Theta_{\mathrm{Kos},R}=\vartheta_R\circ
\varepsilon^{\mathrm{src}}_{\mathrm{bar},R}\).

\item \emph{Inverse limit.}  Verify the Koszul Mittag--Leffler
condition \(R^1\varprojlim=0\) for the inverse systems of source
coalgebras, completed chiral bar/cobar constructions, target
truncations, and the cones of \(b_{X,R}\) and
\(\Theta_{\mathrm{Kos},R}\) \cite[\S3.5]{WeibelHA}.  Form
\[
C_X^{\mathrm{geom}}
=\varprojlim_R C_{X,R}^{\mathrm{geom}},\qquad
\Theta_{\mathrm{Kos}}:
\Omega_E^{\mathrm{ch}}C_X^{\mathrm{geom}}
\xrightarrow{\simeq}\mathsf{Den}_{\Delta_5,E}.
\]

\item Output
\((C_{X,R}^{\mathrm{geom}},F^{\mathrm{co}}_{R,\bullet},
\Delta_R^{\mathrm{ch}},b_{X,R},\Theta_{\mathrm{Kos},R})\) at every
\(R\), with successor compatibility.
\end{enumerate}

\paragraph{Type checks.}
\(C_{X,R}^{\mathrm{geom}}\) is a coaugmented conilpotent chiral
coalgebra on \(E\) with bounded bar-length filtration of length at
most \(N_R\); \(\Delta_R^{\mathrm{ch}}\) is coassociative and counital
on the chosen flag stacks; the bar-cobar inversion is supplied by
Francis--Gaitsgory chiral Koszul; the cobar comparison
\(\Theta_{\mathrm{Kos},R}\) is source-internal, not the target
bar-cobar counit (the source/target separation of
Lemma~\ref{lem:source-target-koszul-separation}).

\paragraph{Obstructions discharged or invoked.}
The Koszul Mittag--Leffler exactness hypothesis is the
\textup{(D0)}-component for the chiral Koszul lane.  The total finite
Koszul residual
\(\mathfrak O_{\mathrm{Kos},R}=
(\mathfrak o^C_R,\mathfrak o^{\mathrm{fil}}_R,
\mathfrak o^{\mathrm{conil}}_R,\mathfrak o^{\Delta,\mathrm{ch}}_R,
\{\mathfrak o^{C,\mathrm{tr}}_{R'R}\},\mathfrak o^b_R,
\mathfrak o^\Omega_R,\mathfrak o^{\mathrm{hyb}}_R,
\mathfrak o^\Pi_R,\mathfrak o^{\mathrm{sep}}_R,
\mathfrak o^{\Prim}_R)\)
vanishes after this step.

\paragraph{Cross-volume firewall.}  The bar coalgebra
\(\bar B_E^{\mathrm{ch}}
(\mathcal F_{X,R}^{\mathrm{hyb}})\) is a \emph{twisting / coupling}
coalgebra (Maurer--Cartan classification), \emph{not} the bulk
\(Z^{\mathrm{der}}_{\mathrm{ch}}(\mathcal F_X^{\mathrm{hyb}})
=\mathrm{ChirHoch}^\bullet(\mathcal F_X^{\mathrm{hyb}},
\mathcal F_X^{\mathrm{hyb}})\); cf.\
\texttt{chiral-bar-cobar-vol2:CLAUDE.md~§6} (universal stage chain
\(\mathsf P\to\mathsf C\to\mathsf S\to\mathsf Z\)).
The chiral Koszul comparison
\(\Theta_{\mathrm{Kos}}:\Omega_E^{\mathrm{ch}}C_X^{\mathrm{geom}}
\xrightarrow{\simeq}\mathsf{Den}_{\Delta_5,E}\) is a comparison
between two algebras of the same level; the homotopy inverse to the
target-internal counit
\(\Omega_E^{\mathrm{ch}}\bar B_E^{\mathrm{ch}}
(\mathsf{Den}_{\Delta_5,E})\to\mathsf{Den}_{\Delta_5,E}\) is not used
in either direction.

\subsection{Construction~8 — Primitive recognition}
\label{app:build-recognition}

\paragraph{Tag.} \(\delta\)-endpoint convergence
(source-to-target identification).

\paragraph{Input.}
The descended source primitive Lie superalgebra
\(P_X^{\Pi,\mathrm{red}}\) of
Construction~\ref{app:build-hall} (after the radical quotient
\(\overline{\mathrm{Rad}\langle\,\cdot,\cdot\rangle_X}\)); the imported
target denominator algebra
\(\fg_{\Delta_5}=G(\mathscr B_{\Delta_5})\) on
\(\La^{2,1}_{II}\) of Definition~\ref{def:bkm-imaginary-simple-roots} and
Definition~\ref{def:bkm-relations}; the increasing cofinal
sequence of finite downward-saturated windows
\(W_1\subset W_2\subset\cdots\) with
\(\bigcup_\nu W_\nu=\mathcal R_+\); the Borcherds--Kac--Moody data
\((H,I,p,\alpha_i,(\,,\,))\) of
Theorem~\ref{thm:primitive-recognition}.

\paragraph{Output.}
The cofinal finite-window primitive-recognition datum
\(\{\textup{(R0),(R1),(R2),(R3),(R4),(R5),(R6),(R7),(R8)}\}\) of
Definition~\ref{def:cofinal-primitive-recognition-datum} for every
\(W_\nu\), compatible under \(W_\nu\subset W_{\nu+1}\); the source
representatives
\(e_i^X,f_i^X,h_i^X\); the source presentation map
\(\pi_W:\widehat{\mathfrak F}(\mathscr B_{\Delta_5})_W
\to P_X^{\Pi,\mathrm{red}}\big|_W\); the kernel equality
\(\ker\pi_W=
(\overline{J_{\mathrm{BK}}+\mathrm{Rad}_{\mathrm{GN}}})_W\); the
isomorphism
\[
\boxed{\quad
P_X^{\Pi,\mathrm{red}}\cong\fg_{\Delta_5},\quad
\Omega_E^{\mathrm{ch}}C_X\simeq\mathsf{Den}_{\Delta_5,E}.
\quad}
\]

\paragraph{Procedure.}
For every cofinal \(W_\nu\):
\begin{enumerate}
\item[\textup{(R0)}] \emph{Compact source labels.}  Verify every
finite source basis vector in \(W_\nu\cup(-W_\nu)\cup\{0\}\) satisfies
the compact source-admissibility condition of
Theorem~\ref{thm:finite-stage-dirac-igusa-pro-object}.

\item[\textup{(R1)}] \emph{Representatives and parities.}  Supply
parity-homogeneous primitive representatives
\(e_i^X\in P^{\Pi,\mathrm{red}}_{X,\alpha_i}\),
\(f_i^X\in P^{\Pi,\mathrm{red}}_{X,-\alpha_i}\),
\(h_i^X=\iota_H^{-1}(\alpha_i)\) for every simple-root index \(i\) with
\(\alpha_i\in W_\nu\), with the Gritsenko--Nikulin parity fibres
\(\mu_{\overline 0}(a),\mu_{\overline 1}(a)\) for imaginary simple
roots and even real representatives for \(\delta_1,\delta_2,\delta_3\).

\item[\textup{(R2)}] \emph{Relations.}  Verify in the protected Hall
super-bracket the Cartan, Chevalley, real Serre, Borcherds
orthogonality, and super sign relations for every relation whose
displayed source and target degrees lie in
\(W_\nu\cup(-W_\nu)\cup\{0\}\).  In the real rank-three subsystem the
Cartan matrix is
\(((\delta_i,\delta_j))=\mathrm{diag}(2)-2(J-\mathrm{diag}(1))\) with
real Serre exponent~3.

\item[\textup{(R3)}] \emph{Full finite parity dimensions.}  For every
\(\beta\in W_\nu\) verify
\(\dim P^{\Pi,\mathrm{red}}_{X,\beta,\overline p}
=\rootmult_{\overline p}^{\mathrm{GN}}(\beta)\) for both parities
\(p=0,1\).  In the first timelike channel
\(\beta=\delta_{123}\),
Proposition~\ref{prop:bkm-delta123-presentation-count} requires
\((d^{\mathrm{GN}}_{(1,1,1),\overline 0},
d^{\mathrm{GN}}_{(1,1,1),\overline 1})=(29,93)\).  Determinant
information \(29-93=-64=f(1,1)\) is consistent but not the
recognition.

\item[\textup{(R4)}] \emph{Hopf pairing and radical.}  Compute the
positive-negative pairing matrices in parity blocks; verify the
kernel agrees with the
Gritsenko--Nikulin/Kac invariant-pairing kernel; verify Lie-ideal and
coideal stability under all finite brackets; verify carrying under
HN transitions to the next window.

\item[\textup{(R5)}] \emph{No extra relations.}  Verify the kernel
equality
\(\ker\pi_\nu=(\overline{J_{\mathrm{BK}}+\mathrm{Rad}_{\mathrm{GN}}})_{\le W_\nu}\) on the finite triangular span generated by
\(W_\nu\cup(-W_\nu)\cup\{0\}\).

\item[\textup{(R6)}] \emph{Generation.}  Verify that every source
primitive root space in \(W_\nu\cup(-W_\nu)\), after the finite
radical quotient, is generated by the Cartan and the supplied
simple-root representatives through brackets whose intermediate
positive degrees stay in \(W_\nu\).

\item[\textup{(R7)}] \emph{Completed PBW compatibility.}  Verify the
PBW-graded isomorphism
\(\mathrm{gr}_{\mathrm{PBW}}U(P_X^{\Pi,\mathrm{red}})_{\le W_\nu}
\xrightarrow{\sim}
\mathrm{gr}_{\mathrm{PBW}}U(\fg_{\Delta_5})_{\le W_\nu}\); verify
strict preservation under \(W_{\nu+1}\to W_\nu\) transition maps.

\item[\textup{(R8)}] \emph{Exact triangular completion.}  Verify
strictness of the transition maps and closedness of images on every
finite window, so that passage to the inverse limit commutes with the
finite kernels, cokernels, radicals, PBW associated gradeds, and
parity pieces.  Equivalently the relevant \(\lim^1\) terms vanish for
these triangular finite-window systems.

\item Output the global isomorphism
\(P_X^{\Pi,\mathrm{red}}\cong\fg_{\Delta_5}\) by passage to the inverse
limit (compatible cofinal datum); together with the chiral Koszul
comparison of Construction~\ref{app:build-koszul} this gives
\(\Omega_E^{\mathrm{ch}}C_X\simeq\mathsf{Den}_{\Delta_5,E}\) and hence
the protected Pfaffian identification
\(\operatorname{Pf}_{\mathrm{prot}}(\mathfrak D_X^{\mathrm{DI}})=\Delta_5\).
\end{enumerate}

\paragraph{Type checks.}
The source presentation map \(\pi:\widehat{\mathfrak F}(\mathscr B_{\Delta_5})\to P_X^{\Pi,\mathrm{red}}\) is
continuous and degree-preserving; \textup{(R5)} gives injectivity
after the completed radical quotient; \textup{(R6)} gives
surjectivity; \textup{(R7)} gives PBW compatibility; \textup{(R8)}
gives exactness in the inverse-limit step.  No claim on signed
superdimensions \(\sdim=f(nm,l)\) substitutes for the two-integer
parity dimensions of \textup{(R3)}.

\paragraph{Obstructions discharged or invoked.}
This step inherits the obstruction tower discharged by
Constructions~1--7 and adds the cofinal exact-completion residual
\textup{(R8)}.  The Beilinson--Drinfeld bar-cobar firewall is
respected throughout: \textup{(R5)} is a source kernel-equality
theorem, not an inversion of the target bar-cobar counit;
\textup{(R7)} is a PBW graded comparison of source and target
enveloping algebras, not a graph-equality argument from the scalar
square.

\paragraph{Cross-volume anchor.}  The output is
\texttt{chiral-bar-cobar-vol2:CLAUDE.md~§9} Construction
Problem~1: the operator-level construction of
\(\mathfrak D_X^{\mathrm{DI}}\) for \(K3\times E\) with
\(\operatorname{Pf}_{\mathrm{prot}}(\mathfrak D_X^{\mathrm{DI}})=\Delta_5\), under the
four obstructions \((D_0),(O1),(O1)^+,(O2)\).  The Drinfeld doubling
\(Y^+(X)\to G(X)=D(Y^+(X))\) of
\texttt{calabi-yau-quantum-groups:CLAUDE.md} requires the Hopf
pairing of \textup{(R4)}, completion by \textup{(R8)}, integral form,
and stable-envelope transport which are supplied by the same
data.  The κ-tuple coordinate \(\kappa_{\mathrm{BKM}}=5\) is the
weight of \(\Delta_5\) and lives at level~4 of the universal stage
chain; it is a consequence of the construction, not an input to it.

\subsection{The eight-step audit}
\label{app:build-audit}

The eight construction steps form a strict ordered sequence of
input/output type discipline.  An audit of any candidate construction of
\(\mathfrak D_X^{\mathrm{DI}}\) reduces to verifying, for every
\(R\in\mathcal R_{\mathrm{HN}}\):
\begin{enumerate}
\item the lattice combinatorics of
Construction~\ref{app:build-charge-window};
\item the finite-type quasi-smooth derived stacks of
Construction~\ref{app:build-finite-moduli};
\item the hybrid local/wrapped factorization datum and the eight-word
two-step flag atlas of Construction~\ref{app:build-hybrid-wrapped};
\item the four Cech--Borel orientation classes and the
\((O1)^+\)~Coxeter coherence of
Construction~\ref{app:build-orientation};
\item the finite Hall bialgebra and primitive matrices of
Construction~\ref{app:build-hall};
\item the rank-one Pfaffian crossing and the half-Hilbert orbit
transport of Construction~\ref{app:build-dirac-pfaffian};
\item the bar coalgebra, conilpotent collision coproduct, and the
Francis--Gaitsgory cobar quasi-isomorphism of
Construction~\ref{app:build-koszul};
\item the cofinal datum
\(\textup{(R0)}\)--\(\textup{(R8)}\) of
Construction~\ref{app:build-recognition};
\end{enumerate}
together with the compatible vanishing of the obstruction tuples
\((\mathfrak O^{E_3-\mathrm{src}}_R,
\mathfrak O_{\mathrm{hyb},R},
\mathfrak O_{\mathrm{Kos},R},
\mathfrak O^{\mathrm{O1}}_R,
\mathfrak O^{\mathrm{O1}^+}_R,
\mathfrak O^{\mathrm{O2}}_R,
\mathfrak O_R^{D_0})\) at every height \(R\), and the
Mittag--Leffler exactness condition \textup{[ML]} of
Definition~\ref{def:O1-strong-reduced-orientation} on the cofinal HN inverse
system.  Each obstruction is attributable to the unique step
in which its components are constructed.

\paragraph{Order rigidity.}  Each of the seven adjacent input/output
dependencies fails type-checking under exchange:
\begin{enumerate}
\item charge window $\to$ finite moduli: $\Lambda_{\mathrm{ample}}$
  parametrises the HN strata; without the window, the moduli stack
  has no finite cofinal subsystem.
\item finite moduli $\to$ hybrid carrier: $\mathrm{Ran}^{\mathrm{hyb}}(E)$
  needs the finite Hall stack as factor on which to graft the
  wrapped/local atlas.
\item hybrid carrier $\to$ reduced orientation: the cosection-twisted
  Joyce--Upmeier line bundle is defined on the carrier, not on the
  raw moduli (the cosection lives on $\mathrm{Ran}^{\mathrm{hyb}}$).
\item reduced orientation $\to$ Hall correspondences: the Pfaffian
  line bundle on the Hall stack is the orientation-twisted sign data;
  without orientation, the Hall product is unsigned.
\item Hall correspondences $\to$ Pfaffian line + Dirac operator: the
  primitive Hall bracket $[-,-]_X$ defines the Hall--Drinfeld
  correspondence whose Pfaffian section is the Dirac--Igusa Pfaffian.
\item Pfaffian line $\to$ chiral Koszul source coalgebra: the source
  $C_X$ is the bar of the augmented hybrid factorization algebra
  pulled back along the Pfaffian comparison
  $\iota_{\mathrm{aut}}$.
\item chiral Koszul source $\to$ primitive recognition: the primitive
  part $P_X^{\Pi,\mathrm{red}}$ is the Frobenius--Serre primitive
  filtration on the source, which the recognition theorem matches to
  $\mathfrak g_{\Delta_5}$.
\end{enumerate}
Two diagnostic exchanges: orientation before hybrid leaves the middle
type-II wall \(\delta_2\leftrightarrow(0,1,1)\) without a
representative, by
Lemma~\ref{app:pfaffian-wall-three-walls}; recognition before the
Koszul cobar leaves the identification
\(\Omega_E^{\mathrm{ch}}C_X\simeq\mathsf{Den}_{\Delta_5,E}\)
unconstructed.  All other adjacent exchanges fail by the same
input/output type criterion.

\paragraph{Conditional theorem.}  Granted the eight constructions and
the four obstructions \((D_0),(O1),(O1)^+,(O2)\),
Theorem~\ref{thm:pfaffian-dirac-finite-stage} reads
\[
\operatorname{Pf}_{\mathrm{prot}}(\mathfrak D_X^{\mathrm{DI}})=\Delta_5,\qquad
\epsilon_o=\nu_{\Delta_5}\big|_{W^{(2)}(\La^{2,1}_{II})},\qquad
P_X^{\Pi,\mathrm{red}}\cong\fg_{\Delta_5}.
\]
The scalar square
\(Z_{\mathrm{BPS}}^{K3\times E}=
(\Phi_{10}^{\mathrm{un}})^{-1}=\Delta_5^{-2}\) is the level-4 trace
of the constructed protected algebra; promotion to the 3d
gravitational path integral is not claimed and requires the
Hall--Borcherds residual of
\texttt{chiral-bar-cobar-vol2}.

%%% End of Appendix F.
